\chapter{连续函数}
在本章中，我们考察分析的拓扑基础，并给出它的一些初步应用。我们主要局限于度量空间的拓扑，因为度量空间的理论是分析中很大一部分内容的框架，同时又足够简单、具体，从而能为初学者尽量减少困难。即便如此，度量空间的概念对于更深入的数学研究而言仍不够一般，因此，只要可能，我们都给出了在一般拓扑空间中依然成立的证明。各定理在一般拓扑空间中成立的程度，将在每一节的末尾加以讨论。这些评论在初读本书时可以略去，它们为读者提供了抽象拓扑的入门介绍。

在第一节中，我们考察度量空间之间的连续函数。特别地，我们利用上一章关于收敛序列的结果来研究连续性。

第 2 节专门讨论开集的概念。这里的一个关键结果是将连续函数刻画为具有如下性质的函数：每个开集的原像都是开集。

在下一节中，我们讨论紧致度量空间。特别地，我们证明，对于度量空间而言，紧致性与序列紧致性是一回事。紧致性的巨大重要性已经在本节所给出的应用中显现出来。例如，利用关于紧致度量空间上连续实值函数的极值定理，我们证明 $ \mathbb{K}^n $ 上的所有范数都是等价的，并给出代数基本定理的一个证明。

在第 4 节中，我们研究连通空间与道路连通空间。特别地，我们证明，对于赋范向量空间的开子集，这两个概念是一致的。作为连通性的一个重要应用，我们证明介值定理的一个广义版本。

在本次涉足抽象拓扑、为后续各章的分析研究奠定基础之后，我们在本章剩余的两

节中转向对实函数的研究。在简短的第五节中，我们讨论实变量单调函数的行为，并特别证明了连续单调函数的反函数定理。

与本章前五节相对抽象的性质形成对照的是，在最后一节（这一节相当长）中，我们研究指数函数及其相关函数：对数、幂函数与三角函数。在这一研究中，我们实际上运用了本章所引入的几乎所有方法与定理。

\section{连续性}

经验表明，尽管一般来说函数可能非常复杂且难以描述，但出现在应用中的函数却具有某些重要的定性性质，其中之一就是连续性。对于函数 $ f: X \to Y $，连续（或不连续）度量的是像 $ f(X) \subseteq Y $ 中的“微小变化”如何由定义域 $ X $ 中相应的“微小变化”引起。为使其有意义，集合 $ X $ 与 $ Y $ 必须被赋予某种额外结构，以便能对“微小变化”给出精确的含义。度量空间显然是具备这种额外结构的集合的候选。

\subsection*{基本性质与例子}

设 $f: X \to Y$ 是度量空间\footnote{在 $ X $ 中的度量 $ d_X $ 与在 $ Y $ 中的度量 $ d_Y $，我们通常都写作 $ d $。} $(X, d_X)$ 与 $(Y, d_Y)$ 之间的函数。则 $f$ 在 $x_0 \in X$ 处 \textbf{连续}，若对 $Y$ 中 $f(x_0)$ 的每个邻域 $V$，存在 $X$ 中 $x_0$ 的邻域 $U$，使得 $f(U) \subseteq V$。

\bookfigure[0.51]{figures/30212f0df5e73c18.jpg}

因此，为证明 $f$ 在 $x_0$ 处的连续性，我们假设给定了 $f(x_0)$ 的任意邻域 $V$，然后证明存在 $x_0$ 的邻域 $U$，使得 $f(U) \subseteq V$，即对所有 $x \in U$ 有 $f(x) \in V$。

函数 $f: X \to Y$ 是 \textbf{连续的}，若它在 $X$ 的每一点处连续。若 $f$ 在 $x_0$ 处不连续，则称 $f$ 在 $x_0$ 处 \textbf{不连续的}。最后，$f$ 是 \textbf{不连续的}，若它在 $X$ 的（至少）一点处不连续，即若 $f$ 不连续。从 $X$ 到 $Y$ 的所有连续函数构成的集合记为 $C(X,Y)$。显然 $C(X,Y)$ 是 $Y^X$ 的子集。

连续性的这个定义使用了邻域的概念，因此相当简单。在具体情形中，下面的等价表述往往更有用。

\qyxffntextdef{1}{记号 $ \delta := \delta(x_0, \varepsilon) $ 表明 $ \delta $ 一般依赖于 $ x_0 \in X $ 与 $ \varepsilon > 0 $。}
\begin{proposition}

函数 $ f: X \to Y $ 在 $ x_0 \in X $ 处连续，当且仅当对每个 $ \varepsilon > 0 $，存在某个\qyxffnmark{1} $ \delta := \delta(x_0, \varepsilon) > 0 $，满足性质

\begin{equation}
d\big(f(x_{0}),f(x)\big)<\varepsilon\quad\text{对所有 }x\in X\text{ 使得 }d(x_{0},x)<\delta.
\end{equation}

\proofhead
‘$\Rightarrow$’ 设 $f$ 在 $x_0$ 处连续，且 $\varepsilon > 0$。则对于邻域 $V := \mathbb{B}_Y(f(x_0), \varepsilon) \in \mathcal{U}_Y(f(x_0))$，存在某个 $U \in \mathcal{U}_X(x_0)$，使得 $f(U) \subseteq V$。由定义，存在某个 $\delta := \delta(x_0, \varepsilon) > 0$，使得 $\mathbb{B}_X(x_0, \delta) \subseteq U$。因此
\[
f\bigl(\mathbb{B}_{X}(x_{0},\delta)\bigr)\subseteq f(U)\subseteq V=\mathbb{B}_{Y}\bigl(f(x_{0}),\varepsilon\bigr).
\]

这些包含关系蕴含 (1.1)。

‘⇐’ 假设 (1.1) 成立，且 $ V \in \mathcal{U}_Y(f(x_0)) $。则存在某个 $ \varepsilon > 0 $，使得 $ \mathbb{B}_Y(f(x_0), \varepsilon) \subseteq V $。由 (1.1)，存在某个 $ \delta > 0 $，使得 $ U := \mathbb{B}_X(x_0, \delta) $ 的像包含于 $ \mathbb{B}_Y(f(x_0), \varepsilon) $，从而也包含于 $ V $。因此 $ f $ 在 $ x_0 $ 处连续。 ■

\end{proposition}

\begin{corollary}

设 $E$ 与 $F$ 为赋范向量空间，且 $X \subseteq E$。则 $f: X \to F$ 在 $x_0 \in X$ 处连续，当且仅当对每个 $\varepsilon > 0$，存在某个 $\delta := \delta(x_0, \varepsilon) > 0$，满足

\begin{equation*}
\begin{array}{r l}{\|f(x)-f(x_{0})\|_{F}<\varepsilon}&{{}{\text{对所有 }}x\in X{\text{ 使得 }}\|x-x_{0}\|_{E}<\delta.}\end{array}
\end{equation*}

\proofhead
这直接由赋范向量空间中度量的定义得出。

假设 $E := F := \mathbb{R}$，且函数 $f: X \to \mathbb{R}$ 由下面的图像给出。
\bookfigure[0.57]{figures/6473b74b93c53473.jpg}

则 $f$ 在 $x_0$ 处连续，因为对每个 $\varepsilon > 0$，存在某个 $\delta > 0$，使得 $U := (x_0 - \delta, x_0 + \delta)$ 的像包含于 $V := (f(x_0) - \varepsilon, f(x_0) + \varepsilon)$。

另一方面，不存在 $\delta>0$ 使得对所有 $x\in(x_1,x_1+\delta)$ 都有 $|f(x)-f(x_1)|<\varepsilon_1$，因此 $f$ 在 $x_1$ 处不连续。

\end{corollary}

\begin{example}

在下面的例子中，X 与 Y 都是度量空间。

\begin{items}
\item

平方根函数 $ \mathbb{R}^+ \to \mathbb{R}^+ $，$ x \mapsto \sqrt{x} $ 是连续的。

\proofhead
设 $ x_0 \in \mathbb{R}^+ $，$ \varepsilon > 0 $。若 $ x_0 = 0 $，我们令 $ \delta := \varepsilon^2 > 0 $。则
\[
\left|\sqrt{x}-\sqrt{x_{0}}\right|=\sqrt{x}<\varepsilon,\qquad x\in\left[0,\delta\right).
\]

否则 $ x_0 > 0 $，我们取 $ \delta := \delta(x_0, \varepsilon) := \min\{\varepsilon \sqrt{x_0}, x_0\} $。则

\[
\left|\sqrt{x}-\sqrt{x_{0}}\right|=\left|\frac{x-x_{0}}{\sqrt{x}+\sqrt{x_{0}}}\right|<\frac{\left|x-x_{0}\right|}{\sqrt{x_{0}}}\leq\varepsilon
\]

对所有 $ x \in (x_0 - \delta, x_0 + \delta) $ 成立。 ■

\item

下取整函数 $ \lfloor \cdot \rfloor $：$ \mathbb{R} \to \mathbb{R} $，$ x \mapsto \lfloor x \rfloor := \max\{k \in \mathbb{Z} ; k \leq x\} $ 在 $ x_0 \in \mathbb{R} \setminus \mathbb{Z} $ 处连续，而在 $ x_0 \in \mathbb{Z} $ 处不连续。

\proofhead
若 $ x_0 \in \mathbb{R} \backslash \mathbb{Z} $，则存在唯一的 $ k \in \mathbb{Z} $，使得 $ x_0 \in (k, k+1) $。若取 $ \delta := \min\{x_0 - k, k+1 - x_0\} > 0 $，则显然有
\[
|\lfloor x\rfloor-\lfloor x_{0}\rfloor|=0,\qquad x\in(x_{0}-\delta,x_{0}+\delta)~.
\]

因此下取整函数 $ \lfloor \cdot \rfloor $ 在 $ x_{0} $ 处连续。
否则，对 $ x_0 \in \mathbb{Z} $，对所有 $ x < x_0 $ 有不等式 $ \left|\lfloor x \rfloor - \lfloor x_0 \rfloor\right| = \lfloor x_0 \rfloor - \lfloor x \rfloor \geq 1 $。故不存在 $ x_0 $ 的邻域 $ U $，使得对所有 $ x \in U $ 都有 $ \left|\lfloor x \rfloor - \lfloor x_0 \rfloor\right| < 1/2 $。即 $ \lfloor \cdot \rfloor $ 在 $ x_0 $ 处不连续。

\item

狄利克雷函数 $ f: \mathbb{R} \to \mathbb{R} $ 定义为

\begin{equation*}
f(x):=\left\{\begin{array}{l l}{1,}&{\quad x\in\mathbb{Q},}\\ {0,}&{\quad x\in\mathbb{R}\backslash\mathbb{Q},}\end{array}\right.
\end{equation*}

处处不连续，即它在每个 $ x_0 \in \mathbb{R} $ 处都不连续。

\proofhead
设 $ x_0 \in \mathbb{R} $。由于有理数 $ \mathbb{Q} $ 与无理数 $ \mathbb{R} \setminus \mathbb{Q} $ 都在 $ \mathbb{R} $ 中稠密（见命题 I.10.8 与 I.10.11），故在 $ x_0 $ 的每个邻域中都有某个 $ x $，使得 $ |f(x) - f(x_0)| = 1 $。因此 $ f $ 在 $ x_0 $ 处不连续。
\item

假设 $ f: X \to \mathbb{R} $ 在 $ x_0 \in X $ 处连续，且 $ f(x_0) > 0 $。则存在 $ x_0 $ 的邻域 $ U $，使得对所有 $ x \in U $ 都有 $ f(x) > 0 $。

\proofhead
令 $ \varepsilon := f(x_0)/2 > 0 $。则存在 $ x_0 $ 的邻域 $ U $，使得
\[
f(x_{0})-f(x)\leq|f(x)-f(x_{0})|<\varepsilon=\frac{f(x_{0})}{2},\qquad x\in U.
\]

因此对所有 $ x \in U $ 都有 $ f(x) > f(x_0)/2 > 0 $。 ■

\item

函数 $ f: X \to Y $ 以利普希茨常数 $ \alpha > 0 $ 利普希茨连续，若

\[
d\big(f(x),f(y)\big)\leq\alpha d(x,y)~,\qquad x,y\in X~.
\]

每个利普希茨连续函数都是连续的。\footnote{逆命题不成立。见习题 18。}

\proofhead
给定 $ x_0 \in X $ 与 $ \varepsilon > 0 $，令 $ \delta := \varepsilon/\alpha $。则 $ f $ 的连续性由命题 1.1 得出。注意，在此情形中 $ \delta $ 与 $ x_0 \in X $ 无关。
\item

任何常值函数 $ X \to Y $，$ x \mapsto y_0 $ 都是利普希茨连续的。

\item

恒等函数 id：$X \to X$，$x \mapsto x$ 是利普希茨连续的。

\item

若 $ E_1, \ldots, E_m $ 为赋范向量空间，则 $ E := E_1 \times \cdots \times E_m $ 关于例 II.3.3(c) 中的积范数 $ \|\cdot\|_\infty $ 是赋范向量空间。典范投影

\[
\operatorname{pr}_{k}\colon E\to E_{k}~,\quad x=(x_{1},\ldots,x_{m})\mapsto x_{k}~,\qquad1\leq k\leq m~,
\]

都是利普希茨连续的。特别地，投影 $ \operatorname{pr}_k : \mathbb{K}^m \to \mathbb{K} $ 是利普希茨连续的。

\proofhead
对 $ x = (x_1, \ldots, x_m) $ 与 $ y = (y_1, \ldots, y_m) $，我们有
\[
\|\operatorname*{p r}_{k}(x)-\operatorname*{p r}_{k}(y)\|_{E_{k}}=\|x_{k}-y_{k}\|_{E_{k}}\leq\|x-y\|_{\infty},
\]

这蕴含 $ \operatorname{pr}_{k} $ 的利普希茨连续性。对于余下的断言，见命题 II.3.12。

\item

函数 $ z \mapsto \mathrm{Re}(z) $、$ z \mapsto \mathrm{Im}(z) $ 与 $ z \mapsto \overline{z} $ 中的每一个都在 $ \mathbb{C} $ 上利普希茨连续。

\proofhead
这由不等式
\[
\max\left\{\mid\operatorname{Re}(z_{1})-\operatorname{Re}(z_{2})\mid,\mid\operatorname{Im}(z_{1})-\operatorname{Im}(z_{2})\mid\right\}\leq\left|z_{1}-z_{2}\right|=\left|\overline{z}_{1}-\overline{z}_{2}\right|,\qquad z_{1},z_{2}\in\mathbb{C},
\]

得出，而该不等式来自命题 I.11.4。 ■

\item

设 E 是赋范向量空间。则范数函数

\[
\|\cdot\|\colon E\to\mathbb{R}~,\quad x\mapsto\|x\|
\]

是利普希茨连续的。

\proofhead
反向三角形不等式
\[
\mid\|x\|-\|y\|\mid\leq\|x-y\|,\qquad x,y\in E,
\]

蕴含该断言。 ■

\item

若 $ A \subseteq X $ 且 $ f: X \to Y $ 在 $ x_0 \in A $ 处连续，则 $ f|A: A \to Y $ 在 $ x_0 $ 处连续。这里 $ A $ 具有由 $ X $ 诱导的度量。

\proofhead
这直接由 f 的连续性以及诱导度量的定义得出。
\end{items}

(I) 设 $ M \subseteq X $ 是 $ X $ 的非空子集。对每个 $ x \in X $，

\[
d(x,M):=\inf_{m\in M}d(x,m)
\]

称为 x 到 M 的距离。距离函数

\[
d(\cdot,M):X\to\mathbb{R}~,\quad x\mapsto d(x,M)
\]

是利普希茨连续的。

\proofhead
设 $ x, y \in X $。由三角形不等式，对每个 $ m \in M $ 有 $ d(x, m) \leq d(x, y) + d(y, m) $。由于对所有 $ m \in M $ 有 $ d(x, M) \leq d(x, m) $，这蕴含
\[
d(x,M)\leq d(x,y)+d(y,m)~,\qquad m\in M~.
\]

对所有 $ m \in M $ 取下确界，得

\[
d(x,M)\leq d(x,y)+d(y,M).
\]

将此等式与 x 和 y 互换后所得的同一等式结合，得

\[
|d(x,M)-d(y,M)|\leq d(x,y),
\]

这证明了 $ d(\cdot, M) $ 的利普希茨连续性。 ■

\begin{items}
\item[（m）]

对任何内积空间 $ (E, (\cdot|\cdot)) $，标量积 $ (\cdot|\cdot): E \times E \to \mathbb{K} $ 是连续的。

\proofhead
设 $ (x, y), (x_0, y_0) \in E \times E $，且 $ \varepsilon \in (0, 1) $。由三角形不等式与柯西-施瓦茨不等式，得
\begin{equation*}
\begin{aligned}
\left|(x|y)-(x_{0}|y_{0})\right|&\leq\left|(x-x_{0}|y)\right|+\left|(x_{0}|y-y_{0})\right|\\&\leq\|x-x_{0}\|\|y\|+\|x_{0}\|\|y-y_{0}\|\\&\leq d\big((x,y),(x_{0},y_{0})\big)\left(\|y\|+\|x_{0}\|\right)\\&\leq d\big((x,y),(x_{0},y_{0})\big)\left(\|x_{0}\|+\|y_{0}\|+\|y-y_{0}\|\right),
\end{aligned}
\end{equation*}

其中 $d$ 是积度量。令 $M := \max\{1, \|x_0\|, \|y_0\|\}$，且 $\delta := \varepsilon/(1 + 2M)$。则对所有 $(x, y) \in \mathbb{B}_{E \times E}((x_0, y_0), \delta)$，由上述不等式得

\[
\left|\left(x\mid y\right)-\left(x_{0}\mid y_{0}\right)\right|<\delta(2M+\delta)<\varepsilon,
\]

这证明了标量积在点 $ (x_{0}, y_{0}) $ 处的连续性。 ■

\item[（n）]

设 $E$ 与 $F$ 为赋范向量空间，且 $X \subseteq E$。则 $f: X \to F$ 在 $x_0 \in X$ 处的连续性与 $E$ 及 $F$ 上等价范数的选取无关。

\proofhead
这容易由推论 1.2 得出。
\item[（o）]

度量空间 $X$ 与 $Y$ 之间的函数 $f$ 是等距的（或称为等距映射），若对所有 $x, x' \in X$ 有 $d(f(x), f(x')) = d(x, x')$，即若 $f$ “保持距离”。显然，这样的函数是利普希茨连续的，并且是从 $X$ 到其像 $f(X)$ 的双射。若 $E$ 与 $F$ 为赋范向量空间且 $T: E \to F$ 线性，则 $T$ 等距当且仅当对所有 $x \in E$ 有 $\|Tx\| = \|x\|$。若此外 $T$ 还是满射，则 $T$ 是从 $E$ 到 $F$ 的等距同构，且 $T^{-1}$ 也是等距的。

\end{items}

\end{example}

\subsection*{序列连续性}

邻域概念对连续性的定义和序列收敛的定义都是核心的。这提示我们，函数的连续性可以用序列来定义：度量空间 $ X $ 与 $ Y $ 之间的函数 $ f: X \to Y $ 在 $ x \in X $ 处称为 \textbf{序列连续的}，若对 $ X $ 中每个满足 $ \lim x_k = x $ 的序列 $ (x_k) $，都有 $ \lim f(x_k) = f(x) $。

\begin{theorem}

（序列准则） 设 $X, Y$ 为度量空间。则函数 $f: X \to Y$ 在 $x$ 处连续，当且仅当它在 $x$ 处序列连续。

\proofhead
‘$\Rightarrow$’ 设 $(x_k)$ 是 $X$ 中满足 $x_k \to x$ 的序列。设 $V$ 是 $Y$ 中 $f(x)$ 的邻域。由假设，存在 $X$ 中 $x$ 的邻域 $U$，使得 $f(U) \subseteq V$。由于 $x_k \to x$，存在某个 $N \in \mathbb{N}$，使得对所有 $k \geq N$ 有 $x_k \in U$。因此对所有 $k \geq N$ 有 $f(x_k) \in V$，即 $f(x_k)$ 收敛到 $f(x)$。

‘⇐’ 反设 $f$ 序列连续但在 $x$ 处不连续。则存在 $f(x)$ 的邻域 $V$，使得 $x$ 的任何邻域 $U$ 都不满足 $f(U) \subseteq V$。特别地，我们有
\[
f\left(\mathbb{B}(x,1/k)\right)\cap V^{c}\neq\emptyset,\qquad k\in\mathbb{N}^{\times}.
\]

因此，对每个 $ k \in \mathbb{N}^\times $，我们可取某个 $ x_k \in X $，使得 $ d(x, x_k) < 1/k $ 且 $ f(x_k) \notin V $。按构造，$ (x_k) $ 收敛到 $ x $，但 $ (f(x_k)) $ 不收敛到 $ f(x) $。这与 $ f $ 的序列连续性矛盾。 $ \blacksquare $

\end{theorem}

设 $f: X \to Y$ 是度量空间之间的连续函数。则对 $X$ 中任何收敛序列 $(x_{k})$，有

\[
\lim f(x_{k})=f(\lim x_{k}).
\]

因此人们说“连续函数与取极限运算可交换”。

\subsection*{连续函数的加法与乘法}

定理 1.4 使我们能够把关于收敛序列的定理应用于连续函数。为此，先引入几个定义是有益的。

设 $M$ 为任意集合，$F$ 为向量空间。设 $f$ 与 $g$ 是满足 $\mathrm{dom}(f)$, $\mathrm{dom}(g) \subseteq M$ 且取值于 $F$ 的函数。则 $f$ 与 $g$ 的 $\text{和}$ 是由下式定义的函数 $f + g$：

\[
f+g\colon\operatorname{d o m}(f+g):=\operatorname{d o m}(f)\cap\operatorname{d o m}(g)\to F~,\quad x\mapsto f(x)+g(x)~.
\]

类似地，对 $ \lambda \in \mathbb{K} $，我们定义 $ \lambda f $ 为\footnote{若 $ f $ 与 $ g $ 定义在整个 $ M $ 上，则 $ f + g $ 与 $ \lambda f $ 的定义与例 I.12.3(e) 中的定义一致。}

\[
\lambda f:\operatorname{dom}(f)\to F,\quad x\mapsto\lambda f(x).
\]

最后，在特殊情形 $F = \mathbb{K}$ 中，我们令

\[
\mathrm{dom}(f\cdot g):=\mathrm{dom}(f)\cap\mathrm{dom}(g),
\]
\[
\mathrm{dom}(f/g):=\mathrm{dom}(f)\cap\left\{x\in\mathrm{dom}(g);g(x)\neq0\right\},
\]

并将 f 与 g 的乘积与商定义为

\[
f\cdot g:\mathrm{dom}(f\cdot g)\to\mathbb{K},\quad x\mapsto f(x)\cdot g(x)
\]

以及

\[
f/g\colon\operatorname{d o m}(f/g)\to\mathbb{K}~,\quad x\mapsto f(x)/g(x)\enspace.
\]

\begin{proposition}

假设 X 是度量空间，F 是赋范向量空间，且

\[
f:\operatorname{dom}(f)\subseteq X\to F~,\quad g:\operatorname{dom}(g)\subseteq X\to F
\]

都在 $ x_0 \in \mathrm{dom}(f) \cap \mathrm{dom}(g) $ 处连续。

\begin{subitems}
\item

$ f + g $ 与 $ \lambda f $ 在 $ x_{0} $ 处连续。

\item

若 $ F = \mathbb{K} $，则 $ f \cdot g $ 在 $ x_0 $ 处连续。

\item

若 $ F = \mathbb{K} $ 且 $ g(x_0) \neq 0 $，则 $ f/g $ 在 $ x_0 $ 处连续。

\proofhead
这些断言由定理 1.4 的序列准则、命题 II.2.2 与注 II.3.1(c)，以及命题 II.2.4(ii)、II.2.6 与例 1.3(d) 得出。
\end{subitems}

\end{proposition}

\qyxffntextdef{1}{在不致引起混淆时，我们常把 $ C(X, \mathbb{K}) $ 写作 $ C(X) $。}
\begin{corollary}

\begin{subitems}
\item

有理函数是连续的。

\item

$n$ 元多项式是连续的（在 $\mathbb{K}^{n}$ 上）。

\item

$ C(X,F) $ 是 $ F^X $ 的子空间，即从 X 到 F 的连续函数\qyxffnmark{1}构成的向量空间。

\proofhead
断言 (i) 与 (iii) 是命题 1.5 的直接结论。对于 (ii)，还需要例 1.3(h)。
\end{subitems}

\end{corollary}

\begin{proposition}

设 $ a = \sum a_k X^k $ 是收敛半径 $ \rho_a $ 为正的幂级数。则由 $ a $ 表示的函数 $ \underline{a} $ 在 $ \rho_a $ $ \mathbb{B} $ 上连续。

\proofhead
设 $ x_0 \in \rho_a \mathbb{B}_{\mathbb{C}} $，$ \varepsilon > 0 $，且 $ |x_0| < r < \rho_a $。由定理 II.9.2，级数 $ \sum |a_k| r^k $ 收敛，故存在某个 $ K \in \mathbb{N} $，使得
\begin{equation}
\sum_{k=K+1}^{\infty}\left|a_{k}\right|r^{k}<\varepsilon/4.
\end{equation}

因此，对 $ |x| \leq r $，我们有

\begin{equation}
\begin{aligned}
|\underline{a}(x)-\underline{a}(x_{0})|&\leq\left|\sum_{k=0}^{K}a_{k} x^{k}-\sum_{k=0}^{K}a_{k} x_{0}^{k}\right|+\sum_{k=K+1}^{\infty}|a_{k}|\left|x\right|^{k}+\sum_{k=K+1}^{\infty}|a_{k}|\left|x_{0}\right|^{k}\\&\leq|p(x)-p(x_{0})|+2\sum_{k=K+1}^{\infty}|a_{k}|r^{k},
\end{aligned}
\end{equation}

其中我们已令

\[
p:=\sum_{k=0}^{K}a_{k}X^{k}\in\mathbb{C}[X]~.
\]

由推论 1.6，存在某个 $ \delta \in (0, r - |x_0|) $，使得

\[
|p(x)-p(x_{0})|<\varepsilon/2~,\qquad|x-x_{0}|<\delta~.
\]

结合 (1.2) 与 (1.3)，这蕴含对所有 $ |x - x_0| < \delta $ 有 $ |\underline{a}(x) - \underline{a}(x_0)| < \varepsilon $。由于 $ \mathbb{B}(x_0, \delta) \subseteq \rho_a \mathbb{B}_{\mathbb{C}} $，我们已证明该断言。 $ \blacksquare $

\end{proposition}

下面这个重要定理往往为某些函数的连续性提供简单的证明。我们将在定理之后的例子中加以说明。

\begin{theorem}

（复合的连续性） 设 $X$、$Y$ 与 $Z$ 为度量空间。假设 $f: X \to Y$ 在 $x \in X$ 处连续，$g: Y \to Z$ 在 $f(x) \in Y$ 处连续。则复合 $g \circ f: X \to Z$ 在 $x$ 处连续。

\proofhead
设 $W$ 是 $Z$ 中 $g \circ f(x) = g(f(x))$ 的邻域。由于 $g$ 在 $f(x)$ 处连续，存在 $Y$ 中 $f(x)$ 的邻域 $V$，使得 $g(V) \subseteq W$。由于 $f$ 在 $x$ 处连续，存在 $X$ 中 $x$ 的邻域 $U$，使得 $f(U) \subseteq V$。因此
\[
g\circ f(U)=g\bigl(f(U)\bigr)\subseteq g(V)\subseteq W~,
\]

由此即得断言。 ■

\end{theorem}

\begin{example}

在下文中，X 是度量空间，E 是赋范向量空间。

\begin{items}
\item

设 $ f: X \to E $ 在 $ x_0 $ 处连续。则 $ f $ 的范数

\[
\left\|f\right\|:X\to\mathbb{R},\quad x\mapsto\left\|f(x)\right\|,
\]

在 $ x_{0} $ 处连续

\proofhead
由例 1.3(j)，$ \|\cdot\|: E \to \mathbb{R} $ 是利普希茨连续的。由于 $ \|f\| = \|\cdot\| \circ f $，该断言由定理 1.8 得出。
\item

设 $ g: \mathbb{R} \to X $ 连续。则函数 $ \widehat{g}: E \to X $，$ x \mapsto g(\|x\|) $ 是连续的。

\proofhead
只需注意到 $ \widehat{g} = g \circ \|\cdot\| $ 是连续函数的复合。
\item

定理 1.8 的逆命题不成立，即 $ g \circ f $ 的连续性并不蕴含 f 或 g 连续。

\proofhead
令 $ Z := [-3/2, -1/2] \cup (1/2, 3/2] $ 与 $ I := [-1, 1] $。定义函数 $ f: Z \to \mathbb{R} $ 与 $ g: I \to \mathbb{R} $ 如下：
\begin{equation*}
f(x):=\left\{\begin{array}{l l}{x+1/2\;,}&{\quad x\in\left[-3/2,-1/2\right],}\\ {x-1/2\;,}&{\quad x\in\left(1/2,3/2\right],}\end{array}\right.
\end{equation*}

以及

\begin{equation*}
g(y):=\left\{\begin{array}{l l}{y-1/2\;,}&{\quad y\in\left[-1,0\right],}\\ {y+1/2\;,}&{\quad y\in\left(0,1\right].}\end{array}\right.
\end{equation*}

不难验证 $ f: Z \to \mathbb{R} $ 连续，而 $ g: I \to \mathbb{R} $ 在 0 处不连续，但复合 $ f \circ g = \mathrm{id}_I $ 与 $ g \circ f = \mathrm{id}_Z $ 都连续。我们把构造 $ f $ 也不连续的类似例子的任务留给读者。

\item

函数 $ f: \mathbb{R} \to \mathbb{R} $，$ x \mapsto 1 / \sqrt{1 + x^2} $ 是连续的。

\proofhead
由于 $ 1/\sqrt{1+x^2} = \sqrt{1/(1+x^2)} $，该断言由推论 1.6.(i)、命题 1.5.(iii)、定理 1.8 与例 1.3(a) 得出。
\item

指数函数 $ \exp: \mathbb{C} \to \mathbb{C} $ 是连续的。

\proofhead
这由命题 1.7 与例 II.9.5(a) 得出。 ■
\end{items}

\end{example}

\begin{proposition}

设 $X$ 为度量空间。则从 $X$ 到 $\mathbb{K}^m$ 的函数 $f = (f_1, \ldots, f_m)$ 在 $x$ 处连续，当且仅当对每个 $k$，$f_k: X \to \mathbb{K}$ 在 $x$ 处连续。特别地，$f: X \to \mathbb{C}$ 在 $x$ 处连续，当且仅当 $\operatorname{Re} f$ 与 $\operatorname{Im} f$ 在 $x$ 处连续。

\proofhead
设 $ (x_n) $ 是 X 中满足 $ x_n \to x $ 的序列。由命题 II.3.14，我们有
\[
f(x_{n})\to f(x)\iff f_{k}(x_{n})\to f_{k}(x),~k=1,\ldots,m~.
\]

该断言现由序列准则得出。 ■

\end{proposition}

\subsection*{单侧连续性}

设 $X$ 是 $\mathbb{R}$ 的子集，且 $x_0 \in X$。$\mathbb{R}$ 的序结构使我们能考虑 $x_0$ 的单侧邻域。具体地，对 $\delta > 0$，集合 $X \cap (x_0 - \delta, x_0]$（或 $X \cap [x_0, x_0 + \delta)$）称为 $x_0$ 的左（或右）$\delta$-邻域。

现在设 $Y$ 为度量空间。则 $f: X \to Y$ 在 $x_0$ 处左（或右）连续，若对 $Y$ 中 $f(x_0)$ 的每个邻域 $V$，存在某个 $\delta > 0$，使得 $f(X \cap (x_0 - \delta, x_0]) \subseteq V$（或 $f(X \cap [x_0, x_0 + \delta)) \subseteq V$）。

与命题 1.1 一样，只需考虑 $Y$ 中 $f(x_0)$ 的 $\varepsilon$-邻域，即 $f: X \to Y$ 在 $x_0$ 处左（或右）连续，当且仅当对每个 $\varepsilon > 0$，存在某个 $\delta > 0$，使得对 $x$ 的左（或右）$\delta$-邻域中的所有 $x$，有 $d(f(x_0), f(x)) < \varepsilon$。

显然，连续函数既左连续又右连续。另一方面，单侧连续性并不蕴含连续性，如下面的例子所示。

\begin{example}

\begin{items}
\item

下取整函数 $ \lfloor \cdot \rfloor $：$ \mathbb{R} \to \mathbb{R} $ 在 $ x \in \mathbb{R} \setminus \mathbb{Z} $ 处连续，而在 $ x \in \mathbb{Z} $ 处右连续但不左连续。

\item

函数

\begin{equation*}
\mathrm{s i g n}:\mathbb{R}\to\mathbb{R}\;,\quad x\mapsto\left\{\begin{array}{r l}{-1\;,}&{\quad x<0\;,}\\ {0\;,}&{\quad x=0\;,}\\ {1\;,}&{\quad x>0\;,}\\ \end{array}\right.
\end{equation*}

在 0 处既非左连续也非右连续。
下一个命题把定理 1.4 的序列准则推广到单侧连续函数。

\end{items}

\end{example}

\begin{proposition}

设 $ Y $ 为度量空间，$ X \subseteq \mathbb{R} $，且 $ f: X \to Y $。则下列命题等价：

\begin{subitems}
\item

f 在 $ x \in X $ 处左（或右）连续。

\item

对 X 中每个满足 $ x_n \to x $ 且 $ x_n \leq x $（或 $ x_n \geq x $）的序列 $ (x_n) $，序列 $ (f(x_n)) $ 收敛到 $ f(x) $。

\proofhead
该断言的证明与定理 1.4 的证明类似。
单侧连续性也可用于刻画连续性。
\end{subitems}

\end{proposition}

\begin{proposition}

设 $ Y $ 为度量空间，$ X \subseteq \mathbb{R} $，且 $ f: X \to Y $。则下列命题等价：

\begin{subitems}
\item

f 在 $ x_{0} $ 处连续。

\item

f 在 $ x_{0} $ 处左连续且右连续。

\proofhead
蕴含关系 $\Rightarrow$ 是显然的。
$\Leftarrow$ 设 $\varepsilon > 0$。由 $f$ 在 $x_0$ 处的左连续性与右连续性，存在正数 $\delta^-$ 与 $\delta^+$，使得对所有 $x \in X \cap (x_0 - \delta^-, x_0]$ 与 $x \in X \cap [x_0, x_0 + \delta^+)$，有 $d(f(x), f(x_0)) < \varepsilon$。令 $\delta := \min\{\delta^-, \delta^+\}$。则对所有 $x \in X \cap (x_0 - \delta, x_0 + \delta)$，有 $d(f(x), f(x_0)) < \varepsilon$。因此 $f$ 在 $x_0$ 处连续。
\end{subitems}

\end{proposition}

\exerciseshead

\begin{exercises}
\item

锯齿形函数： $ \mathbb{R} \to \mathbb{R} $ 定义为

\[
\operatorname{z i g z a g}(x):=\left|\left\lfloor x+1/2\right\rfloor-x\right|\mathrm{~,~}\quad x\in\mathbb{R}\mathrm{~,~}
\]

其中 $ \lfloor\cdot\rfloor $ 是下取整函数。画出锯齿形的图像，并证明以下结论：

\begin{items}
\item

锯齿形(x) = |x| 对所有 |x| ≤ 1/2 成立。

\item

锯齿形 $ (x + n) = \text{zigzag}(x) $，$ x \in \mathbb{R} $，$ n \in \mathbb{Z} $。

\item

锯齿形是连续的。

\end{items}

\item

设 $ q \in \mathbb{Q} $。证明函数 $ (0, \infty) \to (0, \infty) $，$ x \mapsto x^q $ 是连续的。\footnote{在第 6 节中，我们将更一般地研究函数 $ x \mapsto x^q $。}（提示：见习题 II.2.7。）

\item

设 $ \varphi: \mathbb{R} \to (-1,1) $，$ x \mapsto x/(1+|x|) $。证明 $ \varphi $ 是双射，且 $ \varphi $ 与 $ \varphi^{-1} $ 连续。

\item

证明或反驳函数

\begin{equation*}
f:\mathbb{Q}\to\mathbb{R},\quad x\mapsto\left\{\begin{array}{ll}0,&x<\sqrt{2},\\1,&x>\sqrt{2},\end{array}\right.
\end{equation*}

是连续的。

\item

设 $d_1$ 与 $d_2$ 是 $X$ 上的度量，且 $X_j := (X, d_j)$，$j = 1, 2$。则若对每个 $x \in X$ 有 $\mathcal{U}_{X_1}(x) \supseteq \mathcal{U}_{X_2}(x)$，即若每点具有比 $d_2$ 邻域更多的 $d_1$ 邻域，则称 $d_1$ 强于 $d_2$。此时也称 $d_2$ 弱于 $d_1$。

证明以下结论：

\begin{items}
\item

$ d_1 $ 强于 $ d_2 $，当且仅当恒等函数 $ i: X_1 \to X_2 $，$ x \mapsto x $ 是连续的。

\item

$ d_1 $ 与 $ d_2 $ 等价，当且仅当 $ d_1 $ 既强于又弱于 $ d_2 $，即对每个 $ x \in X $，$ \mathcal{U}_{X_1}(x) = \mathcal{U}_{X_2}(x) $。

\end{items}

\item

设 $ f: \mathbb{R} \to \mathbb{R} $ 是加法群 ($ \mathbb{R},+ $) 的连续\footnote{可以证明，$ (\mathbb{R}, +) $ 存在不连续的同态（见第 III 卷习题 IX.5.6）。}同态。证明 $ f $ 是线性的，即存在某个 $ a \in \mathbb{R} $，使得 $ f(x) = ax $，$ x \in \mathbb{R} $。（提示：证明对所有 $ q \in \mathbb{Q} $ 有 $ f(q) = qf(1) $，并利用命题 I.10.8。）

\item

设 $ f: \mathbb{R} \to \mathbb{R} $ 定义为

\begin{equation*}
f(x):=\left\{\begin{array}{c c}{-1,}&{x\geq1,}\\ {1/n,}&{1/(n+1)\leq x<1/n,}\\ {0,}&{x\leq0.}\\ \end{array}\right.n\in\mathbb{N}^{\times},
\end{equation*}

f 在何处连续？在何处左连续？在何处右连续？

\item

假设 $X$ 是度量空间，且 $f, g \in \mathbb{R}^X$ 在 $x_0$ 处连续。证明或反驳\footnote{见例 I.4.4(c)。}

\begin{equation}
|f|,\quad f^{+}:=0\lor f,\quad f^{-}:=0\lor\left(-f\right),\quad f\lor g,\quad f\land g
\end{equation}

都在 $ x_{0} $ 处连续。（提示：例 1.3(j) 与习题 I.8.11。）

\item

设 $ f: \mathbb{R} \to \mathbb{R} $ 与 $ g: \mathbb{R} \to \mathbb{R} $ 定义为

\begin{equation*}
f(x):=\left\{\begin{array}{r l}{1\quad,\quad}&{x\text{ 为有理数},}\\ {-1\quad,\quad}&{x\text{ 为无理数},}\end{array}\right.\quad g(x):=\left\{\begin{array}{r l}{x\quad,\quad}&{x\text{ 为有理数},}\\ {-x\quad,\quad}&{x\text{ 为无理数}.}\end{array}\right.
\end{equation*}

函数 $f$、$g$、$|f|$、$|g|$ 与 $f \cdot g$ 在何处连续？

\item

设 $ f: \mathbb{R} \to \mathbb{R} $ 定义为

\begin{equation*}
f(x):=\left\{\begin{array}{c l}{1/n\;,}&{\quad x\in\mathbb{Q}\text{ 且 }x=m/n\text{ 为最简分数}\;,}\\ {0\;,}&{\quad x\in\mathbb{R}\backslash\mathbb{Q}\;.}\end{array}\right.
\end{equation*}
\end{exercises}

\begin{proof}

$f$ 在每个无理数处连续，而在每个有理数处不连续。\footnote{可以证明，不存在从 $ \mathbb{R} $ 到 $ \mathbb{R} $ 的函数，它在每个有理数处连续而在每个无理数处不连续（见习题 V.4.5）。}（提示：对每个 $x \in \mathbb{Q}$，由命题 I.10.11，存在序列 $x_n \in \mathbb{R} \setminus \mathbb{Q}$，使得 $x_n \to x$。因此 $f$ 不能在 $x$ 处连续。

设 $ x \in \mathbb{R} \setminus \mathbb{Q} $，且 $ \varepsilon > 0 $。则只有有限多个 $ n \in \mathbb{N} $ 满足 $ n \leq 1/\varepsilon $。因此存在某个 $ \delta > 0 $，使得任何满足 $ n \leq 1/\varepsilon $ 的 $ q = m/n $ 都不在 $ (x - \delta, x + \delta) $ 中。即对 $ y = m/n \in (x - \delta, x + \delta) $，有 $ f(y) = f(m/n) = 1/n < \varepsilon $。

\end{proof}

\begin{exercises}
\item[11.]

考虑函数

\begin{equation*}
f:\mathbb{R}^{2}\to\mathbb{R},\quad(x,y)\mapsto\left\{\begin{array}{c l}{x y/(x^{2}+y^{2}),}&{\quad(x,y)\neq(0,0),}\\ {0,}&{\quad(x,y)=(0,0),}\\ \end{array}\right.
\end{equation*}

并对固定的 $ x_0 \in \mathbb{R} $，定义

\[
f_{1}\colon\mathbb{R}\to\mathbb{R}~,\quad x\mapsto f(x,x_{0})~,\qquad f_{2}\colon\mathbb{R}\to\mathbb{R}~,\quad x\mapsto f(x_{0},x)~.
\]

证明以下结论：

\begin{items}
\item

$ f_{1} $ 与 $ f_{2} $ 是连续的。

\item

$f$ 在 $\mathbb{R}^2 \setminus \{(0,0)\}$ 上连续，而在 $(0,0)$ 处不连续。（提示：对零序列 $(x_n)$，考虑 $f(x_n, x_n)$。）

\end{items}

\item[12.]

证明从 $ \mathbb{K}^n $ 到 $ \mathbb{K}^m $ 的任何线性函数都是利普希茨连续的。（提示：利用命题 II.3.12 与适当的范数。）

\item[13.]

假设 $V$ 与 $W$ 是赋范向量空间，且 $f: V \to W$ 是从 $(V,+)$ 到 $(W,+)$ 的连续群同态。证明 $f$ 是线性的。（提示：若 $\mathbb{K} = \mathbb{R}$，$x \in V$ 且 $q \in \mathbb{Q}$，则 $f(qx) = qf(x)$。另见习题 6。）

\item[14.]

设 $ (E, (\cdot|\cdot)) $ 是内积空间，且 $ x_0 \in E $。证明函数

\[
E\to\mathbb{K}~,\quad x\mapsto(x\left|x_{0}\right)~,~\quad E\to\mathbb{K}~,\quad x\mapsto(x_{0}\left|x\right)
\]

是连续的。

\item[15.]

设 $ A \in \mathrm{End}(\mathbb{K}^n) $。证明函数

\[
\mathbb{K}^{n}\to\mathbb{K},\quad x\mapsto(A x\mid x)
\]

是连续的。（提示：利用习题 12 与柯西-施瓦茨不等式。）

\item[16.]

设 $ n \in \mathbb{N}^\times $。矩阵 $ A = [a_{jk}] \in \mathbb{K}^{n \times n} $ 的行列式定义为（见习题 I.9.6）

\[
\det A=\sum_{\sigma\in\mathsf{S}_{n}}(\mathrm{sign}\sigma)a_{1\sigma(1)}\cdot\cdots\cdot a_{n\sigma(n)}.
\]

证明函数

\[
\mathbb{K}^{n\times n}\to\mathbb{K}~,\quad A\mapsto\det A
\]

是连续的（见习题 II.3.14）。（提示：利用双射

\begin{equation*}
\mathbb{K}^{m\times n}\to\mathbb{K}^{m n},\quad\left[\begin{array}{c c c}{a_{11},}&{\ldots,}&{a_{1n}}\\ {\vdots}&{}&{\vdots}\\ {a_{m1},}&{\ldots,}&{a_{m n}}\\ \end{array}\right]\mapsto\left(a_{11},\ldots,a_{1n},a_{21},\ldots,a_{m n}\right)
\end{equation*}

来定义 $ \mathbb{K}^{m\times n} $ 上的自然拓扑。）

\item[17.]

设 $X$ 与 $Y$ 为度量空间，且 $f: X \to Y$。对 $x \in X$，函数

\[
\omega_{f}(x,\cdot):(0,\infty)\to\mathbb{R}~,\quad\varepsilon\mapsto\sup_{y,z\in\mathbb{B}(x,\varepsilon)}d\big(f(y),f(z)\big)
\]

称为 f 的连续模。令

\[
\omega_{f}(x):=\inf_{\varepsilon>0}\omega_{f}(x,\varepsilon).
\]
\end{exercises}

\begin{proof}

f 在 x 处连续，当且仅当 $ \omega_{f}(x)=0 $。

\end{proof}

\begin{exercises}
\item[18.]

证明平方根函数 $ w: \mathbb{R}^+ \to \mathbb{R} $，$ x \mapsto \sqrt{x} $ 连续但不利普希茨连续。证明对每个 $ a > 0 $，$ w \mid [a, \infty) $ 是利普希茨连续的。

\end{exercises}

\section{拓扑基础}

为了更深入地理解连续函数，我们在本节引入拓扑空间的一些基本概念。主要结果是定理 2.20，它把连续函数刻画为拓扑空间之间保持结构的函数。

\subsection*{开集}

在下文中，$X := (X, d)$ 是一个度量空间。若 $X$ 的子集 $A$ 的元素 $a$ 存在一个邻域 $U$ 使得 $U \subseteq A$，则称 $a$ 为 $A$ 的内点。若 $A$ 的每个点都是内点，则称集合 $A$ 是开集。

\begin{remark}

\begin{items}
\item

显然，$ a $ 是 $ A $ 的内点当且仅当存在某个 $ \varepsilon > 0 $ 使得 $ \mathbb{B}(a, \varepsilon) \subseteq A $。

\item

A 是开集当且仅当 A 是其每个点的邻域。

\end{items}

\end{remark}

\begin{example}

开球 $ \mathbb{B}(a, r) $

\proofhead
对 $x_0 \in \mathbb{B}(a, r)$，令 $s := d(x_0, a)$。则 $\varepsilon := r - s$ 是正的。对所有 $x \in \mathbb{B}(x_0, \varepsilon)$，我们有
\[
d(x,a)\leq d(x,x_{0})+d(x_{0},a)<\varepsilon+s=r,
\]
\bookfigure[0.17]{figures/bfb6e42ef8a41865.jpg}

因此 $\mathbb{B}(x_0,\varepsilon)$ 包含于 $\mathbb{B}(a,r)$。这表明 $x_0$ 是 $\mathbb{B}(a,r)$ 的内点。$\blacksquare$

\end{example}

\begin{remark}

\begin{items}
\item

“内点”与“开集”这两个概念依赖于周围的度量空间 $X$。有时把这一点明确说出是有用的，即“$a$ 是 $A$ 关于 $X$ 的内点，或“$A$ 在 $X$ 中是开集”。
例如，$\mathbb{R}$ 中的开球，即开区间 $J$，由前例可知在 $\mathbb{R}$ 中是开集。然而，若我们把 $\mathbb{R}$ 视为嵌入到 $\mathbb{R}^2$ 中，则 $J$ 在 $\mathbb{R}^2$ 中不是开集。

\item

设 $ X = (X, \|\cdot\|) $ 是赋范向量空间，$ \|\cdot\|_1 $ 与 $ \|\cdot\| $ 是 $ X $ 上的等价范数。则由注 II.3.13(d)，

\[
A\text{ 是 }(X,\|\cdot\|)\text{ 中的开集}\Longleftrightarrow A\text{ 是 }(X,\|\cdot\|_{1})\text{ 中的开集}.
\]

因此，若 A 关于某个特定范数是开集，则它关于所有等价范数都是开集。

\item

由例 2.2 可得，度量空间中的每个点都有开邻域。

\end{items}

\end{remark}

\begin{proposition}

设 $ \mathcal{T} := \{O \subseteq X; O \text{ 是开集}\} $ 是开集族。

\begin{subitems}
\item

$ \emptyset, X \in \mathcal{T} $。

\item

若对所有 $ \alpha \in \mathbb{A} $ 有 $ O_\alpha \in \mathcal{T} $，则 $ \bigcup_\alpha O_\alpha \in \mathcal{T} $。也就是说，开集的任意并是开集。

\item

若 $ O_0, \ldots, O_n \in \mathcal{T} $，则 $ \bigcap_{k=0}^{n} O_k \in \mathcal{T} $。也就是说，开集的有限交是开集。

\proofhead
(i) 显然 $ X $ 在 $ \mathcal{T} $ 中，而由注 I.2.1(a)，$\emptyset$ 也是开集。
\end{subitems}

\begin{subitems}
\item[（ii）]

设 $ \mathbf{A} $ 是指标集，对所有 $ \alpha \in \mathbf{A} $ 有 $ O_\alpha \in \mathcal{T} $，且 $ x_0 $ 是 $ \bigcup_\alpha O_\alpha $ 的一个点。则存在某个 $ \alpha_0 \in \mathbf{A} $ 使得 $ x_0 \in O_{\alpha_0} $。由于 $ O_{\alpha_0} $ 是开集，存在 $ x_0 $ 在 $ X $ 中的某个邻域 $ U $ 使得 $ U \subseteq O_{\alpha_0} \subseteq \bigcup_\alpha O_\alpha $。因此 $ \bigcup_\alpha O_\alpha $ 是开集。

\item[（iii）]

设 $O_0, \ldots, O_n \in \mathcal{T}$ 且 $x_0 \in \bigcap_{k=0}^n O_k$。则存在正数 $\varepsilon_k$ 使得对 $k = 0, \ldots, n$ 有 $\mathbb{B}(x_0, \varepsilon_k) \subseteq O_k$。令 $\varepsilon := \min\{\varepsilon_0, \ldots, \varepsilon_n\} > 0$。则 $\mathbb{B}(x_0, \varepsilon)$ 包含于每个 $O_k$，从而 $\mathbb{B}(x_0, \varepsilon) \subseteq \bigcap_{k=0}^n O_k$。
命题 2.4 的性质 (i)–(iii) 涉及集合运算 $ \bigcup $ 与 $ \bigcap $，但不涉及度量。这提示我们对度量空间的概念作如下推广：设 $ M $ 是一个集合，$ \mathcal{T} \subseteq \mathcal{P}(M) $ 是满足 (i)–(iii) 的子集族。则 $ \mathcal{T} $ 称为 $ M $ 上的一个拓扑，$ \mathcal{T} $ 的元素称为关于 $ \mathcal{T} $ 的开集。最后，二元组 $ (M, \mathcal{T}) $ 称为拓扑空间。

\end{subitems}

\end{proposition}

\begin{remark}

\begin{items}
\item

设 $ \mathcal{T} \subseteq \mathcal{P}(X) $ 是命题 2.4 中的集合族。则 $ \mathcal{T} $ 称为由度量 $ d $ 在 $ X $ 上诱导的拓扑。若 $ X $ 是带有由范数诱导的度量的赋范向量空间，则 $ \mathcal{T} $ 称为范数拓扑。

\item

设 $(X,\|\cdot\|)$ 是赋范向量空间，$\|\cdot\|_1$ 是 $X$ 上与 $\|\cdot\|$ 等价的范数。设 $\mathcal{T}_{\|\cdot\|}$ 与 $\mathcal{T}_{\|\cdot\|_1}$ 分别是由 $(X,\|\cdot\|)$ 与 $(X,\|\cdot\|_1)$ 诱导的范数拓扑。由注 2.3(b)，$\mathcal{T}_{\|\cdot\|}$ 与 $\mathcal{T}_{\|\cdot\|_1}$ 重合，也就是说，等价范数在 $X$ 上诱导相同的拓扑。

\end{items}

\end{remark}

\subsection*{闭集}

度量空间 X 的子集 A 称为在 X 中是闭集，若 $ A^{c} $ 在 X 中是开集\footnote{注意，$ A $ 不是开集并\textbf{不}蕴含 $ A $ 是闭集。例如，设 $ X := \mathbb{R} $ 与 $ A := [0, 1) $。则 $ A $ 在 $ \mathbb{R} $ 中既不是开集也不是闭集。}。

\begin{proposition}

\begin{subitems}
\item

$ \emptyset $ 与 X 是闭集。

\item

闭集的任意交是闭集。

\item

闭集的有限并是闭集。

\proofhead
这些断言容易由命题 2.4 与命题 I.2.7(iii) 得出。
\end{subitems}

\end{proposition}

\begin{remark}

\begin{items}
\item

开集的无限交不必是开集。

\proofhead
例如，在 $ \mathbb{R} $ 中有 $ \bigcap_{n=1}^{\infty} \mathbb{B}(0,1/n) = \{0\} $。$ \blacksquare $
\item

闭集的无限并不必是闭集。

\proofhead
例如，在 $ \mathbb{R} $ 中有 $ \bigcup_{n=1}^{\infty}\left[\mathbb{B}(0,1/n)\right]^{c}=\mathbb{R}^{\times} $。$ \blacksquare $
\end{items}

设 $ A \subseteq X $ 且 $ x \in X $。若 $ x $ 在 $ X $ 中的每个邻域都与 $ A $ 有非空交集，则称 $ x $ 为 $ A $ 的聚点。若 $ x $ 在 $ X $ 中的每个邻域都含有 $ A $ 中异于 $ x $ 的点，则称元素 $ x \in X $ 为 $ A $ 的\textbf{极限点}。最后我们置

\[
\overline{A}:=\left\{x\in X;x\text{ 是 }A\text{ 的聚点}\right\}.
\]

显然 A 的每个元素以及 A 的每个极限点都是 A 的聚点。事实上 $ \overline{A} $ 是 A 与 A 的极限点集合之并。

\end{remark}

\begin{proposition}

设 A 是度量空间 X 的子集。

\begin{subitems}
\item

$ A \subseteq \overline{A} $。

\item

$ A = \overline{A} \Leftrightarrow A $ 是闭集。

\proofhead
断言 (i) 是显然的。
\end{subitems}

(ii) “$\Rightarrow$” 设 $x \in A^c = (\overline{A})^c$。由于 $x$ 不是 $A$ 的聚点，存在某个 $U \in \mathcal{U}(x)$ 使得 $U \cap A = \emptyset$。于是 $U \subseteq A^c$，即 $x$ 是 $A^c$ 的内点。从而 $A^c$ 是开集，$A$ 在 $X$ 中是闭集。

“$\Leftarrow$” 设 $A$ 在 $X$ 中是闭集。则 $A^c$ 在 $X$ 中是开集。对任何 $x \in A^c$，存在某个 $U \in \mathcal{U}(x)$ 使得 $U \subseteq A^c$。这意味着 $U$ 与 $A$ 不交，从而 $x$ 不是 $A$ 的聚点，即 $x \in (\overline{A})^c$。因此我们证明了包含关系 $A^c \subseteq (\overline{A})^c$，它等价于 $\overline{A} \subseteq A$。结合 (i)，这推出 $\overline{A} = A$。

集合 A 的极限点是 A 中某些序列的极限。

\end{proposition}

\begin{proposition}

X 的元素 x 是 A 的极限点当且仅当存在 $ A \setminus \{x\} $ 中的序列 $ (x_k) $ 收敛于 x。

\proofhead
设 $x$ 是 $A$ 的极限点。对每个 $k \in \mathbb{N}^\times$，在 $\mathbb{B}(x, 1/k)$ 中选取某个元素 $x_k \neq x$。则 $(x_k)$ 是 $A\backslash\{x\}$ 中的序列，且 $x_k \to x$。

反之，设 $ (x_k) $ 是 $ A \setminus \{x\} $ 中的序列且 $ x_k \to x $。则对 $ x $ 的每个邻域 $ U $，存在某个 $ k \in \mathbb{N} $ 使得 $ x_k \in U $。这意味着 $ x_k \in U \cap (A \setminus \{x\}) $。因此 $ x $ 的每个邻域都含有 $ A $ 中异于 $ x $ 的元素。$ \blacksquare $
\end{proposition}

\begin{corollary}

X 的元素 x 是 A 的聚点当且仅当存在 A 中的序列 $ (x_k) $ 使得 $ x_k \to x $。

\proofhead
若 $x$ 是极限点，则断言由命题 2.9 得出。否则，若 $x$ 是 $A$ 的聚点但不是 $A$ 的极限点，则存在 $x$ 的邻域 $U$ 使得 $U \cap A = \{x\}$。于是 $x$ 在 $A$ 中，而常序列 $(x_k)$（对所有 $k \in \mathbb{N}$ 有 $x_k = x$）具有所要求的性质。

我们现在可以用收敛序列来刻画闭集。
\end{corollary}

\begin{proposition}

对 $ A \subseteq X $，下列陈述等价：

\begin{subitems}
\item

A 是闭集。

\item

A 包含它的所有极限点。

\item

A 中每个在 X 中收敛的序列，其极限都在 A 中。

\proofhead
“(i) $ \Rightarrow $ (ii)” A 的任何极限点也是聚点，从而包含于 $ \overline{A} $。由 (i) 与命题 2.8，$ A = \overline{A} $，从而所有极限点都在 $ A $ 中。
$ ^{'(ii) \Rightarrow (iii)'} $ 设 $ (x_k) $ 是 $ A $ 中的序列且 $ x_k \to x $（在 $ X $ 中）。则由推论 2.10，$ x $ 是 $ A $ 的聚点。这意味着，要么 $ x $ 在 $ A $ 中，要么 $ x $ 是 $ A $ 的极限点，因此由假设，$ x $ 在 $ A $ 中。
“(iii)⇒(i)” 这一蕴含由命题 2.8 与推论 2.10 得出。■
\end{subitems}

\end{proposition}

\subsection*{集合的闭包}

设 A 是度量空间 X 的子集。按下式定义 A 的闭包：

\[
\mathrm{cl}(A):=\mathrm{cl}_{X}(A):=\bigcap_{B\in M}B
\]

其中

\[
M:=\{B\subseteq X\;;\;B\supseteq A\text{ 且 }B\text{ 是 }X\text{ 中的闭集}\}
\]

由于 $X$ 是闭集且包含 $A$，集合 $M$ 非空，此定义是有意义的。由命题 2.6(ii)，$\operatorname{cl}(A)$ 是闭集。由于 $A \subseteq \operatorname{cl}(A)$，$A$ 的闭包正是包含 $A$ 的最小闭集，也就是说，任何包含 $A$ 的闭集也包含 $\operatorname{cl}(A)$。

在下一个命题中我们证明，$A$ 的闭包就是 $A$ 的所有聚点之集，即 $\overline{A} = \mathrm{cl}(A)$。

\begin{proposition}

设 A 是度量空间 X 的子集。则 $ \overline{A} = \mathrm{cl}(A) $。

证明

\begin{subitems}
\item

首先我们证明 $ \overline{A} \subseteq \mathrm{cl}(A) $。若 $ \mathrm{cl}(A) = X $，则陈述显然成立。反之设 $ \mathrm{cl}(A) \neq X $ 且 $ x \in U := (\mathrm{cl}(A))^c $。由于 $ \mathrm{cl}(A) $ 是
闭集，$U$ 是开集，从而是 $x$ 的邻域。由 $A \subseteq \mathrm{cl}(A)$ 可得 $A$ 与 $U$ 不交，即 $x$ 不是 $A$ 的聚点。这推出 $\left(\mathrm{cl}(A)\right)^c \subseteq (\overline{A})^c$，从而 $\overline{A} \subseteq \mathrm{cl}(A)$。

\item

现在我们证明相反的包含关系 $ \operatorname{cl}(A) \subseteq \overline{A} $。同样，情形 $ \overline{A} = X $ 是平凡的。若 $ x \notin \overline{A} $，则存在 $ x $ 的开邻域 $ U $ 使得 $ U \cap A = \emptyset $，即 $ A $ 包含于闭集 $ U^c $。于是 $ x \in U \subseteq (\operatorname{cl}(A))^c $，我们证明了 $ (\overline{A})^c \subseteq (\operatorname{cl}(A))^c $，等价地有 $ \operatorname{cl}(A) \subseteq \overline{A} $。
下面的推论汇集了“$ \overline{A} $ 是包含 A 的最小闭集”这一事实的一些简单推论。

\end{subitems}

\end{proposition}

\begin{corollary}

设 A 与 B 是 X 的子集。

\begin{subitems}
\item

$ A \subseteq B \Rightarrow \overline{A} \subseteq \overline{B} $。

\item

$ \overline{(A)}=\overline{A}. $

\item

$ \overline{A \cup B} = \overline{A} \cup \overline{B} $。

\proofhead
断言 (i) 与 (ii) 直接由命题 2.12 得出。
为证明 (iii)，我们首先注意到，由命题 2.6(iii) 与 2.12，$ \overline{A \cup B} $ 是闭集。由于 $ \overline{A \cup B} $ 包含 $ A \cup B $，命题 2.12 推出 $ \overline{A \cup B} \subseteq \overline{A} \cup \overline{B} $。另一方面 $ \overline{A \cup B} $ 也是闭集。由于 $ A \subseteq \overline{A \cup B} $ 且 $ B \subseteq \overline{A \cup B} $，我们得到包含关系 $ \overline{A} \subseteq \overline{A \cup B} $ 与 $ \overline{B} \subseteq \overline{A \cup B} $。合起来，这些推出 $ \overline{A} \cup \overline{B} \subseteq \overline{A \cup B} $。
这个推论表明，函数 $ h: \mathcal{P}(X) \to \mathcal{P}(X) $，$ A \mapsto \overline{A} $ 是递增的且幂等的，即 $ h \circ h = h $。
\end{subitems}

\end{corollary}

\subsection*{集合的内部}

闭集、聚点与闭包之间的关系，对于开集有一个平行的对应物，我们在本节中加以描述。取代闭包角色的是 A 的内部，定义为

\[
\mathrm{i n t}(A):=\mathrm{i n t}_{X}(A):=\bigcup\{O\subseteq A\;;\;O\text{ 是 }X\text{ 中的开集}\}.
\]

显然 int(A) 是 A 的子集，且由命题 2.4(ii)，int(A) 是开集。因此 int(A) 是 A 的最大开子集。聚点的角色由内点承担，我们定义

\[
\mathring{A}:=\{a\in A\;;\;a\text{ 是 }A\text{ 的内点}\}
\]

于是，与命题 2.12 相对应，我们有：

\begin{proposition}

设 A 是度量空间 X 的子集。则 $ \mathring{A} = \operatorname{int}(A) $

证明

\begin{subitems}
\item

对每个 $ a \in \mathring{A} $，存在 a 的开邻域 $ U $ 使得 $ U \subseteq A $。于是 $ a \in U \subseteq \mathrm{int}(A) $，从而我们证明了 $ \mathring{A} \subseteq \mathrm{int}(A) $。

\item

反之，设 $a \in \mathrm{int}(A)$。则存在 $A$ 的开子集 $O$ 使得 $a \in O$。于是 $O$ 是 $a$ 的包含于 $A$ 的邻域，即 $a$ 是 $A$ 的内点。因此我们有包含关系 $\mathrm{int}(A) \subseteq \mathring{A}$。
下面的推论是这个命题的直接推论。

\end{subitems}

\end{proposition}

\begin{corollary}

设 A 与 B 是 X 的子集。

\begin{subitems}
\item

$ A \subseteq B \Rightarrow \mathring{A} \subseteq \mathring{B} $。

\item

$ (\mathring{A})^{\circ} = \mathring{A} $

\item

A 是开集 $ \Leftrightarrow A = \mathring{A} $。
与闭包的情形类似，函数 $ \mathcal{P}(X) \to \mathcal{P}(X) $，$ A \mapsto \mathring{A} $ 是递增的且幂等的。

\end{subitems}

\end{corollary}

\subsection*{集合的边界}

直观上，我们期望平面中圆盘的边界是包围它的圆周。利用开集与闭集的概念，可以把边界应当是什么这一概念精确化。具体地，对度量空间 $X$ 的子集 $A$，$A$ 的（拓扑）边界定义为 $\partial A := \overline{A} \backslash \mathring{A}$。例如，$X$ 的边界是空的，即 $\partial X = \emptyset$。

\begin{proposition}

设 A 是 X 的子集。

\begin{subitems}
\item

$ \partial A $ 是闭集。

\item

$x$ 在 $\partial A$ 中当且仅当 $x$ 的每个邻域都与 $A$ 和 $A^{c}$ 两者有非空交集。

\proofhead
这些断言直接由 $ \partial A = \overline{A} \cap (\mathring{A})^{c} $ 得出。■
\end{subitems}

\end{proposition}

\subsection*{豪斯多夫条件}

下面的命题表明，在度量空间中，任何两个不同的点都有不交的邻域。

\begin{proposition}

设 $x, y \in X$ 满足 $x \neq y$。则存在 $x$ 的邻域 $U$ 与 $y$ 的邻域 $V$ 使得 $U \cap V = \emptyset$。

\proofhead
由于 $ x \neq y $，我们有 $ \varepsilon := d(x, y)/2 > 0 $。令 $ U := \mathbb{B}(x, \varepsilon) $ 与 $ V := \mathbb{B}(y, \varepsilon) $。假设 $ U \cap V \neq \emptyset $，于是存在某个 $ z \in U \cap V $。则由三角形不等式，
\[
2\varepsilon=d(x,y)\leq d(x,z)+d(z,y)<\varepsilon+\varepsilon=2\varepsilon,
\]

矛盾。因此 U 与 V 不交。

命题 2.17 的断言称为豪斯多夫条件。为证明这个条件，我们本质地利用了度量的存在性。事实上存在命题 2.17 不成立的（非度量）拓扑空间。这种拓扑空间的一个简单例子出现在习题 10 中。

豪斯多夫条件的一个简单推论是

\[
\bigcap\left\{U\;;\;U\in\mathcal{U}_{X}(x)\right\}=\left\{x\right\}\;,\qquad x\in X\;,
\]

这意味着有足够多的邻域来区分度量空间中的点。

\end{proposition}

\qyxffntextdef{1}{这也是命题 2.11(iii) 的一个简单推论（另见注 2.29(d)）。}
\begin{corollary}

度量空间的任何单元素子集都是闭集。

证明\qyxffnmark{1} 固定 $x \in X$。若 $X = \{x\}$，则断言由命题 2.6(i) 得出。否则，若 $y \in \{x\}^c$，则由命题 2.17，存在 $x$ 的邻域 $U$ 与 $y$ 的邻域 $V$ 使得 $U \cap V = \emptyset$。特别地，$\{x\} \cap V \subseteq U \cap V = \emptyset$，从而 $V \subseteq \{x\}^c$。因此 $\{x\}^c$ 是开集。■

\end{corollary}

\subsection*{例子}

我们用一些例子来说明这些新概念，这些例子特别表明，先前定义的“开区间”、“闭区间”、“开球”与“闭球”这些概念与拓扑概念是一致的。

\begin{example}

\begin{items}
\item

开区间 $ (a, b) \subseteq \mathbb{R} $ 在 $ \mathbb{R} $ 中是开集。

\item

闭区间 $ [a, b] \subseteq \mathbb{R} $ 在 $ \mathbb{R} $ 中是闭集。

\item

设 $ I \subseteq \mathbb{R} $ 是区间，$ a := \inf I $ 且 $ b := \sup I $。则

\begin{equation*}
\partial I=\left\{\begin{array}{r l}{\emptyset\quad,}&{\quad I=\mathbb{R}\text{ 或 }I=\emptyset\quad,}\\ {\{a\}\quad,}&{\quad a\in\mathbb{R}\text{ 且 }b=\infty\quad,}\\ {\{b\}\quad,}&{\quad b\in\mathbb{R}\text{ 且 }a=-\infty\quad,}\\ {\{a,b\}\quad,}&{\quad-\infty<a<b<\infty\quad,}\\ {\{a\}\quad,}&{\quad a=b\in\mathbb{R}\quad.}\end{array}\right.
\end{equation*}
\item

闭球 $ \bar{\mathbb{B}}(x,r) $ 是闭集。

\proofhead
若 $X = \bar{\mathbb{B}}(x, r)$ 则无须证明。故设 $\bar{\mathbb{B}}(x, r) \neq X$ 且 $y$ 不在 $\bar{\mathbb{B}}(x, r)$ 中，即 $\varepsilon := d(x, y) - r > 0$。则对 $z \in \mathbb{B}(y, \varepsilon)$，由反向三角形不等式得
\[
d(x,z)\geq d(x,y)-d(y,z)>d(x,y)-\varepsilon=r.
\]

因此球 $ \mathbb{B}(y,\varepsilon) $ 包含于 $ (\bar{\mathbb{B}}(x,r))^c $。由于这对所有 $ y\in(\bar{\mathbb{B}}(x,r))^c $ 成立，$ (\bar{\mathbb{B}}(x,r))^c $ 是开集。■

\item

在任何度量空间中，对 $ r \geq 0 $ 有 $ \overline{\mathbb{B}(x,r)} \subseteq \bar{\mathbb{B}}(x,r) $。若 $ X $ 是赋范向量空间\footnote{存在度量空间，其中 $ \overline{\mathbb{B}(x,r)} $ 是 $ \bar{\mathbb{B}}(x,r) $ 的真子集，如习题 3 所示。} 且 $ r > 0 $，则 $ \overline{\mathbb{B}(x,r)} = \bar{\mathbb{B}}(x,r) $。

\proofhead
第一个断言是 (d) 与命题 2.12 的推论。
对第二个断言，设 $X$ 是赋范向量空间且 $r > 0$。只需证明包含关系 $\bar{\mathbb{B}}(x, r) \subseteq \overline{\mathbb{B}}(x, r)$。反设 $\overline{\mathbb{B}}(x, r) \subset \bar{\mathbb{B}}(x, r)$。选取某个 $y \in \bar{\mathbb{B}}(x, r) \setminus \overline{\mathbb{B}}(x, r)$，并注意到 $d(y, x) = \|y - x\| = r > 0$，因此 $x \neq y$。对 $\varepsilon \in (0, 1)$，定义
\[
x_{\varepsilon}:=x+(1-\varepsilon)(y-x)=\varepsilon x+(1-\varepsilon)y.
\]

则 $ \|x - x_\varepsilon\| = (1 - \varepsilon) \|y - x\| = (1 - \varepsilon) r < r $ 且 $ \|y - x_\varepsilon\| = \varepsilon \|x - y\| = \varepsilon r > 0 $。现在设 $ (\varepsilon_k) $ 是 $ (0,1) $ 中的零序列，且对所有 $ k \in \mathbb{N} $ 有 $ x_k := x_{\varepsilon_k} $。则 $ (x_k) $ 是 $ \mathbb{B}(x, r) $ 中的序列且 $ x_k \to y $。由命题 2.10，$ y $ 是 $ \mathbb{B}(x, r) $ 的聚点，即 $ y \in \overline{\mathbb{B}(x, r)} $。但这与 $ y $ 的选取矛盾。■

\end{items}

\bookfigure[0.15]{figures/d14685d45f078c06.jpg}

\begin{items}
\item[（f）]

在任何赋范向量空间 X 中，

\[
\partial\mathbb{B}(x,r)=\partial\bar{\mathbb{B}}(x,r)=\left\{y\in X;\|x-y\|=r\right\}.
\]

\proofhead
这由 (e) 得出。
\item[（g）]

$ n $-球面 $ S^n := \{x \in \mathbb{R}^{n+1} ; |x| = 1\} $ 在 $ \mathbb{R}^{n+1} $ 中是闭集。

\proofhead
由于 $ S^n = \partial \mathbb{B}^{n+1} $，断言由命题 2.16(i) 得出。■
\end{items}

\end{example}

\subsection*{连续函数的刻画}

我们现在给出先前预告的本节主要结果。

\begin{theorem}

设 $ f: X \to Y $ 是度量空间 X 与 Y 之间的函数。则下列陈述等价：

\begin{subitems}
\item

f 是连续的。

\item

对 Y 中每个开集 O，$ f^{-1}(O) $ 在 X 中是开集。

\item

对 Y 中每个闭集 A，$ f^{-1}(A) $ 在 X 中是闭集。

\proofhead
“(i)⇒(ii)” 设 $O \subseteq Y$ 是开集。若 $f^{-1}(O) = \emptyset$，则断言由命题 2.4(i) 得出。故设 $f^{-1}(O) \neq \emptyset$。由于 $f$ 连续，对每个 $x \in f^{-1}(O)$，存在 $x$ 在 $X$ 中的开邻域 $U_x$ 使得 $f(U_x) \subseteq O$。这推出
\[
x\in U_{x}\subseteq f^{-1}(O),\qquad x\in f^{-1}(O),
\]

由此我们得到

\[
\bigcup_{x\in f^{-1}(O)}U_{x}=f^{-1}(O).
\]

由例 2.2 与命题 2.4(iii)，$ f^{-1}(O) $ 在 X 中是开集。
“(ii)$\Rightarrow$(iii)” 设 $ A \subseteq Y $ 是闭集。则 $ A^c $ 在 $ Y $ 中是开集。由 (ii) 与命题 I.3.8(iv')，$ f^{-1}(A^c) = (f^{-1}(A))^c $ 在 $ X $ 中是开集。因此 $ f^{-1}(A) $ 在 $ X $ 中是闭集。
“(iii)⇒(i)” 设 $ x \in X $。若 $ V $ 是 $ f(x) $ 在 $ Y $ 中的开邻域，则 $ V^c $ 在 $ Y $ 中是闭集。由命题 I.3.8(iv’) 与我们的假设，$ (f^{-1}(V))^c = f^{-1}(V^c) $ 在 $ X $ 中是闭集，即 $ U := f^{-1}(V) $ 在 $ X $ 中是开集。由于 $ x \in U $，$ U $ 是 $ x $ 的满足 $ f(U) \subseteq V $ 的邻域。这意味着 $ f $ 在 $ x $ 处连续。■

\end{subitems}

\end{theorem}

\begin{remark}

根据此定理，一个函数连续当且仅当任何开集的原像是开集，当且仅当任何闭集的原像是闭集。为给出这个重要结果的另一种表述，我们用 $\mathcal{T}_{X}$ 表示度量空间 $X$ 的拓扑，即

\[
\mathcal{T}_{X}:=\left\{O\subseteq X\;;\;O\text{ 是 }X\text{ 中的开集}\right\}.
\]

则

\[
f\colon X\to Y\text{ 是连续的}\Longleftrightarrow f^{-1}\colon\mathcal{T}_{Y}\to\mathcal{T}_{X},
\]

也就是说，$ f: X \to Y $ 连续当且仅当 $ \mathcal{T}_Y $ 在集值函数 $ f^{-1}: \mathcal{P}(Y) \to \mathcal{P}(X) $ 下的像包含于 $ \mathcal{T}_X $。■

下面的例子表明如何用定理 2.20 证明某些集合是开集或闭集。

\end{remark}

\begin{example}

\begin{items}
\item

设 $X$ 与 $Y$ 是度量空间，$f: X \to Y$ 连续。则对每个 $y \in Y$，$f$ 的纤维 $f^{-1}(y)$ 在 $X$ 中是闭集，也就是说，方程 $f(x) = y$ 的解集是闭集。

\proofhead
这由推论 2.18 与定理 2.20 得出。
\item

设 $ k, n \in \mathbb{N}^\times $ 满足 $ k \leq n $。则 $ \mathbb{K}^k $ 在 $ \mathbb{K}^n $ 中是闭集。

\proofhead
若 k = n 则断言是显然的。对 k < n，考虑投影
\[
\operatorname{pr}\colon\mathbb{K}^{n}\to\mathbb{K}^{n-k}~,\quad\left(x_{1},\ldots,x_{n}\right)\mapsto\left(x_{k+1},\ldots,x_{n}\right)\;.
\]

则例 1.3(h) 表明这个函数是连续的。此外 $ \mathbb{K}^k = \mathrm{pr}^{-1}(0) $。因此断言由 (a) 得出。$ \blacksquare $

\item

不等式的解集 设 $ f: X \to \mathbb{R} $ 连续且 $ r \in \mathbb{R} $。则 $ \{x \in X; f(x) \leq r\} $ 在 $ X $ 中是闭集，$ \{x \in X; f(x) < r\} $ 在 $ X $ 中是开集。

\proofhead
显然
\[
\left\{x\in X;f(x)\leq r\right\}=f^{-1}\big((-\infty,r]\big)\quad\text{且}\quad\left\{x\in X;f(x)<r\right\}=f^{-1}\big((-\infty,r)\big).
\]

因此断言由例 2.19(a)、(b) 与定理 2.20 得出。

\item

闭的 n 维单位立方体

\[
I^{n}:=\left\{x\in\mathbb{R}^{n};0\leq x_{k}\leq1,1\leq k\leq n\right\}
\]

在 $ \mathbb{R}^{n} $ 中是闭集。

\proofhead
设 $ \mathrm{pr}_k: \mathbb{R}^n \to \mathbb{R} $，$ (x_1, \ldots, x_n) \mapsto x_k $ 是第 $ k $ 个投影。则
\[
I^{n}=\bigcap_{k=1}^{n}\left(\left\{x\in\mathbb{R}^{n} ;\mathrm{pr}_{k}(x)\leq1\right\}\cap\left\{x\in\mathbb{R}^{n} ;\mathrm{pr}_{k}(x)\geq0\right\}\right).
\]

由 (c)，$ I^{n} $ 是闭集的有限交，从而由命题 2.6，它本身是闭集。

\item

闭集（或开集）的连续像不必是闭集（或开集）。

\proofhead
(i) 设 $X := \mathbb{R}^2$ 且 $A := \{(x, y) \in \mathbb{R}^2 ; xy = 1\}$。由于函数 $\mathbb{R}^2 \to \mathbb{R}$，$(x, y) \mapsto xy$ 是连续的（见命题 1.5(ii)），由 (a) 可得集合 $A$ 在 $X$ 中是闭集。尽管投影 $\mathrm{pr}_1 : \mathbb{R}^2 \to \mathbb{R}$ 是连续的，但 $\mathrm{pr}_1(A) = \mathbb{R}^{\times}$ 在 $\mathbb{R}$ 中不是闭集。
\end{items}

(ii) 对第二个断言，设 $ X := Y := \mathbb{R} $，$ O := (-1,1) $ 且 $ f: \mathbb{R} \to \mathbb{R} $，$ x \mapsto x^2 $。则 $ O $ 在 $ \mathbb{R} $ 中是开集且 $ f $ 连续，但 $ f(O) = [0,1) $ 在 $ \mathbb{R} $ 中不是开集。$ \blacksquare $

\end{example}

\subsection*{连续延拓}

设 $X$ 与 $Y$ 是度量空间。设 $D \subseteq X$，$f: D \to Y$ 连续，且 $a \in X$ 是 $D$ 的极限点。若 $D$ 不是闭集，则 $a$ 可能不在 $D$ 中，从而 $f$ 在 $a$ 处无定义。在本节中我们考虑是否能定义 $f(a)$，使得 $f$ 在 $D \cup \{a\}$ 上连续。若这样的延拓存在，则对 $D$ 中任何收敛于 $a$ 的序列 $(x_n)$，$(f(x_n))$ 收敛于 $f(a)$。因此，对（不必连续的）函数 $f: D \to Y$ 与 $D$ 的极限点 $a$，我们定义

\begin{equation}
\lim_{x\to a}f(x)=y
\end{equation}

若 $y \in Y$ 满足：对 $D$ 中每个收敛于 $a$ 的序列 $(x_n)$，序列 $(f(x_n))$ 在 $Y$ 中收敛于 $y$。

\begin{remark}

\begin{items}
\item

下列陈述等价：

\begin{subitems}
\item

$ \lim_{x \to a} f(x) = y $。

\item

对 $ y $ 在 $ Y $ 中的每个邻域 $ V $，存在 $ a $ 在 $ X $ 中的邻域 $ U $ 使得 $ f(U \cap D) \subseteq V $。

\proofhead
“(i)⇒(ii)” 我们证明逆否命题。设存在 $y$ 在 $Y$ 中的邻域 $V$，使得对 $a$ 在 $X$ 中的每个邻域 $U$ 都有 $f(U \cap D) \not\subseteq V$。特别地，
\[
f\bigl(\mathbb{B}_{X}(a,1/n)\cap D\bigr)\cap V^{c}\neq\emptyset\;,\qquad n\in\mathbb{N}^{\times}\;.
\]

因此，对每个 $n \in \mathbb{N}^\times$，我们可以选取某个 $x_n \in \mathbb{B}_X(a,1/n) \cap D$ 使得 $f(x_n) \in V^c$。特别地，序列 $(x_n)$ 在 $D$ 中且收敛于 $a$。由于对每个 $n$ 有 $f(x_n) \notin V$，$(f(x_n))$ 不可能收敛于 $y$。
“(ii)⇒(i)” 设 $ (x_n) $ 是 $ D $ 中的序列且在 $ X $ 中有 $ x_n \to a $，$ V $ 是 $ y $ 在 $ Y $ 中的邻域。由假设，存在 a 的某个邻域 $ U $ 使得 $ f(U \cap D) \subseteq V $。由于 $ (x_n) $ 收敛于 $ a $，存在某个 $ N \in \mathbb{N} $ 使得对所有 $ n \geq N $ 有 $ x_n \in U $。因此对所有 $ n \geq N $，像 $ (f(x_n)) $ 包含于 $ V $。这意味着 $ f(x_n) \to y $。■

\end{subitems}

\item

若 $ a \in D $ 是 $ D $ 的极限点，则

\[
\lim_{x\to a}f(x)=f(a)\Longleftrightarrow f\text{ 在 }a\text{ 处连续}.
\]

\proofhead
这由 (a) 得出。
\end{items}

\end{remark}

\begin{proposition}

设 $X$ 与 $Y$ 是度量空间，$D \subseteq X$，$f: D \to Y$ 连续。设 $a \in D^c$ 是 $D$ 的极限点，且存在某个 $y \in Y$ 使得 $\lim_{x \to a} f(x) = y$。则

\begin{equation*}
\overline{f}\colon D\cup\{a\}\to Y,\quad x\mapsto\left\{\begin{array}{c l}{f(x),}&{\quad x\in D,}\\ {y,}&{\quad x=a,}\\ \end{array}\right.
\end{equation*}

是 $f$ 到 $D \cup \{a\}$ 上的连续延拓。

\proofhead
我们只需证明 $ \overline{f}\colon D\cup\{a\}\to Y $ 在 $ a $ 处连续。但这直接由注 2.23 得出。$ \blacksquare $
\end{proposition}

对特殊情形 $ X \subseteq \mathbb{R} $，我们可以如下定义单侧极限。设 $ D \subseteq X $，$ f: D \to Y $ 是函数，$ a \in X $ 是 $ D \cap (-\infty, a] $（或 $ D \cap [a, \infty) $）的极限点。则我们定义\footnote{在不致引起混淆时，我们也写作 $ f(a-) := \lim_{x \to a-} f(x) $ 与 $ f(a+) := \lim_{x \to a+} f(x) $。}左（或右）极限

\[
\lim_{x\to a-}f(x)\quad(\text{或}\lim_{x\to a+}f(x))
\]

与 $\lim_{x \to a} f(x)$ 类似，只允许满足 $x_n < a$（或 $x_n > a$）的序列。类似地，若对每个满足 $x_n \to \infty$（或 $x_n \to -\infty$）的序列 $(x_n)$ 都有 $f(x_n) \to y$，我们写作 $y = \lim_{x \to \infty} f(x)$（或 $y = \lim_{x \to -\infty} f(x)$）。

\begin{example}

\begin{items}
\item

设 $ X := \mathbb{R} $，$ D := \mathbb{R} \setminus \{1\} $，$ n \in \mathbb{N}^\times $，且 $ f: D \to \mathbb{R} $ 由 $ f(x) := (x^n - 1)/(x - 1) $ 定义。则

\[
\lim_{x\to1}f(x)=\lim_{x\to1}\frac{x^{n}-1}{x-1}=n.
\]

\proofhead
由习题 I.8.1(b) 我们有
\[
\frac{x^{n}-1}{x-1}=1+x+x^{2}+\cdots+x^{n-1}.
\]

断言由此以及 $ \mathbb{R} $ 中多项式的连续性得出。

\item

对 $X := \mathbb{C}$ 与 $D := \mathbb{C}^\times$，

\[
\lim_{z\to0}\frac{\exp(z)-1}{z}=1.
\]

\proofhead
由 $ \exp(z) = \sum z^k / k! $ 我们得到
\[
\frac{\exp(z)-1}{z}-1=\frac{z}{2}\Big[1+\frac{z}{3}+\frac{z^{2}}{3\cdot4}+\frac{z^{3}}{3\cdot4\cdot5}+\cdots\Big].
\]

因此，对所有满足 $ |z| < 1 $ 的 $ z \in \mathbb{C}^\times $，我们有不等式

\[
\left|\frac{\exp(z)-1}{z}-1\right|\leq\frac{|z|}{2}\bigl[1+|z|+|z^{2}|+|z^{3}|+\cdots\bigr]=\frac{|z|}{2(1-|z|)}.
\]

断言于是由

\[
\lim_{z\to0}\left(\frac{|z|}{2(1-|z|)}\right)=0
\]

得出，它是注 2.23(b) 与 $ |z|/(1-|z|) $ 在 z=0 处连续性的推论。

\item

设 $X := D := Y := \mathbb{R}$，且对 $n \in \mathbb{N}$ 有 $f(x) := x^n$。则

\begin{equation*}
\lim_{x\to\infty}x^{n}=\left\{\begin{array}{cl}1,&n=0,\\\infty,&n\in\mathbb{N}^{\times},\end{array}\right.
\end{equation*}

且

\begin{equation*}
\lim_{x\to-\infty}x^{n}=\left\{\begin{array}{r}1,\\ \infty,\\ -\infty,\end{array}\quad\begin{array}{l}n=0,\\ n\in2\mathbb{N}^{\times},\\ n\in2\mathbb{N}+1.\end{array}\right.
\end{equation*}
\item

因为 $ \lim_{x \to 0^-} 1/x = -\infty $ 且 $ \lim_{x \to 0+} 1/x = \infty $，函数 $ \mathbb{R}^\times \to \mathbb{R} $，$ x \mapsto 1/x $ 不能延拓为 $ \mathbb{R} $ 上的连续函数。

\end{items}

\end{example}

\subsection*{相对拓扑}

设 $X$ 是度量空间，$Y$ 是 $X$ 的子集。则 $Y$ 本身关于由 $X$ 诱导的度量 $d_Y := d|Y \times Y$ 是度量空间，因此“在 $(Y, d_Y)$ 中是开集”与“在 $(Y, d_Y)$ 中是闭集”都是良定义的概念。

有另一种定义 $Y$ 的开子集的方法，它完全避免使用度量。这个定义只要求 $X$ 是拓扑空间。具体地，若存在 $X$ 中的开集 $O$（或 $X$ 中的闭集 $A$）使得 $M = O \cap Y$（或 $M = A \cap Y$），则 $Y$ 的子集 $M$ 在 $Y$ 中是 $\text{开的}$（或 $\text{闭的}$）。若 $M \subseteq Y$ 在 $Y$ 中是开集（或闭集），我们也说 $M$ 在 $Y$ 中是相对开的（或相对闭的）。利用这些定义，容易看出 $X$ 的拓扑结构在 $Y$ 上诱导出拓扑结构。

\bookfigure[0.24]{figures/866df4896fc298cc.jpg}

因此我们有定义 Y 的开子集的两种方式。下一个命题表明这些定义是等价的。

\begin{proposition}

设 $X$ 是度量空间且 $M \subseteq Y \subseteq X$。则 $M$ 在 $Y$ 中是开集（或闭集）当且仅当 $M$ 在 $(Y, d_Y)$ 中是开集（或闭集）。

\proofhead
不失一般性，我们可以设 M 非空。

\begin{subitems}
\item

设 $M$ 在 $Y$ 中是开集。则存在 $X$ 中的某个开集 $O$ 使得 $M = O \cap Y$。于是，对每个 $x \in M$，存在某个 $r > 0$ 使得 $\mathbb{B}_X(x, r) \subseteq O$。由于

\[
\mathbb{B}_{Y}(x,r)=\mathbb{B}_{X}(x,r)\cap Y\subseteq O\cap Y=M,
\]

x 是 M 关于 $ (Y, d_{Y}) $ 的内点。从而 M 在 $ (Y, d_{Y}) $ 中是开集。

\item

现在设 $M$ 在 $(Y, d_Y)$ 中是开集。对每个 $x \in M$，存在某个 $r_x > 0$ 使得 $\mathbb{B}_Y(x, r_x) \subseteq M$。令 $O := \bigcup_{x \in M} \mathbb{B}_X(x, r_x)$。则由例 2.2 与命题 2.4(ii)，$O$ 是 $X$ 的开子集。此外，由命题 I.2.7(ii)，

\[
O\cap Y=\Big(\bigcup_{x\in M}\mathbb{B}_{X}(x,r_{x})\Big)\cap Y=\bigcup_{x\in M}\big(\mathbb{B}_{X}(x,r_{x})\cap Y\big)=\bigcup_{x\in M}\mathbb{B}_{Y}(x,r_{x})=M.
\]

因此 M 在 X 中是开集。

\item

接下来设 $M$ 在 $Y$ 中是闭集，即存在 $X$ 中的闭集 $A$ 使得 $M = Y \cap A$。因为 $Y \setminus M = Y \cap A^c$，由 (i) 可得 $Y \setminus M$ 在 $(Y, d_Y)$ 中是开集。因此 $M$ 在 $(Y, d_Y)$ 中是闭集。

\item

最后，若 $M$ 在 $(Y, d_Y)$ 中是闭集，则 $Y \setminus M$ 在 $(Y, d_Y)$ 中是开集。由 (ii)，$Y \setminus M$ 在 $Y$ 中是开集，于是存在 $X$ 中的开集 $O$ 使得 $O \cap Y = Y \setminus M$。这推出 $M = Y \cap O^c$，从而 $M$ 在 $Y$ 中是闭集。■

\end{subitems}

\end{proposition}

\begin{corollary}

若 $ M \subseteq Y \subseteq X $，则 M 在 Y 中是开集当且仅当 $ Y \setminus M $ 在 Y 中是闭集。

\end{corollary}

\begin{example}

\begin{items}
\item

设 $ X := \mathbb{R}^2 $，$ Y := \mathbb{R} \times \{0\} $ 且 $ M := (0, 1) \times \{0\} $。则 $ M $ 在 $ Y $ 中是开集，但在 $ X $ 中不是。

\item

设 $X := \mathbb{R}$ 且 $Y := (0, 2]$。则 $(1, 2]$ 在 $Y$ 中是开集但在 $X$ 中不是，$(0, 1]$ 在 $Y$ 中是闭集但在 $X$ 中不是。

\end{items}

\end{example}

\subsection*{一般拓扑空间}

尽管度量空间是我们大部分讨论的自然框架，但在后面的章节——以及其他书中——一般拓扑空间也很重要。为此，分析本节的诸定义与命题，找出哪些在任何拓扑空间中都成立，是有益的。我们在下面的注中做这件事。

\begin{remark}

设 $ X = (X, \mathcal{T}) $ 是拓扑空间。

\begin{items}
\item

如上，若 $ A^c $ 是开集，即若 $ A^c \in \mathcal{T} $，则称 $ A \subseteq X $ 是闭集。聚点、极限点与 $ \overline{A} $ 的定义保持不变。则显然命题 2.6 与 2.8 仍然成立。

\item

若存在开集 $O$ 使得 $A \subseteq O \subseteq U$，则称子集 $U \subseteq X$ 为 $X$ 的子集 $A$ 的\textbf{邻域}。若 $A = \{x\}$，则 $U$ 称为 $x$ 的\textbf{邻域}。$x$ 的所有邻域之集我们仍记作 $\mathcal{U}(x)$，或更精确地记作 $\mathcal{U}_X(x)$。显然每个点都有开邻域。若 $x$ 的某个邻域包含于 $A$，则称点 $x$ 为 $A \subseteq X$ 的\textbf{内部}点。显然这些定义与已为度量空间引入的定义是一致的。
最后，$ A \subseteq X $ 的内部 $ \mathring{A} $ 与边界 $ \partial A $ 的定义与度量空间的情形完全相同。于是容易验证命题 2.12 与 2.14，以及推论 2.13 与 2.15 仍然成立。因此我们有 $ \overline{A} = \mathrm{cl}(A) $ 与 $ \mathring{A} = \mathrm{int}(A) $。

\item

命题 2.9 与 2.11，以及推论 2.10 在一般拓扑空间中不成立。尽管如此，下列结论总是成立：若 $A$ 是闭集且 $(x_k)$ 是 $A$ 中的收敛序列且 $\lim x_k = x$，则 $x$ 在 $A$ 中。当然，这里序列的 $\text{收敛}$ 与收敛序列的 $\text{极限}$ 的定义与 II.1 节中完全相同。对命题 2.9 证明的分析表明，用到的是度量空间的下列性质：

\begin{equation}
\left.\begin{gathered} \text{对每个点 } x \in X \text{，存在 } x \text{ 的邻域序列 } (U_k) \text{，使得} \\ \text{对 } x \text{ 的任何邻域 } U \text{，存在某个 } k \in \mathbb{N} \text{ 使得 } U_k \subseteq U\text{。} \end{gathered}\right\}
\end{equation}

对度量空间，取 $ U_k := \mathbb{B}(x, 1/k) $ 即可。
如上所述的邻域序列 $ (U_{k}) $ 称为 x 的可数邻域基。使 (2.2) 成立的拓扑空间称为满足第一可数性公理。

\item

我们已经注意到，命题 2.17 在一般拓扑空间中不成立。满足豪斯多夫条件的拓扑空间称为豪斯多夫空间。命题 2.17 表明任何度量空间都是豪斯多夫空间。
在豪斯多夫空间中，推论 2.18 以完全相同的证明成立，因此每个单元素集都是闭集。此外，豪斯多夫空间中的收敛序列有唯一的极限。

\item

拓扑空间之间函数的连续性的定义与 1 节中完全相同。因此关于复合的连续性的定理 1.8 仍然成立。当 X 是任意拓扑空间时，命题 1.5 与 1.10 成立，但证明必须修改，以便更直接地使用连续性的定义（见习题 19）。
最后，定理 2.20——本节中最重要的定理——对任意拓扑空间成立。因此拓扑空间之间的函数连续当且仅当开集（或闭集）的原像是开集（或闭集）。若 X 是拓扑空间且 Y 是豪斯多夫空间，则例 2.22(a) 与 (c) 仍然成立（为什么？）。

\item

若 X 与 Y 是任意拓扑空间，则定理 1.4 证明的第一部分表明，从 X 到 Y 的任何连续函数也是序列连续的。同一证明的第二部分表明，若 X 满足第一可数性公理，则其逆成立。

\item

设 $X$ 与 $Y$ 是拓扑空间，$a \in X$ 是 $D \subseteq X$ 的极限点。则对 $f: D \to Y$，极限

\begin{equation}
\lim_{x\to a}f(x)
\end{equation}

\begin{subitems}
\item

只有当 X 满足第一可数性公理（更精确地，若 a 有可数邻域基）时才能如在 (2.1) 中那样定义。此时，注 2.23(a) 仍然成立。若 X 是任意拓扑空间，则用注 2.23(a) 的 (ii) 作为 (2.3) 的定义。无论哪种情形，注 2.23(b) 与命题 2.24 都成立。

\end{subitems}

\item

若 Y 是拓扑空间 X 的子集，则相对开（即在 Y 中是开集）与相对闭（即在 Y 中是闭集）的概念如前定义。则

\[
\mathcal{T}_{Y}:=\left\{B\subseteq Y;B\text{ 是 }Y\text{ 中的开集}\right\}
\]

是 $Y$ 上的拓扑，称为 $Y$ 关于 $X$ 的 $\text{相对}$（或 $\text{诱导}$）$\text{拓扑}$。因此 $(Y, \mathcal{T}_Y)$ 本身是拓扑空间，从而是 $X$ 的 $\text{拓扑子空间}$。容易看出，$A \subseteq Y$ 是相对闭的当且仅当 $A$ 在 $(Y, \mathcal{T}_Y)$ 中是闭集，即若 $A^c \in \mathcal{T}_Y$（见推论 2.27）。此外，若 $X$ 是豪斯多夫空间（或满足第一可数性公理），则 $(Y, \mathcal{T}_Y)$ 也是。若 $i := i_Y : Y \to X$，$y \mapsto y$ 是例 I.3.2(b) 的包含映射，则对所有 $A \subseteq X$ 有 $i^{-1}(A) = A \cap Y$。因此，若 $Y$ 有某个其他拓扑 $\mathcal{T}_Y'$，则 $i : (Y, \mathcal{T}_Y') \to X$ 连续当且仅当 $\mathcal{T}_Y'$ 强于相对拓扑 $\mathcal{T}_Y$。

\item

设 $X$ 与 $Y$ 是拓扑空间，$A$ 是 $X$ 的带有相对拓扑的子集。若 $f: X \to Y$ 在 $x_0 \in A$ 处连续，则 $f|A: A \to Y$ 在 $x_0$ 处连续（见例 1.3(k)）。

\proofhead
这由 $ f|A = f\circ i_A $ 与 (h) 得出。■
\end{items}

\end{remark}

\exerciseshead

\begin{exercises}
\item

对度量空间 $X$ 的下列子集 $M$，确定 $\overline{M}$、$\mathring{M}$、$\partial M$ 以及 $M$ 的所有极限点之集 $M'$：

\begin{items}
\item

$ M = (0, 1] $，$ X = \mathbb{R} $。

\item

$ M = (0,1] \times \{0\} $，$ X = \mathbb{R}^2 $。

\item

$ M = \{1/n; n \in \mathbb{N}^\times\}, X = \mathbb{R} $。

\item

$ M = \mathbb{Q} $，$ X = \mathbb{R} $。

\item

$ M = \mathbb{R} \setminus \mathbb{Q} $，$ X = \mathbb{R} $。

\end{items}

\item

设 $ \mathbb{Q} $ 带有自然度量，$ S := \{ x \in \mathbb{Q} ; -\sqrt{2} < x < \sqrt{2} \} $。证明或否定下列陈述：

\begin{items}
\item

S 在 Q 中是开集。

\item

S 在 Q 中是闭集。

\end{items}

\item

设 X 是非空集合，d 是 X 上的离散度量。证明下列结论：

\begin{items}
\item

X 的每个子集都是开集，即 $ \mathcal{P}(X) $ 是 $ (X, d) $ 的拓扑。

\item

一般地，$ \bar{\mathbb{B}}(x,r)=\overline{\mathbb{B}(x,r)} $ 并不成立。

\end{items}

\item

对 $S := \left\{ (x, y) \in \mathbb{R}^2 ; \, x^2 + y^2 < 1 \right\} \setminus \left([0,1) \times \{0\}\right)$，确定 $(\overline{S})^0$。是否有 $(\overline{S})^0 = S$？

\item

设 $X$ 是度量空间且 $A \subseteq X$。证明 $\mathring{A} = X \setminus \overline{(X \setminus A)}$。

\item

设 $ X_j $（$ j = 1, \ldots, n $）是度量空间，$ X := X_1 \times \cdots \times X_n $。证明下列结论：

\begin{items}
\item

若对所有 $j$ 有 $O_j$ 在 $X_j$ 中是开集，则 $O_1 \times \cdots \times O_n$ 在 $X$ 中是开集。

\item

若对所有 $j$ 有 $A_{j}$ 在 $X_{j}$ 中是闭集，则 $A_{1} \times \cdots \times A_{n}$ 在 $X$ 中是闭集。

\end{items}

\item

设 $ h: \mathcal{P}(X) \to \mathcal{P}(X) $ 是具有下列性质的函数：

\begin{subitems}
\item

$ h(\emptyset) = \emptyset $，

\item

$ h(A) \supseteq A $，$ A \in \mathcal{P}(X) $，

\item

$ h(A \cup B) = h(A) \cup h(B) $，$ A, B \in \mathcal{P}(X) $，

\item

$ h \circ h = h $。

\end{subitems}

\begin{items}
\item

令 $ \mathcal{T}_h := \{ A^c \in \mathcal{P}(X) ; \, h(A) = A \} $，并证明 $ (X, \mathcal{T}_h) $ 是拓扑空间。

\item

给定拓扑空间 $ (X, \mathcal{T}) $，求一个满足 (i)-(iv) 且 $ \mathcal{T}_h = \mathcal{T} $ 的函数 $ h: \mathcal{P}(X) \to \mathcal{P}(X) $。

\end{items}

\item

设 $X$ 是度量空间且 $A, B \subseteq X$。证明或否定 $(A \cup B)^\circ = \mathring{A} \cup \mathring{B}$ 与 $(A \cap B)^\circ = \mathring{A} \cap \mathring{B}$。

\item

考虑由 $ \delta(x, y) := |x - y|/(1 + |x - y|) $（见习题 II.1.9）给出的 $ \mathbb{R} $ 上的度量。证明集合 $ A_n := [n, \infty) $（$ n \in \mathbb{N} $）在 $ (\mathbb{R}, \delta) $ 中是闭的且有界的\footnote{另见习题 3.5。}，且对每个 $ k \in \mathbb{N} $ 有 $ \bigcap_{n=0}^{k} A_n \neq \emptyset $，而 $ \bigcap A_n = \emptyset $。

\item

设 $X := \{1, 2, 3, 4, 5\}$ 且

\[
\mathcal{T}:=\left\{\varnothing,X,\{1\},\{3,4\},\{1,3,4\},\{2,3,4,5\}\right\}.
\]
\end{exercises}

\begin{proof}

$ (X, \mathcal{T}) $ 是拓扑空间，并确定 $ \{2, 4, 5\} $ 的闭包。

\end{proof}

\begin{exercises}
\item[11.]

设 $ \mathcal{T}_1 $ 与 $ \mathcal{T}_2 $ 是集合 $ X $ 上的拓扑。证明或否定 $ \mathcal{T}_1 \cup \mathcal{T}_2 $ 与 $ \mathcal{T}_1 \cap \mathcal{T}_2 $ 是 $ X $ 上的拓扑。

\item[12.]

设 X 与 Y 是度量空间。证明

\[
f\colon X\to Y\text{ 是连续的}\Longleftrightarrow f(\overline{A})\subseteq\overline{f(A)},\;A\subseteq X\;.
\]
\item[13.]

设 $ A $ 与 $ B $ 是度量空间 $ X $ 的闭子集。设 $ Y $ 是度量空间，$ g: A \to Y $ 与 $ h: B \to Y $ 是连续函数，满足

\[
g\mid A\cap B=h\mid A\cap B\quad\text{若}\quad A\cap B\neq\varnothing.
\]

证明函数

\begin{equation*}
f\colon A\cup B\to Y~,\quad x\mapsto\left\{\begin{array}{l l}{g(x)~,\quad}&{x\in A~,}\\ {h(x)~,\quad}&{x\in B~,}\\ \end{array}\right.
\end{equation*}

是连续的。

\item[14.]

度量空间 $(X,d)$ 与 $(Y,\delta)$ 之间的函数 $f: X \to Y$ 称为开的，若 $f(\mathcal{T}_d) \subseteq \mathcal{T}_\delta$，即若开集的像是开集。若对任何闭集 $A$ 都有 $f(A)$ 是闭集，则函数 $f$ 称为闭的。设 $d$ 表示 $\mathbb{R}$ 上的自然度量，$\delta$ 表示离散度量。证明下列结论：

\begin{items}
\item

id: $ (\mathbb{R}, d) \to (\mathbb{R}, \delta) $ 是开的且闭的，但不连续。

\item

id: $ (\mathbb{R}, \delta) \to (\mathbb{R}, d) $ 是连续的，但既不开也不闭。

\end{items}

\item[15.]

设 $f: \mathbb{R} \to \mathbb{R}$，$x \mapsto \exp(x)$ zigzag$(x)$（见习题 1.1）。则 $f$ 是连续的，但既不开也不闭。（提示：考虑习题 II.8.10，并确定 $f\big((-\infty,0)\big)$ 与 $f\big(\{\, -(2n+1)/2 \;;\; n\in\mathbb{N}\,\}\big)$。）

\item[16.]

证明函数

\begin{equation*}
f\colon[0,2]\to[0,2],\quad x\mapsto\left\{\begin{array}{c l}{0,}&{\quad x\in[0,1],}\\ {x-1,}&{\quad x\in(1,2],}\\ \end{array}\right.
\end{equation*}

是连续的且闭的，但不开。

\item[17.]

设 $ S^1 := \{(x, y) \in \mathbb{R}^2; x^2 + y^2 = 1\} $ 是 $ \mathbb{R}^2 $ 中的单位圆，带有自然度量。证明函数

\begin{equation*}
f\colon S^{1}\to[0,2),\quad(x,y)\mapsto\left\{\begin{array}{c l}{0,}&{\quad y\geq0,}\\ {1+x,}&{\quad y\leq0,}\\ \end{array}\right.
\end{equation*}

是闭的，但既不连续也不开。

\item[18.]

设 X 与 Y 是度量空间，且

\[
p\colon X\times Y\to X~,\quad(x,y)\mapsto x
\]

是到 X 上的典范投影。则 p 是连续的且开的，但一般不是闭的。

\item[19.]

对任意拓扑空间 X 证明命题 1.5 与 1.10。

\item[20.]

设 X 与 Y 是度量空间，$ f: X \to Y $。证明（见习题 1.17）

\[
A_{n}:=\left\{x\in X\;;\;\omega_{f}(x)\geq1/n\right\}
\]

对每个 $ n \in \mathbb{N}^{\times} $ 都是闭集。

\item[21.]

设 $X$ 是度量空间且 $A \subseteq X$。证明下列结论：

\begin{subitems}
\item

若 A 是完备的，则 A 在 X 中是闭集。其逆一般不正确。

\item

若 X 是完备的，则 A 完备当且仅当 A 在 X 中是闭集。

\end{subitems}

\end{exercises}

\section{紧致性}

我们已经看到，开集的连续像未必是开的，闭集的连续像未必是闭的。在接下来的两节中，我们研究拓扑空间的某些性质，这些性质与上述情形相反，在连续函数下得以保持。这些性质具有深远的重要意义，并且对实值函数的研究特别有用。

\subsection*{覆盖}

下面，$X := (X, d)$ 是一个度量空间。

集合族 $\{A_\alpha \subseteq X; \alpha \in \mathbb{A}\}$ 称为子集 $K \subseteq X$ 的一个\textbf{覆盖}，若 $K \subseteq \bigcup_\alpha A_\alpha$。若每个 $A_\alpha$ 在 $X$ 中都是开集，则称该覆盖为\textbf{开的}。若 $K$ 的每个开覆盖都有一个有限子族也构成 $K$ 的覆盖，则称子集 $K \subseteq X$ 为\textbf{紧致的}。换句话说，若 $K$ 的每个开覆盖都有一个有限\textbf{子覆盖}，则 $K \subseteq X$ 是紧致的。

\begin{example}

\begin{items}
\item

设 $ (x_k) $ 是 $ X $ 中收敛于极限 $ a $ 的序列。则集合 $ K := \{a\} \cup \{x_k\;;\;k \in \mathbb{N}\} $ 是紧致的。

\proofhead
设 $\{O_\alpha; \alpha \in \mathbb{A}\}$ 是 $K$ 的一个开覆盖。则存在 $\alpha$ 与 $\alpha_k \in \mathbb{A}$，使得 $a \in O_\alpha$ 且对所有 $k \in \mathbb{N}$ 有 $x_k \in O_{\alpha_k}$。因为 $\lim x_k = a$，存在某个 $N \in \mathbb{N}$，使得对所有 $k > N$ 有 $x_k \in O_\alpha$。则 $\{O_{\alpha_k}; 0 \leq k \leq N\} \cup \{O_\alpha\}$ 是所给 $K$ 的覆盖的一个有限子覆盖。
\item

一般地，若极限 a 不包含在 K 中，则 (a) 的断言为假。

\proofhead
设 $X := \mathbb{R}$ 且 $A := \{1/k; k \in \mathbb{N}^\times\}$。令 $O_1 := (1/2,2)$，并对所有 $k \geq 2$，令 $O_k := (1/(k+1), 1/(k-1))$。则 $\{O_k; k \in \mathbb{N}^\times\}$ 是 $A$ 的一个开覆盖，且具有如下性质：每个 $O_k$ 恰好包含 $A$ 的一个元素。因此 $\{O_k; k \in \mathbb{N}\}$ 没有 $A$ 的有限子覆盖。
\item

自然数集合 N 在 $ \mathbb{R} $ 中不是紧致的。

\proofhead
同样只需构造 $\mathbb{N}$ 的一个开覆盖 $\{O_k; k \in \mathbb{N}\}$，使得每个 $O_k$ 恰好包含一个自然数，例如，对所有 $k \in \mathbb{N}$，$O_k := (k - 1/3, k + 1/3)$。
\end{items}

\end{example}

\begin{proposition}

任意紧致集 $ K \subseteq X $ 在 X 中是闭的且是有界的。

\proofhead
设 $ K \subseteq X $ 是紧致的。

\begin{subitems}
\item

首先证明 $K$ 在 $X$ 中是闭的。显然只需考虑情形 $K \neq X$，因为 $X$ 在 $X$ 中是闭的。于是假设 $x_0$ 在 $K^c$ 中。由于豪斯多夫性质，对每个 $y \in K$，存在开邻域 $U_y \in \mathcal{U}(y)$ 与 $V_y \in \mathcal{U}(x_0)$，使得 $U_y \cap V_y = \emptyset$。因为 $\{U_y ; \, y \in K\}$ 是 $K$ 的一个开覆盖，存在 $K$ 中有限多个点 $y_0, \ldots, y_m$，使得 $K \subseteq \bigcup_{j=0}^m U_{y_j} =: U$。
由命题 2.4，$U$ 与 $V := \bigcap_{j=0}^{m} V_{y_j}$ 是开的且不交。因此 $V$ 是 $x_0$ 的一个邻域，使得 $V \subseteq K^c$，即 $x_0$ 是 $K^c$ 的内点。由于这对每个 $x_0 \in K^c$ 都成立，$K^c$ 是开集且 $K$ 是闭集。

\item

为验证 $K$ 的有界性，取定 $X$ 中某个 $x_0$。由于由例 2.2，$\mathbb{B}(x_0, k)$ 是开的且 $K \subseteq \bigcup_{k=1}^{\infty} \mathbb{B}(x_0, k) = X$，$K$ 的紧致性蕴含存在 $k_0, \ldots, k_m \in \mathbb{N}$，使得 $K \subseteq \bigcup_{j=0}^{m} \mathbb{B}(x_0, k_j)$。特别地，$K \subseteq \mathbb{B}(x_0, N)$，其中 $N := \max\{k_0, \ldots, k_m\}$。因此 $K$ 是有界的。

\end{subitems}

\end{proposition}

\subsection*{紧致集的刻画}

命题 3.2 的逆命题在一般度量空间中为假（见习题 15），因此紧致集不仅仅是闭集与有界集。取而代之，我们在下一个定理中用聚点给出紧致性的一个刻画。为证明该定理，我们需要下述概念，它在定理 3.10 中再次出现：$ X $ 的子集 $ K $ 是\textbf{完全有界的}，若对每个 $ r > 0 $，存在 $ m \in \mathbb{N} $ 与 $ x_0, \ldots, x_m \in K $，使得 $ K \subseteq \bigcup_{k=0}^{m} \mathbb{B}(x_k, r) $。显然任何完全有界集都是有界的。

\begin{theorem}

子集 $K \subseteq X$ 是紧致的，当且仅当 $K$ 中每个序列都在 $K$ 中有聚点。

证明

\begin{subitems}
\item

首先假设 $K$ 是紧致的，且 $K$ 中存在一个在 $K$ 中没有聚点的序列。于是对每个 $x \in K$，存在 $x$ 的开邻域 $U_x$，它至多包含该序列的有限多项。因为 $\{U_x; x \in K\}$ 是 $K$ 的一个开覆盖，存在 $x_0, \ldots, x_m \in K$，使得 $\{U_{x_k}; k = 0, \ldots, m\}$ 是 $K$ 的一个覆盖。因此 $K$ 至多包含该序列的有限多项。这一矛盾表明 $K$ 中每个序列都在 $K$ 中有聚点。

\item

逆命题的证明分两步进行：

\end{subitems}

\begin{items}
\item

设 K 是 X 的子集，具有如下性质：K 中每个序列都在 K 中有聚点。我们断言 K 是完全有界的。
反设 $K$ 不是完全有界的。则存在某个 $r > 0$，具有如下性质：对任何有限集 $x_0, \ldots, x_m \in K$，$K$ 都不包含在 $\bigcup_{k=0}^{m}\mathbb{B}(x_k, r)$ 中。特别地，存在某个 $x_0 \in K$，使得 $K$ 不包含在 $\mathbb{B}(x_0, r)$ 中。因此存在某个 $x_1 \in K \setminus \mathbb{B}(x_0, r)$。由于 $K$ 不包含在 $\mathbb{B}(x_0, r) \cup \mathbb{B}(x_1, r)$ 中，存在某个 $x_2 \in K \setminus \left[\mathbb{B}(x_0, r) \cup \mathbb{B}(x_1, r)\right]$。迭代这一过程，我们构造 $K$ 中的序列 $(x_k)$，使得对所有 $n$，$x_{n+1}$ 不在 $\bigcup_{k=0}^{n}\mathbb{B}(x_k, r)$ 中。由假设，序列 $(x_k)$ 在 $K$ 中有聚点 $x$，因此特别地，存在 $m, N \in \mathbb{N}^\times$，使得 $d(x_N, x) < r/2$ 且 $d(x_{N+m}, x) < r/2$。三角不等式蕴含 $d(x_N, x_{N+m}) < r$，即 $x_{N+m}$ 在 $\mathbb{B}(x_N, r)$ 中。这与序列 $(x_k)$ 的上述性质矛盾，于是我们证明了 $K$ 是完全有界的。

\item

现在设 $\{O_\alpha : \alpha \in \mathbb{A}\}$ 是 $K$ 的一个开覆盖。与我们的断言相反，假设 $\{O_\alpha : \alpha \in \mathbb{A}\}$ 没有有限子覆盖。由于 $K$ 是完全有界的，对每个 $k \in \mathbb{N}^\times$，存在以 $K$ 中点为圆心、半径为 $1/k$ 的有限个开球，它们构成 $K$ 的一个覆盖。则其中一个开球，记作 $B_k$，具有如下性质：$\{O_\alpha : \alpha \in \mathbb{A}\}$ 的任意有限子集都不是 $K \cap B_k$ 的覆盖。对 $k \in \mathbb{N}^\times$，设 $x_k$ 是 $B_k$ 的圆心。由假设，序列 $(x_k)$ 在 $K$ 中有聚点 $\overline{x}$。
现在设 $ \overline{\alpha} \in \mathbb{A} $ 使得 $ \overline{x} \in O_{\overline{\alpha}} $。由于 $ O_{\overline{\alpha}} $ 是开集，存在某个 $ \varepsilon > 0 $，使得 $ \mathbb{B}(\overline{x}, \varepsilon) \subseteq O_{\overline{\alpha}} $。由于 $ \overline{x} $ 是序列 $ (x_k) $ 的聚点，存在某个 $ M > 2/\varepsilon $，使得 $ d(x_M, \overline{x}) < \varepsilon/2 $。于是对每个 $ x \in B_M $，我们有

\[
d(x,\overline{{x}})\leq d(x,x_{M})+d(x_{M},\overline{{x}})<\frac{1}{M}+\frac{\varepsilon}{2}<\frac{\varepsilon}{2}+\frac{\varepsilon}{2}=\varepsilon,
\]

即 $ B_M \subseteq \mathbb{B}(\overline{x}, \varepsilon) \subseteq O_{\overline{\alpha}} $。这与我们对 $ B_M $ 的选取矛盾，因此覆盖 $ \{O_\alpha : \alpha \in \mathsf{A}\} $ 必有有限子覆盖。■

\end{items}

\end{theorem}

\subsection*{序列紧致性}

若 K 中每个序列都有一个收敛于 K 的某个元素的子序列，则称子集 $ K \subseteq X $ 是序列紧致的。

序列的聚点与收敛子序列之间的关系（见命题 II.1.17）使我们能用序列紧致性重新表述定理 3.3。

\begin{theorem}

度量空间的子集是紧致的，当且仅当它是序列紧致的。

作为定理 3.3 的一个重要应用，我们描述 $ \mathbb{K}^n $ 的紧致子集。

\end{theorem}

\begin{theorem}

（海涅-博雷尔）$ \mathbb{K}^n $ 的子集是紧致的，当且仅当它是闭的且是有界的。特别地，区间是紧致的，当且仅当它是闭的且是有界的。

\proofhead
由命题 3.2，任何紧致集都是闭的且是有界的。逆命题由波尔查诺-魏尔斯特拉斯定理（见定理 II.5.8）、命题 2.11 与定理 3.4 得出。
\end{theorem}

\subsection*{紧致空间上的连续函数}

下述定理表明紧致性在连续函数下得以保持。

\begin{theorem}

设 $X$ 与 $Y$ 是度量空间，$f: X \to Y$ 连续。若 $X$ 是紧致的，则 $f(X)$ 是紧致的。也就是说，紧致集的连续像是紧致的。

\proofhead
设 $\{O_\alpha; \alpha \in \mathsf{A}\}$ 是 $f(X)$ 在 $Y$ 中的一个开覆盖。由定理 2.20，对每个 $\alpha \in \mathsf{A}$，$f^{-1}(O_\alpha)$ 是 $X$ 的开子集。因此 $\{f^{-1}(O_\alpha); \alpha \in \mathsf{A}\}$ 是紧致空间 $X$ 的一个开覆盖，且存在 $\alpha_0, \ldots, \alpha_m \in \mathsf{A}$，使得 $X = \bigcup_{k=0}^m f^{-1}(O_{\alpha_k})$。由此可得 $f(X) \subseteq \bigcup_{k=0}^m O_{\alpha_k}$，即 $\{O_{\alpha_0}, \ldots, O_{\alpha_m}\}$ 是 $\{O_\alpha; \alpha \in \mathsf{A}\}$ 的一个有限子覆盖。因此 $f(X)$ 是紧致的。
\end{theorem}

\begin{corollary}

设 X 与 Y 是度量空间，$ f: X \to Y $ 连续。若 X 是紧致的，则 $ f(X) $ 是有界的。

\proofhead
这直接由定理 3.6 与命题 3.2 得出。
\end{corollary}

\subsection*{极值定理}

对于实值函数，定理 3.6 有一个重要推论：紧致集上的实值连续函数能取得它的最小值与最大值。

\begin{corollary}

（极值定理）设 $X$ 是紧致度量空间，$f: X \to \mathbb{R}$ 连续。则存在 $x_0, x_1 \in X$，使得

\[
f(x_{0})=\min_{x\in X}f(x)\quad\text{且}\quad f(x_{1})=\max_{x\in X}f(x).
\]

\proofhead
由定理 3.6 与命题 3.2 我们知道 $f(X)$ 在 $\mathbb{R}$ 中是闭的且是有界的。因此 $m := \inf\left(f(X)\right)$ 与 $M := \sup\left(f(X)\right)$ 在 $\mathbb{R}$ 中存在。由命题 I.10.5，$f(X)$ 中存在序列 $(y_n)$ 与 $(z_n)$，它们在 $\mathbb{R}$ 中分别收敛于 $m$ 与 $M$。由于 $f(X)$ 是闭集，命题 2.11 蕴含 $m$ 与 $M$ 在 $f(X)$ 中，即存在 $x_0, x_1 \in X$，使得 $f(x_0) = m$ 且 $f(x_1) = M$。

这一结果的重要性可见于下列例子。
\end{corollary}

\begin{example}

\begin{items}
\item

$ \mathbb{K}^{n} $ 上的所有范数都等价。

\begin{subitems}
\item

\proofhead
(i) 设 $ |\cdot| $ 是 $ \mathbb{K}^n $ 上的欧几里得范数，$ \|\cdot\| $ 是 $ \mathbb{K}^n $ 上的任意范数。则只需证明这两个范数的等价性，即存在正常数 $ C $，使得
\begin{equation}
C^{-1}\left|x\right|\leq\left\|x\right\|\leq C\left|x\right|,\qquad x\in\mathbb{K}^{n}.
\end{equation}
\item

令 $ S := \{ x \in \mathbb{K}^n ; \, |x| = 1 \} $。由例 1.3(j) 我们知道函数 $ |\cdot| : \mathbb{K}^n \to \mathbb{R} $ 是连续的，因此由例 2.22(a)，$ S $ 在 $ \mathbb{K}^n $ 中是闭的。当然，$ S $ 在 $ \mathbb{K}^n $ 中也是有界的。由海涅-博雷尔定理，$ S $ 是 $ \mathbb{K}^n $ 的紧致子集。

\item

我们接下来证明 $f: S \to \mathbb{R}$，$x \mapsto \|x\|$ 是连续的。\footnote{这里不能使用例 1.3(j)。为什么？} 设 $\{e_k; 1 \leq k \leq n\}$ 是 $\mathbb{K}^n$ 的标准基。对每个 $x = (x_1, \ldots, x_n) \in \mathbb{K}^n$，我们有 $x = \sum_{k=1}^n x_k e_k$（见例 I.12.4(a) 与注 I.12.5）。由 $\|\cdot\|$ 的三角不等式可得

\begin{equation}
\left\|x\right\|=\left\|\sum_{k=1}^{n}x_{k}e_{k}\right\|\leq\sum_{k=1}^{n}\left|x_{k}\right|\left\|e_{k}\right\|\leq C_{0}\left|x\right|\mathrm{,}\qquad x\in\mathbb{K}^{n},
\end{equation}

其中我们已令 $C_0 := \sum_{k=1}^n \|e_k\|$，并使用了不等式 $|x_k| \leq |x|$。这证明了 (3.1) 的第二个不等式。此外，由 (3.2) 与 $\|\cdot\|$ 的反向三角不等式，可得

\[
|f(x)-f(y)|=\left|\|x\|-\|y\|\right|\leq\|x-y\|\leq C_{0}\left|x-y\right|,\qquad x,y\in S,
\]

这证明了 f 的利普希茨连续性。

\item

对所有 $ x \in S $，我们有 $ f(x) > 0 $，因此由极值定理，我们知道 $ m := \min f(S) $ 是正的，即

\begin{equation}
0<m=\min f(S)\leq f(x)=\|x\|,\qquad x\in S.
\end{equation}

最后设 $ x \in \mathbb{K}^n \setminus \{0\} $。则 $ x/|x| $ 在 $ S $ 中，因此由 (3.3)，我们有 $ m \leq \|x/|x|\| $。即，

\begin{equation}
m\left|x\right|\leq\left\|x\right\|\;,\qquad x\in\mathbb{K}^{n}
\end{equation}

断言现在由 (3.2) 与 (3.4) 以及 $ C := \max\{C_0, 1/m\} $ 得出。■

\end{subitems}

\item

代数学基本定理\footnote{代数学基本定理对实数域 $ \mathbb{R} $ 不成立，例如 $ p = 1 + X^2 $。}任何非常数多项式 $ p \in \mathbb{C}[X] $ 在 $ \mathbb{C} $ 中都有一个零点。

\proofhead
(i) 设 p 是这样一个多项式。不失一般性，我们可以假设 p 的首项系数为 1，从而把 p 写成如下形式
\[
p=X^{n}+a_{n-1}X^{n-1}+\cdots+a_{1}X+a_{0}
\]

其中 $ n \in \mathbb{N}^{\times} $ 且 $ a_k \in \mathbb{C} $。若 $ n = 1 $，结论显然，因此我们假设 $ n \geq 2 $。令

\[
R:=1+\sum_{k=0}^{n-1}\left|a_{k}\right|.
\]

则对每个满足 $ |z| > R \geq 1 $ 的 $ z \in \mathbb{C} $，我们有

\begin{equation*}
\begin{aligned}
\left|p(z)\right|&\geq\left|z\right|^{n}-\left|a_{n-1}\right|\left|z\right|^{n-1}-\cdots-\left|a_{1}\right|\left|z\right|-\left|a_{0}\right|\\&\geq\left|z\right|^{n}-\left(\left|a_{n-1}\right|+\cdots+\left|a_{1}\right|+\left|a_{0}\right|\right)\left|z\right|^{n-1}\\&=\left|z\right|^{n-1}\left(\left|z\right|-(R-1)\right)\geq\left|z\right|^{n-1}>R^{n-1}\geq R.
\end{aligned}
\end{equation*}

因此 $p$ 在球 $\bar{\mathbb{B}}_{\mathbb{C}}(0,R)$ 之外的绝对值大于 $R$。因为 $|p(0)| = |a_0| < R$，这意味着

\[
\operatorname*{i n f}_{z\in\mathbb{C}}|p(z)|=\operatorname*{i n f}_{|z|\leq R}|p(z)|.
\]

\begin{subitems}
\item

我们接下来考虑函数

\[
|p|:\bar{\mathbb{B}}_{\mathbb{C}}(0,R)\to\mathbb{R}~,\quad z\mapsto|p(z)|~,
\]

它是连续函数 $ |\cdot| $ 与 $ p $ 的复合的限制，因此是连续的（见例 1.3(k) 与 1.9(a)，以及推论 1.6）。由海涅-博雷尔定理与例 2.19(d)，闭球 $ \bar{\mathbb{B}}_{\mathbb{C}}(0,R) $ 是紧致的。因此，把极值定理应用于 $ |p| $，存在某个 $ z_0 \in \bar{\mathbb{B}}_{\mathbb{C}}(0,R) $，使得函数 $ |p| $ 在 $ z_0 $ 处取最小值。

\item

假设 $p$ 在 $\bar{\mathbb{B}}_{\mathbb{C}}(0,R)$ 中没有零点。则特别地，$p(z_0) \neq 0$，且 $q := p(X + z_0)/p(z_0)$ 是次数为 $n$ 的多项式，使得

\begin{equation}
|q(z)|\geq1\;,\quad z\in\mathbb{C}\;,\qquad\text{且}\qquad q(0)=1\;.
\end{equation}

因此我们可以把 q 写成如下形式

\[
q=1+\alpha X^{k}+X^{k+1}r
\]

其中 $ \alpha \in \mathbb{C}^{\times} $，$ k \in \{1, \ldots, n-1\} $ 且 $ r \in \mathbb{C}[X] $ 适当。

\item

在这一点上，我们利用复数根的存在性，这一结果我们稍后在第 6 节证明（当然，不使用代数学基本定理）。特别地，该定理说，存在某个 $ z_1 \in \mathbb{C} $\footnote{注意，这一断言对 $ \mathbb{R} $ 不成立。事实上，这是证明中唯一用到 $ \mathbb{C} $ 的特殊性质的地方。}使得 $ z_1^k = -1/\alpha $。因此

\[
q(t z_{1})=1-t^{k}+t^{k+1}z_{1}^{k+1}r(t z_{1}),\qquad t\in[0,1],
\]

从而

\[
|q(t z_{1})|\leq1-t^{k}+t^{k}\cdot t\left|z_{1}^{k+1}r(t z_{1})\right|,\qquad t\in[0,1]~.
\]
\item

最后我们考虑函数

\[
h\colon[0,1]\to\mathbb{R}~,\quad t\mapsto|z_{1}^{k+1}r(t z_{1})|~.
\]

不难看出 $h$ 是连续的（见命题 1.5(ii)、推论 1.6、定理 1.8 与例 1.9(a)）。由海涅-博雷尔定理与推论 3.7，存在某个 $M \geq 1$，使得

\begin{equation}
h(t)=|z_{1}^{k+1}r(t z_{1})|\leq M~,\qquad t\in[0,1]~.
\end{equation}

如果我们在 (3.6) 中使用这个界，我们得到

\[
|q(t z_{1})|\leq1-t^{k}(1-t M)\leq1-t^{k}/2<1~,\qquad t\in\left(0,1/(2M)\right)~,
\]

这与 (3.5) 的第一个断言矛盾。因此 p 在 $ \bar{\mathbb{B}}_{\mathbb{C}}(0, R) $ 中必有一个零点。 $ \blacksquare $

\end{subitems}

推论 设

\[
p=a_{n}X^{n}+a_{n-1}X^{n-1}+\cdots+a_{1}X+a_{0}
\]

其中 $ a_0, \ldots, a_n \in \mathbb{C} $，$ a_n \neq 0 $ 且 $ n \geq 1 $。则存在 $ z_1, \ldots, z_n \in \mathbb{C} $，使得

\[
p=a_{n}\prod_{k=1}^{n}(X-z_{k}).
\]

因此 $\mathbb{C}$ 上每个多项式 $p$ 恰好有 $\deg(p)$ 个（计重数）零点。

\proofhead
由代数学基本定理，对某个 $ z_1 \in \mathbb{C} $ 有 $ p(z_1) = 0 $。由定理 I.8.17，存在某个 $ p_1 \in \mathbb{C}[X] $，使得 $ p = (X - z_1)p_1 $ 且 $ \deg(p_1) = \deg(p) - 1 $。一个简单的归纳论证完成证明。
\item

设 A 与 K 是度量空间的不交子集，其中 K 是紧致的，A 是闭的。则从 K 到 A 的距离 $ d(K, A) $ 是正的，即

\[
d(K,A):=\inf_{k\in K}d(k,A)>0.
\]

\proofhead
由例 1.3(k) 与 (l)，实值函数 $ d(\cdot, A) $ 在 K 上连续，因此由极值定理，存在某个 $ k_0 \in K $，使得 $ d(k_0, A) = d(K, A) $。假设
\[
d(k_{0},A)=\inf_{a\in A}d(k_{0},a)=0.
\]

则存在 $A$ 中的序列 $(a_k)$，使得对所有 $k \to \infty$ 有 $d(k_0, a_k) \to 0$。因此序列 $(a_k)$ 收敛于 $k_0$。因为 $A$ 是闭集，$k_0$ 在 $A$ 中，这与 $A \cap K = \emptyset$ 矛盾。因此我们有 $d(k_0, A) = d(K, A) > 0$。

\item

K 的紧致性在 (c) 中是必要的。

\proofhead
集合 $ A := \mathbb{R} \times \{0\} $ 与 $ B := \{(x, y) \in \mathbb{R}^2; xy = 1\} $ 在 $ \mathbb{R}^2 $ 中是闭的但不是紧致的。由于对 $ n \in \mathbb{N}^\times $ 有 $ d((n,0), (n,1/n)) = 1/n $，我们有 $ d(A, B) = 0 $。 $ \blacksquare $
\end{items}

\end{example}

\subsection*{完全有界性}

鉴于紧致性概念的实际重要性已由上述例子充分展示，我们现在给出紧致集的另一个刻画，它使用完备性与完全有界性。

\begin{theorem}

度量空间的子集是紧致的，当且仅当它是完备的且是完全有界的。

\proofhead
‘$\Rightarrow$’ 设 $ K \subseteq X $ 是紧致的，$ (x_j) $ 是 $ K $ 中的柯西序列。由于 $ K $ 是序列紧致的，$ (x_j) $ 有一个在 $ K $ 中收敛的子序列。因此由命题 II.6.4，序列 $ (x_j) $ 本身在 $ K $ 中收敛。这蕴含 $ K $ 是完备的。

对每个 $ r > 0 $，集合 $ \{\mathbb{B}(x, r) : x \in K\} $ 是 $ K $ 的一个开覆盖。由于 $ K $ 是紧致的，这个覆盖有有限子覆盖。因此我们证明了 $ K $ 是完全有界的。

‘⇐’ 设 $K$ 是完备的且是完全有界的。设 $(x_j)$ 是 $K$ 中的序列。由于 $K$ 是完全有界的，对每个 $n \in \mathbb{N}^\times$，存在以 $K$ 中点为圆心、半径为 $1/n$ 的有限个开球，它们构成 $K$ 的一个覆盖。特别地，$(x_j)$ 有一个子序列 $(x_{1,j})_{j \in \mathbb{N}}$，它包含在一个半径为 1 的球中。则 $(x_{1,j})_{j \in \mathbb{N}}$ 有一个子序列 $(x_{2,j})_{j \in \mathbb{N}}$，它包含在一个半径为 $1/2$ 的球中。进一步，$(x_{2,j})_{j \in \mathbb{N}}$ 有一个子序列 $(x_{3,j})_{j \in \mathbb{N}}$，它包含在一个半径为 $1/3$ 的球中。

迭代这一构造，对每个 $ n \in \mathbb{N}^\times $，得到 $ (x_{n,j})_{j \in \mathbb{N}} $ 的子序列 $ (x_{n+1,j})_{j \in \mathbb{N}} $，它包含在一个半径为 $ 1/(n+1) $ 的球中。

现在对所有 $ n \in \mathbb{N}^\times $ 令 $ y_n := x_{n,n} $。容易验证 $ (y_n) $ 是 $ K $ 中的柯西序列（见注 3.11(a)）。由于 $ K $ 是完备的，序列 $ (y_n) $ 在 $ K $ 中收敛。

因此序列 $ (x_{j}) $ 有一个子序列，即 $ (y_{n}) $，它在 K 中收敛。这表明 K 是序列紧致的，并且由定理 3.4，K 是紧致的。 $ \blacksquare $
\end{theorem}

\begin{remark}

\begin{items}
\item

在前述证明的第二部分，我们使用了一个在许多其他情形中都有用的技巧：从给定的序列 $ (x_{0,j})_{j\in\mathbb{N}} $ 出发，逐次选取子序列 $ (x_{n+1,j})_{j\in\mathbb{N}} $，使得对所有 $ n\in\mathbb{N} $，$ (x_{n+1,j})_{j\in\mathbb{N}} $ 是 $ (x_{n,j}) $ 的子序列。然后构造对角序列：对每个 $ n\in\mathbb{N} $，选取第 $ n $ 个子序列的第 $ n $ 个元素。

\begin{equation*}
\begin{array}{r l r lll}&{\scriptstyle\bigcirc\scriptstyle{x_{0,0}},}&{\scriptstyle x_{0,1},}&{\scriptstyle x_{0,2},}&{\scriptstyle x_{0,3},}&{\scriptstyle \cdots}\\ &{\scriptstyle x_{1,0},}&{\scriptstyle\bigcirc\scriptstyle{x_{1,1}},}&{\scriptstyle x_{1,2},}&{\scriptstyle x_{1,3},}&{\scriptstyle \cdots}\\ &{\scriptstyle x_{2,0},}&{\scriptstyle x_{2,1},}&{\scriptstyle\bigcirc\scriptstyle{x_{2,2}},}&{\scriptstyle x_{2,3},}&{\scriptstyle \cdots}\\ &{\scriptstyle x_{3,0},}&{\scriptstyle x_{3,1},}&{\scriptstyle x_{3,2},}&{\scriptstyle\bigcirc\scriptstyle{x_{3,3}},}&{\scriptstyle \cdots}\\ &{\vdots}&{\vdots}&{\vdots}&{\vdots}&{\vdots\vdots}\\ &{\qquad}&{\qquad}&{\vdots}&{\vdots}&{\vdots\vdots}\end{array}
\end{equation*}

对角序列 $ (y_n) := (x_{n,n})_{n \in \mathbb{N}} $ 显然具有如下性质：对每个 $ N \in \mathbb{N} $，$ (y_n)_{n \geq N} $ 是 $ (x_{N,j})_{j \in \mathbb{N}} $ 的子序列，因此它与每个子序列 $ (x_{n,j})_{j \in \mathbb{N}} $ 具有同样的‘在无穷远处’的性质。

\item

度量空间 X 的子集 K 是紧致的，当且仅当 K 连同诱导度量是一个紧致度量空间。

\proofhead
这是相对拓扑的定义与命题 2.26 的直接推论。
由于注 3.11(b)，本可以只对 X 而不是对 X 的子集 K 表述定理 3.3 与 3.4。然而，在应用中通常给定一个‘底层的’度量空间 X，例如 X 常常是巴拿赫空间，此时要研究的正是 X 的某些子集。因此上述稍长的表述‘更接近实际’。
\end{items}

\end{remark}

\subsection*{一致连续性}

设 $X$ 与 $Y$ 是度量空间，$f: X \to Y$ 连续。则由命题 1.1，对每个 $x_0 \in X$ 与每个 $\varepsilon > 0$，存在某个 $\delta(x_0, \varepsilon) > 0$，使得对每个满足 $d(x, x_0) < \delta$ 的 $x \in X$，我们有 $d(f(x_0), f(x)) < \varepsilon$。正如我们在命题 1.1 之后所注明的，并在例 1.3(a) 中明确看到的，数 $\delta(x_0, \varepsilon)$ 一般依赖于 $x_0 \in X$。另一方面，例 1.3(e) 表明存在这样的连续函数，对其而言数 $\delta$ 可以不依赖于 $x_0 \in X$ 而选取。这样的函数称为一致连续的，并具有重大的实际重要性。具体地，函数 $f: X \to Y$ 称为一致连续的，若对每个 $\varepsilon > 0$，存在某个 $\delta(\varepsilon) > 0$，使得

\[
d\big(f(x),f(y)\big)<\varepsilon\quad\text{对所有 }x,y\in X\text{ 使得 }d(x,y)<\delta(\varepsilon).
\]

\begin{example}

\begin{items}
\item

利普希茨连续函数是一致连续的（见例 1.3(e)）。

\item

函数 $ r: (0, \infty) \to \mathbb{R} $，$ x \mapsto 1/x $ 是连续的，但不是一致连续的。

\proofhead
由于 $r$ 是有理函数的限制，它当然是连续的。现在设 $\varepsilon > 0$。假设存在某个 $\delta := \delta(\varepsilon) > 0$，使得对所有满足 $|x - y| < \delta$ 的 $x, y \in (0,1)$ 有 $|r(x) - r(y)| < \varepsilon$。选取 $x := \delta/(1 + \delta\varepsilon)$ 与 $y := x/2$。则 $x, y \in (0,1)$ 且 $|x - y| = \delta/\left[2(1 + \delta\varepsilon)\right] < \delta$ 且 $|r(x) - r(y)| = (1 + \delta\varepsilon)/\delta > \varepsilon$。这与我们对 $\delta$ 的选取矛盾。
下述重要定理表明，在许多情形中，连续函数自动是一致连续的。
\end{items}

\end{example}

\begin{theorem}

假设 X 与 Y 是度量空间，且 X 是紧致的。若 $f: X \to Y$ 是连续的，则 $f$ 是一致连续的。也就是说，紧致集上的连续函数是一致连续的。

\proofhead
假设 $f$ 是连续的但不是一致连续的。则存在某个 $\varepsilon > 0$，具有如下性质：对每个 $\delta > 0$，存在 $x, y \in X$，使得 $d(x, y) < \delta$ 但 $d(f(x), f(y)) \geq \varepsilon$。特别地，$X$ 中存在序列 $(x_n)$ 与 $(y_n)$，使得
\[
d(x_{n},y_{n})<1/n\quad\text{且}\quad d\bigl(f(x_{n}),f(y_{n})\bigr)\geq\varepsilon~,\qquad n\in\mathbb{N}^{\times}.
\]

由于 $X$ 是紧致的，由定理 3.4，$(x_n)$ 存在子序列 $(x_{n_k})_{k\in\mathbb{N}}$，使得 $\lim_{k\to\infty}x_{n_k}=\overline{x}\in X$。对 $(y_n)$ 的相应子序列 $(y_{n_k})_{k\in\mathbb{N}}$，我们有

\[
d(\overline{{x}},y_{n_{k}})\leq d(\overline{{x}},x_{n_{k}})+d(x_{n_{k}},y_{n_{k}})\leq d(\overline{{x}},x_{n_{k}})+1/n_{k}~,\qquad k\in\mathbb{N}^{\times}~.
\]

因此 $(y_{n_k})_{k \in \mathbb{N}}$ 也收敛于 $\overline{x}$。由于 $f$ 是连续的，这两个序列的像收敛于 $f(\overline{x})$，特别地，存在某个 $K \in \mathbb{N}$，使得

\[
d\big(f(x_{n_{K}}),f(\overline{{x}})\big)<\varepsilon/2\quad\text{且}\quad d\big(f(y_{n_{K}}),f(\overline{{x}})\big)<\varepsilon/2.
\]

这导致矛盾

\[
\varepsilon\leq d\big(f(x_{n_{K}}),f(y_{n_{K}})\big)\leq d\big(f(x_{n_{K}}),f(\overline{{x}})\big)+d\big(f(\overline{{x}}),f(y_{n_{K}})\big)<\varepsilon.
\]

因此 f 是一致连续的。

\end{theorem}

\subsection*{一般拓扑空间中的紧致性}

正如上一节末尾那样，我们想简要考虑一般拓扑空间的情形。诚然，一般情形不再简单，我们这里的讨论必须限于对结果的描述。关于（集合论的）拓扑的证明与更深入的探讨，强烈推荐 Dugundji 的著作 [Dug66]。

\begin{remark}

\begin{items}
\item

设 $ X = (X, \mathcal{T}) $ 是一个拓扑空间。若 $ X $ 是豪斯多夫空间且 $ X $ 的每个开覆盖都有有限子覆盖，则 $ X $ 是紧致的。若 $ X $ 是豪斯多夫空间且每个序列都有收敛子序列，则空间 $ X $ 是序列紧致的。若拓扑子空间 $ (Y, \mathcal{T}_Y) $ 是紧致的（或序列紧致的），则子集 $ Y \subseteq X $ 是紧致的（或序列紧致的）。由命题 2.17 与 2.26 以及注 3.11(b)，这些定义推广了度量空间中紧致子集与序列紧致子集的概念。

\item

豪斯多夫空间 $ X $ 的任何紧致子集 $ K $ 都是闭的。对每个 $ x_0 \in K^c $，$ X $ 中存在不交的开集 $ U $ 与 $ V $，使得 $ K \subseteq U $ 且 $ x_0 \in V $。换句话说，豪斯多夫空间的紧致子集与不在该子集中的点可以被开邻域分离。

\proofhead
这由命题 3.2 证明的第一部分得出 ■
\item

紧致空间的任何闭子集都是紧致的。

\proofhead
见习题 2。
\item

设 X 是紧致的，Y 是豪斯多夫的。则任何连续函数 $ f: X \to Y $ 的像都是紧致的。

\proofhead
定理 3.6 的证明与相对拓扑的定义表明 $ f(X) $ 的每个开覆盖都有有限子覆盖。由于豪斯多夫空间的子空间本身是豪斯多夫空间，断言成立。
\item

在一般拓扑空间中，紧致性与序列紧致性是不同的概念。也就是说，紧致空间未必是序列紧致的，序列紧致空间也未必是紧致的。

\item

在一般拓扑空间中，一致连续性没有定义，因为上面给出的定义本质地使用了度量。

\end{items}

\end{remark}

\exerciseshead

\begin{exercises}
\item

设 $ X_j $，$ j = 1, \ldots, n $，是度量空间。证明 $ X_1 \times \cdots \times X_n $ 是紧致的，当且仅当每个 $ X_j $ 都是紧致的。

\item

设 X 是紧致度量空间，Y 是 X 的子集。证明 Y 是紧致的，当且仅当 Y 是闭的。

\item

设 X 与 Y 是度量空间。若 f 与 $ f^{-1} $ 都连续，则称双射 $ f: X \to Y $ 为同胚。证明下列命题：

\begin{items}
\item

若 $f: X \to Y$ 是同胚，则对所有 $x \in X$ 有 $\mathcal{U}(f(x)) = f(\mathcal{U}(x))$，即‘$f$ 把邻域映为邻域’。

\item

假设 $X$ 是紧致的且 $f: X \to Y$ 是连续的。

\end{items}

\begin{subitems}
\item

f 是闭映射（见习题 2.14）。

\item

若 f 是双射，则它是同胚。

\end{subitems}

\item

若非空集合的子集族 M 的每个有限子族都有非空交集，则称 M 具有有限交性质。

证明下列命题等价：

\begin{items}
\item

X 是紧致度量空间。

\item

$ X $ 的每个具有有限交性质的闭子集族 $ \mathcal{A} $ 都有非空交集，即 $ \bigcap \mathcal{A} \neq \emptyset $。

\end{items}

\item

设 $ (A_j) $ 是 $ X $ 的非空闭子集的序列，且对所有 $ j \in \mathbb{N} $ 有 $ A_j \supseteq A_{j+1} $。证明：若 $ A_0 $ 是紧致的，则 $ \bigcap A_j \neq \emptyset $。\footnote{参见习题 2.9。}

\item

设 E 与 F 是有限维赋范向量空间，$ A: E \to F $ 是线性的。证明 A 是利普希茨连续的。（提示：例 3.9(a)。）

\item

证明所有实正交矩阵的集合 $ O(n) $ 是 $ \mathbb{R}^{(n^2)} $ 的紧致子集。

\item

设

\[
C_{0}:=[0,1]~,\quad C_{1}:=C_{0}\backslash(1/3,2/3)~,\quad C_{2}:=C_{1}\setminus\bigl((1/9,2/9)\cup(7/9,8/9)\bigr)~,\quad\ldots.
\]

一般地，$ C_{n+1} $ 由去掉组成 $ C_n $ 的 $ 2^n $ 个区间中每一个的中间的开三分之一而形成。交集 $ C := \bigcap C_n $ 称为\textbf{康托尔集}。证明下列命题：

\begin{items}
\item

C 是紧致的且具有空的内部。

\item

$C$ 由 $[0,1]$ 中所有这样的数组成，其三进制展开为 $\sum_{k=1}^{\infty} a_k 3^{-k}$，其中 $a_k \in \{0,2\}$。

\item

C 的每个点都是 C 的极限点，即 C 是完全集。

\item

对具有三进制展开 $ \sum_{k=1}^{\infty} a_k 3^{-k} $ 的 $ x \in C $，定义

\[
\varphi(x):=\sum_{k=1}^{\infty}a_{k}2^{-(k+1)}.
\]

则 $ \varphi: C \to [0,1] $ 是递增的、满射的且连续的。

\item

C 是不可数的。

\item

$ \varphi $ 有一个连续延拓 $ f:[0,1] \to [0,1] $，它在 $ [0,1]\backslash C $ 中的每个区间上都是常数。函数 $ f $ 称为 $ C $ 的\textbf{康托尔函数}。

\end{items}

\item

设 $X$ 是度量空间。若对 $X$ 中每个满足 $\lim x_n = a$ 的序列 $(x_n)$，都有 $f(a) \leq \varliminf f(x_n)$，则称函数 $f: X \to \mathbb{R}$ 在 $a \in X$ 处下半连续。若 $-f$ 在 $a$ 处下半连续，则称 $f$ 在 $a$ 处上半连续。最后，若 $f$ 在 $X$ 的每个点处都下半连续（或上半连续），则称 $f$ 是下半连续的（或上半连续的）。

\begin{items}
\item

证明下列命题的等价性：

\begin{subitems}
\item

f 是下半连续的。

\item

对每个 $ a \in X $ 与 $ \varepsilon > 0 $，存在某个 $ U \in \mathcal{U}(a) $，使得对所有 $ x \in U $ 有 $ f(x) > f(a) - \varepsilon $。

\item

对每个 $ \alpha \in \mathbb{R} $，$ f^{-1}\left((\alpha, \infty)\right) $ 是开集。

\item

对每个 $ \alpha \in \mathbb{R} $，$ f^{-1}\big((-\infty, \alpha]\big) $ 是闭集。

\end{subitems}

\item

f 是连续的，当且仅当 f 是下半连续且上半连续的。

\item

设 $ \chi_A $ 是 $ A \subseteq X $ 的特征函数。则 $ A $ 是开集，当且仅当 $ \chi_A $ 是下半连续的。

\item

设 $X$ 是紧致的，$f: X \to \mathbb{R}$ 是下半连续的。则 $f$ 取得它的最小值，即存在某个 $x \in X$，使得对所有 $y \in X$ 有 $f(x) \leq f(y)$。（提示：考虑 $X$ 中满足 $f(x_n) \to \inf f(X)$ 的序列 $(x_n)$。）

\end{items}

\item

设 $ f, g\colon [0,1] \to \mathbb{R} $ 由下式定义

\begin{equation*}
f(x):=\left\{\begin{array}{c l}{1/n\;,}&{\quad x\in\mathbb{Q}\text{，其中 }x=m/n\text{ 为既约分数}\;,}\\ {0\;,}&{\quad x\notin\mathbb{Q}\;,}\end{array}\right.
\end{equation*}

且

\begin{equation*}
g(x):=\left\{\begin{array}{c l}{(-1)^{n}n/(n+1)\;,}&{\quad x\in\mathbb{Q}\text{，其中 }x=m/n\text{ 为既约分数}\;,}\\ {0\;,}&{\quad x\notin\mathbb{Q}\;.}\\ \end{array}\right.
\end{equation*}

证明或否定下列命题：

\begin{items}
\item

f 是上半连续的。

\item

f 是下半连续的。

\item

g 是上半连续的。

\item

g 是下半连续的。

\end{items}

\item

设 $X$ 是度量空间，$f:[0,1)\to X$ 连续。证明：若 $\lim_{t\to1}f(t)$ 存在，则 $f$ 是一致连续的。

\item

下列函数中哪些

\[
f\colon(0,\infty)\to\mathbb{R}~,\quad t\mapsto(1+t^{2})^{-1}~,\qquad g\colon(0,\infty)\to\mathbb{R}~,\quad t\mapsto t^{-2}
\]

是一致连续的？

\item

证明赋范向量空间的有限维子空间是闭的。

（提示：设 $E$ 是赋范向量空间，$F$ 是 $E$ 的有限维子空间。设 $(v_n)$ 是 $F$ 中的序列，$v \in E$ 使得在 $E$ 中 $\lim v_n = v$。由于注 I.12.5、命题 1.10 与波尔查诺-魏尔斯特拉斯定理，存在 $(v_n)$ 的某个子序列 $(v_{n_k})_{k \in \mathbb{N}}$ 与 $w \in F$，使得在 $F$ 中 $\lim_k v_{n_k} = w$。现在利用命题 2.11 与 2.17 证明 $v = w \in F$。）

\item

假设 $X$ 是度量空间，$f: X \to \mathbb{R}$ 是有界的。证明 $\omega_f: X \to \mathbb{R}$ 是上半连续的（见习题 1.17 与 2.20）。

\item

证明 $ \ell_{\infty} $ 中的闭单位球（见注 II.3.6(a)）不是紧致的。

（提示：考虑‘单位向量’ $ e_n $ 构成的序列 $ (e_n) $，其中对所有 $ j \in \mathbb{N} $ 有 $ e_n(j) := \delta_{nj} $。）

\end{exercises}

\section{连通性}

直观上很明显，$ \mathbb{R} $ 中的一个开区间是“连通的”，但若去掉一个点，它就变成“不连通的”。在本节中，我们将使这个直观的连通性概念更加精确。在这样做的过程中，我们再次发现拓扑起着本质的作用。

\subsection*{定义与基本性质}

度量空间 X 称为连通的，若 X 不能表示成两个不交的非空开子集的并。因此 X 是连通的当且仅当

\[
\nexists O_{1},O_{2}\subseteq X,\text{ 开、非空、满足 }O_{1}\cap O_{2}=\varnothing\text{ 且 }O_{1}\cup O_{2}=X.
\]

X 的子集 M 称为在 X 中连通的，若 M 关于由 X 诱导的度量是连通的。

\begin{example}

\begin{items}
\item

显然，空集和任何单元素集都是连通的。

\item

自然数集 N 不是连通的。

\proofhead
由例 2.19(a) 与定理 2.26，子集 $ O_1 := \{0\} = \mathbb{N} \cap (-\infty, 1/2) $ 和 $ O_2 := \{1, 2, 3, \ldots\} = \mathbb{N} \cap (1/2, \infty) $ 在 $ \mathbb{N} $ 中是开集。当然，由于 $ O_1 \cap O_2 = \emptyset $ 且 $ O_1 \cup O_2 = \mathbb{N} $，这表明 $ \mathbb{N} $ 不是连通的。
\item

有理数集 Q 在 $ \mathbb{R} $ 中不是连通的。

\proofhead
子集 $O_1 := \left\{ x \in \mathbb{Q} ; x < \sqrt{2} \right\}$ 与 $O_2 := \left\{ x \in \mathbb{Q} ; x > \sqrt{2} \right\}$ 是开的、非空的，且满足 $O_1 \cap O_2 = \emptyset$ 和 $O_1 \cup O_2 = \mathbb{Q}$。■
\end{items}

\end{example}

\begin{proposition}

对任何度量空间 X，下列各条等价：

\begin{subitems}
\item

X 是连通的。

\item

X 是 X 中唯一既开又闭的非空子集。

\proofhead
“(i)⇒(ii)” 令 $O$ 是 $X$ 中一个既开又闭的非空子集。则 $O^{c}$ 在 $X$ 中也是既开又闭的，并且当然有 $O \cap O^{c} = \emptyset$ 与 $X = O \cup O^{c}$。由于 $X$ 是连通的且由假设 $O$ 非空，由此可得 $O^{c}$ 必为空。因此 $O = X$。
“(ii)⇒(i)” 设 $O_1$ 与 $O_2$ 是 $X$ 的非空开子集，满足 $O_1 \cap O_2 = \emptyset$ 且 $O_1 \cup O_2 = X$。则 $O_1 = O_2^c$ 在 $X$ 中是非空的、既开又闭的，因此由假设，$O_1 = O_2^c = X$。这蕴含 $O_2 = \emptyset$，矛盾。
\end{subitems}

\end{proposition}

\begin{remark}

这个命题常按如下方式用于证明关于连通集的陈述：假设我们要证明连通集中的每个元素 x 集合 $X$ 具有性质 $E$，即 $E(x)$ 对所有 $x \in X$ 成立。令

\[
O:=\left\{x\in X;E(x)\text{ 为真 }\right\}.
\]

那么只需证明集合 O 非空、既开又闭，因为这样一来，由命题 4.2，O = X。■

\end{remark}

\subsection*{$ \mathbb{R} $ 中的连通性}

下一个命题刻画了 $ \mathbb{R} $ 的所有连通子集，并且给出了我们第一个非平凡连通集的具体例子。

\qyxffntextdef{1}{若 $a$ 与 $b$ 是实数，则 $(a, b) \cup \{a, b\} = [a, b]$。}
\begin{theorem}

$ \mathbb{R} $ 的一个子集是连通的当且仅当它是一个区间。

\proofhead
由例 4.1(a)，我们可以假设该子集含有多于一个元素。

“$\Rightarrow$” 令 $X \subseteq \mathbb{R}$ 是连通的。

\begin{subitems}
\item

令 $ a := \inf(X) \in \mathbb{R} $ 和 $ b := \sup(X) \in \mathbb{R} $。由于 $ X $ 至少有两个元素，区间 $ (a, b) $ 非空，并且\qyxffnmark{1} $ X \subseteq (a, b) \cup \{a, b\} $。

\item

我们首先证明包含关系 $(a,b) \subseteq X$。反设 $(a,b)$ 不包含于 $X$。则存在某个 $c \in (a,b)$ 不在 $X$ 中。令 $O_1 := X \cap (-\infty, c)$ 和 $O_2 := X \cap (c, \infty)$。则由命题 2.26，$O_1$ 与 $O_2$ 在 $X$ 中是开集。当然，$O_1$ 与 $O_2$ 不交，且它们的并是 $X$。根据我们对 $a$、$b$、$c$ 的选取，存在元素 $x, y \in X$ 使得 $x < c$ 且 $y > c$。这意味着 $x$ 属于 $O_1$ 且 $y$ 属于 $O_2$，从而 $O_1$ 与 $O_2$ 非空。因此 $X$ 不是连通的，这与我们的假设矛盾。

\item

由于我们已经证明了包含关系 $(a,b)\subseteq X\subseteq(a,b)\cup\{a,b\}$，$X$ 是一个区间。
“⇐” (i) 反设 $X$ 是一个区间，并且存在 $X$ 的开子集、非空子集 $O_1$ 与 $O_2$，使得 $O_1 \cap O_2 = \emptyset$ 且 $O_1 \cup O_2 = X$。选取 $x \in O_1$ 和 $y \in O_2$，并首先考虑 $x < y$ 的情形。由于 $\mathbb{R}$ 是序完备的，$z := \sup(O_1 \cap [x, y])$ 是一个良定义的实数。

\end{subitems}

(ii) 元素 $z$ 不可能在 $O_1$ 中，因为 $O_1$ 在 $X$ 中是开集且 $X$ 是区间，从而存在某个 $\varepsilon > 0$ 使得 $[z, z + \varepsilon) \subseteq O_1 \cap [x, y]$。这与 $z$ 的上确界性质矛盾。类似地，$z$ 不可能在 $O_2$ 中，因为否则存在某个 $\varepsilon > 0$ 使得

\[
(z-\varepsilon,z]\subseteq O_{2}\cap[x,y]~,
\]

这与 $ O_1 \cap O_2 = \emptyset $ 以及 $ z $ 的定义矛盾。于是 $ z \notin O_1 \cup O_2 = X $。另一方面，由于 $ X $ 是区间，$ [x, y] $ 包含于 $ X $。这导致矛盾：$ z \in [x, y] \subseteq X $ 且 $ z \notin X $。$ y < x $ 的情形可类似地证明。

\end{theorem}

\subsection*{广义介值定理}

连通集具有这样的性质：它们在连续函数下的像也是连通的。这个重要的事实可以利用第 2 节的结果轻松证明。

\begin{theorem}

设 $X$ 与 $Y$ 是度量空间，$f: X \to Y$ 连续。若 $X$ 是连通的，则 $f(X)$ 也是连通的。也就是说，连通集的连续像是连通的。

\proofhead
反设 $f(X)$ 不是连通的。则存在 $f(X)$ 的非空子集 $V_1$ 与 $V_2$，使得 $V_1$ 与 $V_2$ 在 $f(X)$ 中是开集，$V_1 \cap V_2 = \emptyset$ 且 $V_1 \cup V_2 = f(X)$。由命题 2.26，存在 $Y$ 中的开集 $O_j$ 使得对 $j = 1,2$ 有 $V_j = O_j \cap f(X)$。令 $U_j := f^{-1}(O_j)$。则由定理 2.20，对 $j = 1,2$，$U_j$ 在 $X$ 中是开集。此外
\[
U_{1}\cup U_{2}=X~,\quad U_{1}\cap U_{2}=\varnothing\qquad\text{ 且 }\qquad U_{j}\neq\varnothing~,\quad j=1,2~,
\]

这对连通集 X 是不可能的。

\end{theorem}

\begin{corollary}

区间的连续像是连通的。

我们将在接下来的两节中说明，定理 4.4 与 4.5 是研究实函数的极其有用的工具。我们现在已经注意到这些定理的如下简单推论。

\end{corollary}

\begin{theorem}

（广义介值定理） 设 $X$ 是连通度量空间，$f: X \to \mathbb{R}$ 连续。则 $f(X)$ 是一个区间。特别地，$f$ 取到任意两个给定函数值之间的每个值。

\proofhead
这直接由定理 4.4 与 4.5 推出。
\end{theorem}

\subsection*{道路连通性}

设 $\alpha, \beta \in \mathbb{R}$ 满足 $\alpha < \beta$。连续函数 $w: [\alpha, \beta] \to X$ 称为连接端点 $w(\alpha)$ 与 $w(\beta)$ 的连续道路。

\bookfigure[0.51]{figures/06920d636f7edf9c.jpg}

度量空间 $X$ 称为道路连通的，若对每一对 $(x, y) \in X \times X$，存在 $X$ 中连接 $x$ 与 $y$ 的连续道路。度量空间的子集称为道路连通的，若它关于诱导度量是道路连通的度量空间。

\bookfigure[0.20]{figures/9880b37cd7303e2e.jpg}

\begin{proposition}

任何道路连通空间都是连通的。

\proofhead
反设存在一个道路连通但不连通的度量空间 $X$。则存在 $X$ 中的非空开集 $O_1$、$O_2$，使得 $O_1 \cap O_2 = \emptyset$ 且 $O_1 \cup O_2 = X$。选取 $x \in O_1$ 与 $y \in O_2$。由假设，存在道路 $w: [\alpha, \beta] \to X$ 使得 $w(\alpha) = x$ 且 $w(\beta) = y$。令 $U_j := w^{-1}(O_j)$。则由定理 2.20，$U_j$ 在 $[a, \beta]$ 中是开集。现在我们有 $\alpha$ 在 $U_1$ 中、$\beta$ 在 $U_2$ 中，以及 $U_1 \cap U_2 = \emptyset$ 和 $U_1 \cup U_2 = [\alpha, \beta]$，因此区间 $[a, \beta]$ 不是连通的。这与定理 4.4 矛盾。

设 $E$ 是赋范向量空间，$a, b \in E$。$E$ 的线性结构允许我们考虑 $E$ 中的“直线”道路：
\begin{equation}
v:[0,1]\to E~,\quad t\mapsto(1-t)a+t b~.
\end{equation}

我们用 $\left[\!\!\left[a, b\right]\!\!\right]$ 表示道路 v 的像。

$E$ 的子集 $X$ 称为 $\mathbf{凸}$ 的，若对每一对 $(a,b)\in X\times X$，$\left[\!\!\left[a,b\right]\!\!\right]$ 都包含于 $X$。

\bookfigure[0.21]{figures/2074e462412f57a3.jpg}

凸

\bookfigure[0.18]{figures/26da150c7ad9c445.jpg}
\end{proposition}

\begin{remark}

设 E 是赋范向量空间。

\begin{items}
\item

E 的每个凸子集都是道路连通的且连通的。

\proofhead
设 $X$ 是凸的，$a, b \in X$。则 (4.1) 定义了 $X$ 中连接 $a$ 与 $b$ 的一条道路。因此 $X$ 是道路连通的。于是命题 4.7 蕴含 $X$ 是连通的。
\item

对所有 $ a \in E $ 与 $ r > 0 $，球 $ \mathbb{B}_E(a, r) $ 和 $ \bar{\mathbb{B}}_E(a, r) $ 都是凸的。

\proofhead
对 $ x, y \in \mathbb{B}_E(a, r) $ 与 $ t \in [0, 1] $，我们有
\begin{equation*}
\begin{aligned}
\|(1-t)x+ty-a\|&=\|(1-t)(x-a)+t(y-a)\|\\&\leq(1-t)\left\|x-a\right\|+t\left\|y-a\right\|<(1-t)r+tr=r.
\end{aligned}
\end{equation*}

这个不等式蕴含 $ \left[\!\!\left[x, y\right]\!\!\right] $ 在 $ \mathbb{B}_E(a, r) $ 中。第二个断言可类似地证明。

\item

$ \mathbb{R} $ 的子集是凸的当且仅当它是区间。

\proofhead
设 $X \subseteq \mathbb{R}$ 是凸的。则由 (a)，$X$ 是连通的，从而由定理 4.4，$X$ 是区间。区间是凸的这一断言是显然的。
在 $\mathbb{R}^2$ 中，存在不是凸的连通集的简单例子。即便如此，在这些情形中，集合中的任意一对点似乎都可以用由有限多条直线段组成的道路连接起来。下面的定理表明，只要集合是开集，这一结论不仅在 $\mathbb{R}^2$ 中成立，而且在任何赋范向量空间中都成立。
设 $X$ 是赋范向量空间的子集。函数 $w: [\alpha, \beta] \to X$ 称为 $X$ 中的 $\text{折线道路}$\footnote{函数 $w: [\alpha, \beta] \to X$ 显然在每一点左连续且右连续，因此由命题 1.12 连续。于是折线道路特别地是一条道路。}，若存在 $n \in \mathbb{N}$ 与实数 $\alpha_0, \ldots, \alpha_{n+1}$ 使得 $\alpha = \alpha_0 < \alpha_1 < \cdots < \alpha_{n+1} = \beta$ 且
\end{items}

\bookfigure[0.24]{figures/2a4c1c292edcf470.jpg}
\[
w\big((1-t)\alpha_{j}+t\alpha_{j+1}\big)=(1-t)w(\alpha_{j})+t w(\alpha_{j+1})
\]

对所有 $ t \in [0,1] $ 和 $ j = 0, \ldots, n $ 成立。

\end{remark}

\begin{theorem}

设 X 是赋范向量空间的非空开连通子集。则 X 中任意一对点都可以用 X 中的一条折线道路连接。

\proofhead
设 $ a \in X $，并且

$ M := \{ x \in X ; \text{ 存在一条折线道路，在 } X \text{ 中连接 } x \text{ 与 } a \} $。

我们现在应用注 4.3 中描述的证明技巧。

\begin{subitems}
\item

因为 $ a \in M $，集合 M 非空。

\item

我们接着证明 $M$ 在 $X$ 中是开集。设 $x \in M$。由于 $X$ 是开集，存在某个 $r > 0$ 使得 $\mathbb{B}(x, r) \subseteq X$。由注 4.9(b)，对每个 $y \in \mathbb{B}(x, r)$，集合 $\left[\!\!\left[x, y\right]\!\!\right]$ 都包含于 $\mathbb{B}(x, r)$，从而也包含于 $X$。由于 $x \in M$，存在折线道路 $w: [\alpha, \beta] \to X$ 使得 $w(\alpha) = a$ 且 $w(\beta) = x$。

\end{subitems}

\bookfigure[0.47]{figures/c2d13d414199db3f.jpg}

现在按如下方式定义 $ \widetilde{w}: [\alpha, \beta + 1] \to X $：

\begin{equation*}
\widetilde{w}(t):=\left\{\begin{array}{c l}{w(t),}&{\quad t\in[\alpha,\beta]~,}\\ {(t-\beta)y+(\beta+1-t)x,}&{\quad t\in(\beta,\beta+1]~.}\\ \end{array}\right.
\end{equation*}

则 $ \tilde{w} $ 是 X 中连接 a 和 y 的折线道路。这表明 $ \mathbb{B}(x,r) $ 包含于 M，x 是 M 的内点，且 M 在 X 中是开集。

\begin{subitems}
\item[（iii）]

还要证明 $M$ 是闭集。设 $y \in X \backslash M$。由于 $X$ 是开集，存在某个 $r > 0$ 使得 $\mathbb{B}(y, r)$ 包含于 $X$。集合 $\mathbb{B}(y, r)$ 与 $M$ 必不交，因为若 $x \in \mathbb{B}(y, r) \cap M$，则由 (ii) 的论证，$X$ 中会存在连接 $a$ 与 $y$ 的折线道路，从而 $y$ 在 $M$ 中，与假设矛盾。于是 $y$ 是 $X \backslash M$ 的内点，并且由于 $y \in X \backslash M$ 是任意的，$X \backslash M$ 是开集。这蕴含 $M$ 在 $X$ 中是闭集。

\end{subitems}

\end{theorem}

\begin{corollary}

赋范向量空间的开子集是连通的当且仅当它是道路连通的。

\proofhead
这由命题 4.8 与定理 4.10 推出。
\end{corollary}

\subsection*{一般拓扑空间中的连通性}

在本节末尾，我们分析上述证明对度量存在性的依赖程度。

\begin{remark}

\begin{items}
\item

“连通”与“道路连通”的定义只依赖于拓扑，并且不使用度量。因此它们在任何拓扑空间中都是有效的。命题 4.2、4.5 与 4.8 也同样成立。特别地，当 X 是任意拓扑空间时，广义介值定理（定理 4.7）成立。

\item

存在连通但不是道路连通的空间的例子。正因如此，定理 4.10 特别有用。

\end{items}

\end{remark}

\exerciseshead

在下文中，X 是度量空间。

\begin{exercises}
\item

证明下列各条的等价性：

\begin{items}
\item

X 是连通的。

\item

不存在连续满射 $ X \to \{0,1\} $。

\end{items}

\item

设对指标集 $ A $ 中的每个 $ \alpha $，$ C_\alpha \subseteq X $ 是连通的。证明：若对所有 $ \alpha, \beta \in A $ 有 $ C_\alpha \cap C_\beta \neq \emptyset $，则 $ \bigcup_\alpha C_\alpha $ 是连通的。也就是说，两两相交的连通集的任意并是连通的。（提示：使用习题 1 并用反证法证明。）

\item

通过例子说明，连通集的交一般来说不是连通的。

\item

设 $ X_j $（$ j = 1, \ldots, n $）是度量空间。证明积 $ X_1 \times \cdots \times X_n $ 是连通的当且仅当每个 $ X_j $ 都是连通的。（提示：把 $ X \times Y $ 写成形如 $ (X \times \{y\}) \cup (\{x\} \times Y) $ 的集合的并。）

\item

证明连通集的闭包也是连通的。（提示：考虑连续函数 $ f: \overline{A} \to \{0,1\} $，并利用 $ f(\overline{A}) \subseteq \overline{f(A)} $（见习题 2.12）。）

\item

给定一个元素 $ x \in X $，集合

\[
K(x):=\bigcup_{Y\in M}Y\quad\text{ 其中 }\quad M:=\{Y\subseteq X\;;\;Y\text{ 是连通的且 }x\in Y\}
\]

由习题 2，是 X 中包含 x 的最大连通子集，因此称为 x 在 X 中的连通分支。证明下列各条：

\begin{items}
\item

$ \{K(x) : x \in X\} $ 是 $ X $ 的一个划分，也就是说，每个 $ x \in X $ 恰好包含在 $ X $ 的一个连通分支中。

\item

每个连通分支都是闭的。

\end{items}

\item

确定 Q 在 $ \mathbb{R} $ 中的所有连通分支。

\item

设 $ E = (E, \|\cdot\|) $ 是满足 $ \dim(E) \geq 2 $ 的赋范向量空间。证明 $ E \setminus \{0\} $ 与单位球面 $ S := \{x \in E; \|x\| = 1\} $ 都是连通的。

\item

证明下列度量空间 X 与 Y 不同胚（见习题 3.3）：

\begin{items}
\item

$ X := S^{1} $，$ Y := [0, 1] $。

\item

$ X := \mathbb{R} $，$ Y := \mathbb{R}^n $，$ n \geq 2 $。

\item

$ X := (0,1) \cup (2,3) $，$ Y := (0,1) \cup (2,3] $。
（提示：在每种情形中，从 X 中去掉一个或两个点。）

\end{items}

\item

证明所有实正交 $ n \times n $ 矩阵构成的集合 $ O(n) $ 不是连通的。（提示：函数 $ O(n) \to \{-1,1\} $，$ A \mapsto \det A $ 是连续的且是满射的（见习题 1.16）。）

\item

对 $ b_{j,k} \in \mathbb{R} $，$ 1 \leq j, k \leq n $，考虑双线性形式

\[
B:\mathbb{R}^{n}\times\mathbb{R}^{n}\to\mathbb{R},\quad(x,y)\mapsto\sum_{j,k=1}^{n}b_{j,k}x_{j}y_{k}.
\]

若对所有 $x \in \mathbb{R}^n \setminus \{0\}$ 有 $B(x,x) > 0$（或 $B(x,x) < 0$），则称 $B$ 是正定的（或负定的）。若 $B$ 既不正定也不负定，则它是不定的。证明下列各条：

\begin{items}
\item

若 $B$ 是不定的，则存在某个 $x \in S^{n-1}$ 使得 $B(x,x) = 0$。

\item

若 $B$ 是正定的，则存在某个 $\beta > 0$ 使得 $B(x, x) \geq \beta |x|^2$，$x \in \mathbb{R}^n$。（提示：对 (a) 用介值定理。对 (b) 用极值定理。）

\end{items}

\item

设 $E$ 是向量空间。设 $x_1, \ldots, x_n \in E$ 与 $\alpha_1, \ldots, \alpha_n \in \mathbb{R}^+$ 满足 $\sum_{j=1}^n \alpha_j = 1$。则称 $\sum_{j=1}^n \alpha_j x_j$ 为 $x_1, \ldots, x_n$ 的凸组合。证明下列各条：

\begin{items}
\item

E 的凸子集的任意交都是凸的。

\item

E 的子集 M 是凸的当且仅当 M 对凸组合封闭，也就是说，M 的点的每个凸组合都在 M 中。

\item

若 $E$ 是赋范向量空间且 $M \subseteq E$ 是凸的，则 $\bar{M}$ 与 $\overline{M}$ 也都是凸的。

\end{items}

\end{exercises}

\section{$ \mathbb{R} $ 上的函数}

我们关于连续性的抽象展开，在应用于 $ \mathbb{R} $ 上的实值函数时尤为富有成效。这当然是 $ \mathbb{R} $ 丰富结构的直接结果。

\subsection*{波尔查诺介值定理}

将广义介值定理应用于实值函数，便得到这一重要定理的波尔查诺原初版本。

\begin{theorem}

（波尔查诺介值定理）设 $ I \subseteq \mathbb{R} $ 是区间，$ f: I \to \mathbb{R} $ 连续。则 $ f(I) $ 是区间。也就是说，区间的连续像是区间。

\proofhead
这由定理 4.4 与定理 4.7 得出。■
\end{theorem}

在下文中，$ I $ 表示 $ \mathbb{R} $ 中的一个非空区间。

\begin{example}

\begin{items}
\item

若 $ f $ 不连续，或不定义在区间上，则波尔查诺介值定理中的断言不成立。这可由下面的函数图像加以说明：

\end{items}

\bookfigure[0.19]{figures/11db4795e29cd5e9.jpg}
\bookfigure[0.19]{figures/724bbdc316f8074a.jpg}

\begin{items}
\item[（b）]

若 $f: I \to \mathbb{R}$ 连续，且存在 $a, b \in I$ 使得 $f(a) < 0 < f(b)$，则存在某个介于 $a$ 与 $b$ 之间的 $\xi$，使得 $f(\xi) = 0$。

\end{items}

\bookfigure[0.44]{figures/7466f704e1f50f48.jpg}

\begin{items}
\item[（c）]

每个次数为奇数的多项式 $ p \in \mathbb{R}[X] $ 都有实零点。

\proofhead
不失一般性，可将 $ p $ 写成如下形式
\[
p=X^{2n+1}+a_{2n}X^{2n}+\cdots+a_{0}
\]

其中 $ n \in \mathbb{N} $，$ a_k \in \mathbb{R} $。于是

\[
p(x)=x^{2n+1}\Big(1+\frac{a_{2n}}{x}+\cdots+\frac{a_{0}}{x^{2n+1}}\Big)~,\qquad x\in\mathbb{R}^\times~.
\]

对于充分大的 $ R > 0 $，有

\[
1+\frac{a_{2n}}{R}+\cdots+\frac{a_{0}}{R^{2n+1}}\geq1-\frac{|a_{2n}|}{R}-\cdots-\frac{|a_{0}|}{R^{2n+1}}\geq\frac{1}{2},
\]

因此 $p(R) \geq R^{2n+1}/2 > 0$ 且 $p(-R) \leq -R^{2n+1}/2 < 0$。由于多项式函数连续，断言由（b）得出。■

\end{items}

\end{example}

\subsection*{单调函数}

$ \mathbb{R} $ 的序完备性对单调函数有着深远的影响。作为第一个例子，我们证明一个单调但未必连续的实函数在区间端点处的左、右极限的存在性。

\begin{proposition}

设 $ f: I \to \mathbb{R} $ 单调，$ \alpha := \inf I $，$ \beta := \sup I $。则

\begin{equation*}
\operatorname*{l i m}_{x\to\alpha+}f(x)=\left\{\begin{array}{c l}{\operatorname*{i n f}f(I),}&{\quad\text{若 }f\text{ 递增},}\\ {\operatorname*{s u p}f(I),}&{\quad\text{若 }f\text{ 递减},}\\ \end{array}\right.
\end{equation*}

且

\begin{equation*}
\operatorname*{l i m}_{x\to\beta-}{f(x)}=\left\{\begin{array}{c l}{\operatorname*{s u p}{f(I)}~,}&{\quad\text{若 }f\text{ 递增}~,}\\ {\operatorname*{i n f}{f(I)}~,}&{\quad\text{若 }f\text{ 递减}~.}\\ \end{array}\right.
\end{equation*}

\proofhead
设 $f$ 递增，且 $b := \sup f(I) \in \bar{\mathbb{R}}$。由 $b$ 的定义，对每个 $\gamma < b$，存在某个 $x_\gamma \in I$ 使得 $f(x_\gamma) > \gamma$。由于 $f$ 递增，有
\[
\gamma<f(x_{\gamma})\leq f(x)\leq b~,\qquad x\geq x_{\gamma}~.
\]

于是注 2.23 关于左极限的类似结论蕴含 $ \lim_{x \to \beta^-} f(x) = b $。区间左端点处以及递减函数情形的断言类似地证明。

为研究实函数的不连续点与连续延拓，我们需要下面的引理。

\end{proposition}

\begin{lemma}

设 $ D \subseteq \mathbb{R} $，$ t \in \mathbb{R} $，且

\[
\overline{D}_{t}:=\overline{D\cap(-\infty,t)}\cap\overline{D\cap(t,\infty)}.
\]

若 $\overline{D}_{t}$ 非空，则 $\overline{D}_{t}=\{t\}$，且存在 $D$ 中的序列 $(r_{n})$、$(s_{n})$ 使得

\[
r_{n}<t~,~s_{n}>t~,\quad n\in\mathbb{N}~,\qquad\text{且}\qquad\operatorname*{l i m}r_{n}=\operatorname*{l i m}s_{n}=t~.
\]

\proofhead
设 $ \overline{D}_t \neq \emptyset $，且 $ \tau \in \overline{D}_t $。则由 $ \overline{D}_t $ 的定义以及命题 2.9，存在 $ D $ 中的序列 $ (r_n) $ 与 $ (s_n) $，使得

\begin{subitems}
\item

$ r_n < t $，$ n \in \mathbb{N} $，且 $ \lim r_n = \tau $，（ii）$ s_n > t $，$ n \in \mathbb{N} $，且 $ \lim s_n = \tau $。由命题 II.2.7，（i）蕴含 $ \tau \leq t $，（ii）蕴含 $ \tau \geq t $。于是 $ \tau = t $，所有断言得证。

\end{subitems}

\end{lemma}

\begin{example}

\begin{items}
\item

设 $ D $ 是区间。则

\begin{equation*}
\overline{{D}}_{t}=\left\{\begin{array}{c c}{\{t\},}&{\quad t\in\mathring{D},}\\ {\varnothing,}&{\quad t\notin\mathring{D}.}\\ \end{array}\right.
\end{equation*}
\item

若 $ D = \mathbb{R}^\times $，则对每个 $ t \in \mathbb{R} $，有 $ \overline{D}_t = \{t\} $。■

\end{items}

现在考虑函数 $f: D \to X$，其中 $X = (X, d)$ 是度量空间，$D$ 是 $\mathbb{R}$ 的子集。设 $t_0 \in \mathbb{R}$ 满足 $\overline{D}_{t_0} \neq \emptyset$。若单侧极限 $f(t_0+) = \lim_{t \to t_0+} f(t)$ 与 $f(t_0-) = \lim_{t \to t_0-} f(t)$ 存在且不相等，则称 $t_0$ 为 $f$ 的一个 $\text{跳跃\textbf{不连续点}}$，并称 $d(f(t_0+), f(t_0-))$ 为该跳跃不连续点在 $t_0$ 处的 $\text{大小}$。

\bookfigure[0.26]{figures/471c65f02f45b142.jpg}
\bookfigure[0.26]{figures/9ba17e23b4aa3f85.jpg}
\end{example}

\begin{proposition}

若 $ f: I \to \mathbb{R} $ 单调，则 $ f $ 除至多可数个跳跃不连续点外连续。

\proofhead
只需考虑递增函数 $f: I \to \mathbb{R}$ 的情形。对 $t_0 \in \mathring{I}$，把命题 5.3 应用于限制函数 $f|I \cap (-\infty, t_0)$ 与 $f|I \cap (t_0, \infty)$ 中的每一个，即得 $\lim_{t \to t_0+} f(t)$ 与 $\lim_{t \to t_0-} f(t)$ 存在。由命题 1.12 与 1.13，只需证明集合
\[
M:=\left\{t_{0}\in\mathring{I};\;f(t_{0}-)\neq f(t_{0}+)\right\}
\]

可数。对每个 $t \in M$，有 $f(t-) < f(t+)$，故可选取某个 $r(t) \in \mathbb{Q} \cap (f(t-), f(t+))$。这定义了一个函数

\[
r\colon M\to\mathbb{Q}~,\quad t\mapsto r(t)~,
\]

由于 $f$ 递增，该函数必为单射。于是 $M$ 与 $\mathbb{Q}$ 的一个子集等势。特别地，由命题 I.6.7 与 I.9.4，$M$ 可数。

\end{proposition}

\subsection*{连续单调函数}

下面的重要定理表明，任何严格单调的连续函数都是单射，并具有定义在其像上的连续单调反函数。

\begin{theorem}

（单调函数的反函数定理）设 $ I \subseteq \mathbb{R} $ 是非空区间，$ f: I \to \mathbb{R} $ 连续且严格递增（或严格递减）。

\begin{subitems}
\item

$ J := f(I) $ 是区间。

\item

$f: I \to J$ 是双射。

\item

$ f^{-1} $：$ J \to I $ 连续且严格递增（或严格递减）。

\proofhead
断言（i）由定理 5.1 得出，（ii）是 $ f $ 严格单调性的直接结果。
为证（iii），设 $f$ 严格递增，并令 $g := f^{-1} : J \to I$。若 $s_1, s_2 \in J$ 满足 $s_1 < s_2$，则 $g(s_1) < g(s_2)$；因为否则
\[
s_{1}=f\big(g(s_{1})\big)\geq f\big(g(s_{2})\big)=s_{2}.
\]

于是 $ g $ 严格递增。
为证 $ g: J \to I $ 的连续性，只需考虑 $ J $ 含有多于一个点的情形，因为否则断言显然。设 $ g $ 在 $ s_0 \in J $ 处不连续。则存在 $ \varepsilon > 0 $ 与 $ J $ 中的序列 $ (s_n) $，使得

\begin{equation}
|s_{n}-s_{0}|<1/n\quad\text{且}\quad|g(s_{n})-g(s_{0})|\geq\varepsilon\;,\qquad n\in\mathbb{N}^{\times}\;.
\end{equation}

于是对所有 $ n \in \mathbb{N}^\times $，有 $ s_n \in [s_0 - 1, s_0 + 1] $；又因 $ g $ 递增，存在 $ \alpha, \beta \in \mathbb{R} $ 使得 $ \alpha < \beta $，且

\[
t_{n}:=g(s_{n})\in[\alpha,\beta]~.
\]

由波尔查诺-魏尔斯特拉斯定理，序列 $(t_n)$ 有收敛子序列 $(t_{n_k})_{k\in\mathbb{N}}$。设 $t_0$ 为该子序列的极限。则 $f$ 的连续性蕴含，当 $k \to \infty$ 时 $f(t_{n_k}) \to f(t_0)$。但由（5.1）的第一个断言，我们也知道 $f(t_{n_k}) = s_{n_k}$ 收敛于 $s_0$。于是 $s_0 = f(t_0)$，从而

\[
g(s_{n_{k}})=t_{n_{k}}\to t_{0}=g(s_{0})\ (k\to\infty)~.
\]

这与（5.1）的第二个断言矛盾，证明完成。■

\end{subitems}

\end{theorem}

\begin{example}

\begin{items}
\item

对每个 $ n \in \mathbb{N}^* $，函数

\[
\mathbb{R}^{+}\to\mathbb{R}^{+}\ ,\quad x\mapsto\sqrt[n]{x}
\]

连续\footnote{另见习题 II.2.7。}且严格递增。此外，$ \lim_{x \to \infty} \sqrt[n]{x} = \infty $。

\proofhead
对 $ n \in \mathbb{N}^x $，令 $ f: \mathbb{R}^+ \to \mathbb{R}^+ $ 由 $ t \mapsto t^n $ 定义。作为多项式函数的限制，$ f $ 连续。若 $ 0 \leq s < t $，则
\[
f(t)-f(s)=t^{n}-s^{n}=t^{n}\Big(1-\Big(\frac{s}{t}\Big)^{n}\Big)>0,
\]

这表明 $f$ 严格递增。最后 $\lim_{t \to \infty} f(t) = \infty$，故所有断言由定理 5.7 得出。■

\item

一般来说，若 $ I $ 不是区间，则定理 5.7（iii）中的连续性断言不成立。

\proofhead
例 1.9（c）中的函数 $ f: Z \to \mathbb{R} $ 连续且严格递增，但其反函数不连续。
单调函数的反函数定理的更多重要应用将在下一节出现。
\end{items}

\end{example}

\exerciseshead

在下文中，$ I $ 是含有多于一个点的紧致区间。

\begin{exercises}
\item

设 $ f: I \to I $ 连续。证明 $ f $ 有不动点，即存在某个 $ \xi \in I $ 使得 $ f(\xi) = \xi $。

\item

设 $ f: I \to \mathbb{R} $ 连续且为单射。证明 $ f $ 严格单调。

\item

设 $D$ 是 $\mathbb{R}$ 的开子集，$f: D \to \mathbb{R}$ 连续且为单射。证明 $f: D \to f(D)$ 是同胚。\footnote{见习题 3.3。}

\item

设 $ \alpha: \mathbb{N} \to \mathbb{Q} $ 是双射，且对 $ x \in \mathbb{R} $，令 $ N_x $ 为集合 $ \{k \in \mathbb{N} : \alpha(k) \leq x\} $。设 $ (y_n) $ 是 $ (0, \infty) $ 中的序列，满足 $ \sum y_n < \infty $。定义

\[
f:\mathbb{R}\to\mathbb{R},\quad x\mapsto\sum_{k\in N_{x}}y_{k}.
\]

证明下列结论：\footnote{本习题表明命题 5.6 不能加强。}

\begin{items}
\item

$ f $ 严格单调。

\item

$ f $ 在每个无理数处连续。

\item

在每个有理数 $ q $ 处，存在大小为 $ y_n $ 的跳跃不连续点，其中 $ n = \alpha^{-1}(q) $。

\end{items}

\item

考虑函数

\begin{equation*}
f\colon[0,1]\to[0,1],\quad x\mapsto\left\{\begin{array}{c l}{x\quad,\quad}&{x\text{ 为有理数}\quad,}\\ {1-x\quad,\quad}&{x\text{ 为无理数}\quad.}\end{array}\right.
\end{equation*}

证明下列结论：

\begin{items}
\item

$ f $ 是双射。

\item

$ f $ 在 $[0,1]$ 的任何子区间上都不单调。

\item

$ f $ 仅在 $ x = 1/2 $ 处连续。

\end{items}

\item

设 $ f_{0} := $ 锯齿函数（见习题 1.1），且

\[
F(x):=\sum_{n=0}^{\infty}4^{-n}f_{0}(4^{n}x)~,\qquad x\in\mathbb{R}~.
\]

证明下列结论：

\begin{items}
\item

$ F $ 良定义。

\item

$ F $ 在任何区间上都不单调。

\item

$ F $ 连续。
（提示：（a）对每个 $ x \in \mathbb{R} $，为 $ \sum 4^{-n} f_0(4^n x) $ 找出一个收敛的优级数。

\end{items}

\begin{items}
\item[（b）]

对所有 $ x \in \mathbb{R} $ 与 $ n \in \mathbb{N} $，令 $ f_n(x) := 4^{-n} f_0(4^n x) $。对 $ k \in \mathbb{Z} $ 与 $ m \in \mathbb{N}^* $，令 $ a := k \cdot 4^{-m} $，$ h := 4^{-2m-1} $。则

\[
f_{n}(a)=0,\quad n\geq m,\qquad\text{且}\qquad f_{n}(a\pm h)=0,\quad n\geq2m+1,
\]

因此 $ F(a \pm h) - F(a) \geq h $。最后用 $ k \cdot 4^{-m} $（$ k \in \mathbb{Z} $，$ m \in \mathbb{N}^\times $）逼近任意 $ x \in \mathbb{R} $。

\item[（c）]

对 $ x, y \in \mathbb{R} $ 与 $ m \in \mathbb{N}^\times $，有 $ |F(x) - F(y)| \leq \sum_{k=0}^{m} |f_k(x) - f_k(y)| + 4^{-m}/3 $。

\end{items}

\item

设 $f: I \to \mathbb{R}$ 单调。证明 $\omega_f(x) = |f(x+) - f(x-)|$，其中 $\omega_f(x)$ 如习题 1.17 中所定义。

\end{exercises}

\section{指数函数与相关函数}

在本节（相当长的一节）中，我们研究数学中最重要的函数之一——指数函数。其重要性已体现在它与三角函数和对数函数的密切关系中，我们也将考察后两者。

\subsection*{欧拉公式}

在第 II 章中，我们利用指数级数定义了指数函数。

\[
\exp(z):=e^{z}:=\sum_{n=0}^{\infty}\frac{z^{n}}{n!}=1+z+\frac{z^{2}}{2!}+\frac{z^{3}}{3!}+\cdots~,\qquad z\in\mathbb{C}~.
\]

用记号 $ e^{z} $ 表示 $ \exp(z) $ 的用法由例 II.8.12(b) 提供依据。与这个级数相关联的是余弦级数

\[
\sum(-1)^{n}\frac{z^{2n}}{(2n)!}=1-\frac{z^{2}}{2!}+\frac{z^{4}}{4!}-+\cdots,
\]

以及正弦级数

\[
\sum(-1)^{n}\frac{z^{2n+1}}{(2n+1)!}=z-\frac{z^{3}}{3!}+\frac{z^{5}}{5!}-+\cdots.
\]

我们将证明：与指数级数类似，余弦级数与正弦级数处处绝对收敛。由这些级数定义的函数，

\[
\cos:\mathbb{C}\to\mathbb{C}~,\quad z\mapsto\sum_{n=0}^{\infty}(-1)^{n}\frac{z^{2n}}{(2n)!}
\]

以及

\[
\sin:\mathbb{C}\to\mathbb{C}\;,\quad z\mapsto\sum_{n=0}^{\infty}(-1)^{n}\frac{z^{2n+1}}{(2n+1)!}\;,
\]

分别称为余弦函数与正弦函数。\footnote{我们稍后将会看到，这些定义给出熟悉的三角函数。}

\begin{theorem}

\begin{subitems}
\item

指数级数、余弦级数与正弦级数具有无穷收敛半径。

\item

函数 $ \exp $、$ \cos $、$ \sin $ 在实变元上取实值。

\item

指数函数的加法定理成立：

\[
e^{w+z}=e^{w}e^{z}~,\qquad w,z\in\mathbb{C}~.
\]
\item

欧拉公式成立：

\begin{equation}
e^{i z}=\cos z+i\sin z~,\qquad z\in\mathbb{C}~.
\end{equation}
\item

函数 $ \exp $、$ \cos $ 与 $ \sin $ 在 C 上连续。

\proofhead
(i) 在例 II.8.7(c) 中，我们已经证明了指数级数具有收敛半径 $ \infty $。因此阿达马公式给出
\[
\infty=\frac{1}{\varlimsup_{n\to\infty}\sqrt[n]{1/n!}}=\varliminf_{n\to\infty}\sqrt[n]{n!}.
\]

由定理 II.5.7，序列 $ ( \sqrt[n]{n!} )_{n \in \mathbb{N}} $ 及其所有子序列都收敛到 $ \infty $。因此

\[
\frac{1}{\varlimsup_{n\to\infty}\sqrt[2n]{1/(2n)!}}=\lim_{n\to\infty}\sqrt[2n]{(2n)!}=\infty
\]

以及

\[
\frac{1}{\varlimsup_{n\to\infty}\sqrt[2n+1]{1/(2n+1)!}}=\lim_{n\to\infty}\sqrt[2n+1]{(2n+1)!}=\infty,
\]

从而由阿达马公式，余弦级数与正弦级数具有无穷收敛半径。

\end{subitems}

\begin{subitems}
\item[（ii）]

因为 $ \mathbb{R} $ 是域，所以当 z 为实数时，上述级数的所有部分和均为实数。由于 $ \mathbb{R} $ 在 $ \mathbb{C} $ 中是闭的，断言得证。

\item[（iii）]

这已在例 II.8.12(a) 中证明。

\item[（iv）]

对 $ n \in \mathbb{N} $，我们有

\[
i^{2n}=(i^{2})^{n}=(-1)^{n}\quad\text{且}\quad i^{2n+1}=i\cdot i^{2n}=i\cdot(-1)^{n}.
\]

因此，由命题 II.7.5。

\[
e^{i z}=\sum_{n=0}^{\infty}\frac{(i z)^{n}}{n!}=\sum_{k=0}^{\infty}\frac{(i z)^{2k}}{(2k)!}+\sum_{k=0}^{\infty}\frac{(i z)^{2k+1}}{(2k+1)!}=\cos z+i\sin z
\]

对所有 $ z \in \mathbb{C} $。

\item[（v）]

这由命题 1.7 可得。

\end{subitems}

\end{theorem}

\begin{remark}

\begin{items}
\item

余弦是偶函数，正弦是奇函数，即\footnote{见习题 II.9.7。}

\begin{equation}
\cos(z)=\cos(-z)\quad\text{且}\quad\sin(z)=-\sin(-z),\qquad z\in\mathbb{C}.
\end{equation}
\item

由 (a) 与欧拉公式 (6.1) 可得

\begin{equation}
\cos(z)=\frac{e^{i z}+e^{-i z}}{2}\;,\quad\sin(z)=\frac{e^{i z}-e^{-i z}}{2i}\;,\qquad z\in\mathbb{C}\;.
\end{equation}
\item

对 $ w, z \in \mathbb{C} $，我们有

\[
e^{z}\neq0,\quad e^{-z}=1/e^{z},\quad e^{z-w}=e^{z}/e^{w},\quad\overline{{e^{z}}}=e^{\overline{{z}}}.
\]

\proofhead
由加法定理可得 $ e^z e^{-z} = e^{z-z} = e^0 = 1 $，前三个断言由此得证。
由例 1.3(i)，函数 $\mathbb{C} \to \mathbb{C}$，$w \mapsto \overline{w}$ 连续。于是定理 1.4 蕴含
\[
\overline{e^{z}}=\lim_{n\to\infty}\sum_{k=0}^{n}\frac{z^{k}}{k!}=\lim_{n\to\infty}\sum_{k=0}^{n}\frac{\overline{z}^{k}}{k!}=e^{\overline{z}}
\]

对所有 $ z \in \mathbb{C} $。■

\item

对所有 $ x \in \mathbb{R} $，$ \cos(x) = \mathrm{Re}(e^{ix}) $ 且 $ \sin(x) = \mathrm{Im}(e^{ix}) $。

\proofhead
这由欧拉公式与定理 6.1(ii) 可得。■
\end{items}

在下面的命题中，我们用“三角函数”这一名称指称余弦与正弦。这一用法在注 6.18 之后得到说明。

\end{remark}

\qyxffntextdef{1}{在不致产生误解时，通常把 $\cos(z)$ 与 $\sin(z)$ 写作 $\cos z$ 与 $\sin z$。}
\begin{proposition}

（三角函数的加法定理）对所有 $ z, w \in \mathbb{C} $，我们有\qyxffnmark{1}

\begin{subitems}
\item
\[
\cos(z\pm w)=\cos z\cos w\mp\sin z\sin w,
\]
\[
\sin(z\pm w)=\sin z\cos w\pm\cos z\sin w.
\]
\item

$ \sin z - \sin w = 2 \cos \frac{z + w}{2} \sin \frac{z - w}{2} $，

\[
\cos z-\cos w=-2\sin\frac{z+w}{2}\sin\frac{z-w}{2}.
\]

\proofhead
(i) (6.3) 中的公式与指数函数的加法定理给出
\begin{equation*}
\begin{aligned}
\cos z\cos w-\sin z\sin w&=\frac{1}{4}\big\{(e^{iz}+e^{-iz})(e^{iw}+e^{-iw})+(e^{iz}-e^{-iz})(e^{iw}-e^{-iw})\big\}\\&=\frac{1}{2}\big\{e^{i(z+w)}+e^{-i(z+w)}\big\}=\cos(z+w)
\end{aligned}
\end{equation*}

对所有 $ z, w \in \mathbb{C} $。再利用 (6.2)，可得

\[
\cos(z-w)=\cos z\cos w+\sin z\sin w~,\qquad z,w\in\mathbb{C}~.
\]
\end{subitems}

(i) 中的第二个公式可类似证明。

(ii) 对 $ z, w \in \mathbb{C} $，令 $ u := (z + w)/2 $ 与 $ v := (z - w)/2 $。则 $ u + v = z $ 且 $ u - v = w $，于是利用 (i)，我们得到

\begin{equation*}
\begin{aligned}
\sin z-\sin w&=\sin(u+v)-\sin(u-v)=2\cos u\sin v\\&=2\cos\frac{z+w}{2}\sin\frac{z-w}{2}.
\end{aligned}
\end{equation*}

(ii) 中的第二个公式可类似证明。

\end{proposition}

\begin{corollary}

对 $ z \in \mathbb{C} $，$ \cos^2 z + \sin^2 z = 1 $。

\proofhead
在命题 6.3(i) 中令 z = w，我们得到
\[
\cos^{2}z+\sin^{2}z=\cos(z-z)=\cos(0)=1,
\]

这证明了断言。■

\end{corollary}

若把 $ z \in \mathbb{C} $ 写成 $ z = x + iy $ 的形式，其中 $ x, y \in \mathbb{R} $，则 $ e^z = e^x e^{iy} $。这一简单观察表明，指数函数完全由实指数函数 $ \exp_{\mathbb{R}} := \exp | \mathbb{R} \text{ 以及 } \exp \text{ 在 } i\mathbb{R} \text{ 上的限制} $ 所确定，即由 $ \exp_{i\mathbb{R}} := \exp | i\mathbb{R} $ 所确定。因此，为理解“复”指数函数 $ \exp: \mathbb{C} \to \mathbb{C} $，我们先分别研究这两个函数。

\subsection*{实指数函数}

我们在下一个命题中汇集函数 $ \exp_{\mathbb{R}} $ 最重要的定性性质。

\begin{proposition}

\begin{subitems}
\item

若 x < 0，则 0 < $ e^{x} $ < 1。若 x > 0，则 1 < $ e^{x} $ < $ \infty $。

\item

$ \exp_{\mathbb{R}}: \mathbb{R} \to \mathbb{R}^+ $ 严格递增。

\item

对每个 $ \alpha \in \mathbb{Q} $，

\[
\lim_{x\to\infty}\frac{e^{x}}{x^\alpha}=\infty,
\]

即指数函数的增长比任何幂函数都快。

\item

$ \lim_{x \to -\infty} e^x = 0 $。

\proofhead
(i) 由
\[
e^{x}=1+\sum_{n=1}^{\infty}\frac{x^{n}}{n!},\qquad x\in\mathbb{R},
\]

可见对所有 $ x > 0 $，$ e^x > 1 $。若 $ x < 0 $，则 $ -x > 0 $，从而 $ e^{-x} > 1 $。这意味着 $ e^x = e^{-(-x)} = 1/e^{-x} \in (0,1) $。

\end{subitems}

\begin{subitems}
\item[（ii）]

设 $ x, y \in \mathbb{R} $ 满足 $ x < y $。由于 $ e^x > 0 $ 且 $ e^{y-x} > 1 $，由此可得

\[
e^{y}=e^{x+(y-x)}=e^{x}e^{y-x}>e^{x}.
\]
\item[（iii）]

只需考虑 $ \alpha > 0 $ 的情形。令 $ n := \lfloor \alpha \rfloor + 1 $。由指数级数可得 $ e^x > x^{n+1}/(n+1) $! 对所有 $ x > 0 $ 成立。因此

\[
\frac{e^{x}}{x^{\alpha}}>\frac{e^{x}}{x^{n}}>\frac{x}{(n+1)!},\qquad x>0,
\]

这证明了断言。

\item[（iv）]

若在 (iii) 中令 $ \alpha = 0 $，我们得到 $ \lim_{x \to \infty} e^x = \infty $。因此

\[
\lim_{x\to-\infty}e^{x}=\lim_{y\to\infty}e^{-y}=\lim_{y\to\infty}\frac{1}{e^{y}}=0,
\]

于是所有断言都得证。

\end{subitems}

\end{proposition}

\subsection*{对数与幂函数}

由命题 6.5 我们有

$ \exp_{\mathbb{R}}: \mathbb{R} \to \mathbb{R}^+ \text{ 是连续的且严格递增的，并且 } \exp(\mathbb{R}) = (0, \infty) $

因此，由定理 5.7，实指数函数有定义在 $(0,\infty)$ 上的连续且严格递增的反函数。这个反函数称为（自然）对数，记作 log，即

\[
\operatorname{l o g}:=\left(\operatorname{e x p}_{\mathbb{R}}\right)^{-1}:(0,\infty)\to\mathbb{R}~.
\]
\bookfigure[0.20]{figures/48f6700599a70d14.jpg}

特别地，$ \log 1 = 0 $ 且 $ \log e = 1 $。

\begin{theorem}

（对数函数的加法定理）对所有 $ x, y \in (0, \infty) $，

\[
\log(xy)=\log x+\log y\quad\text{ 且 }\quad\log(x/y)=\log x-\log y.
\]

\proofhead
设 $x, y \in (0, \infty)$。对 $a := \log x$ 与 $b := \log y$，我们有 $x = e^a$ 与 $y = e^b$。于是指数函数的加法定理蕴含 $xy = e^a e^b = e^{a+b}$ 与 $x/y = e^a / e^b = e^{a-b}$，断言由此得证。
\end{theorem}

\qyxffnnumdef{1}
\begin{proposition}

对所有 $ a > 0 $ 与 $ r \in \mathbb{Q} $，\qyxffnmark{1}

\begin{equation}
a^{r}=e^{r\log a}.
\end{equation}

\proofhead
由定义，我们有 $ a = e^{\log a} $，于是定理 6.1(iii) 蕴含对所有 $ n \in \mathbb{N} $，$ a^n = (e^{\log a})^n = e^{n(\log a)} $。此外，
\[
a^{-n}=(e^{\log a})^{-n}=\frac{1}{(e^{\log a})^{n}}=\frac{1}{e^{n\log a}}=e^{-n\log a},\qquad n\in\mathbb{N}.
\]

现在令 $x := e^{\frac{1}{n} \log a}$。则 $x^n = e^n(\frac{1}{n} \log a) = e^{\log a} = a$，于是由命题 I.10.9，对所有 $n \in \mathbb{N}^\times$，$e^{\frac{1}{n} \log a} = a^{\frac{1}{n}}$。

现在设 $ r \in \mathbb{Q} $。则存在 $ p \in \mathbb{Z} $ 与 $ q \in \mathbb{N}^\times $ 使得 $ r = p/q $。由上我们有

\[
a^{r}=a^{\frac{p}{q}}=\left(a^{\frac{1}{q}}\right)^{p}=\left(e^{\frac{1}{q}\log a}\right)^{p}=e^{\frac{p}{q}\log a}=e^{r\log a},
\]

这就完成了 (6.4) 的证明。■

\end{proposition}

\qyxffntextof{1}{注意，(6.4) 的左端是正数 $a$ 的 $r$ 次幂（其定义见注 I.10.10(d)），而右端是指数函数在 $r\log a \in \mathbb{R}$ 处的值。还要注意，当 $a=e$ 时，(6.4) 化为例 II.8.12(b)。}

设 $a > 0$。迄今为止，我们只对有理指数 $r$ 定义了 $a^r$，并且对这样的指数我们已经证明了 $a^r = e^{r \log a}$。由于 $e^{r \log a}$ 对任意实数 $r \in \mathbb{R}$ 都有定义，这提示了一个显然的推广。具体地，我们定义

\[
a^{x}:=e^{x\log a}~,\qquad x\in\mathbb{R}~,\quad a>0~.
\]

\begin{proposition}

对所有 $ a, b > 0 $ 与 $ x, y \in \mathbb{R} $，

\[
a^{x}a^{y}=a^{x+y},\quad\frac{a^{x}}{a^{y}}=a^{x-y},\quad a^{x}b^{x}=(ab)^{x},\quad\frac{a^{x}}{b^{x}}=\left(\frac{a}{b}\right)^{x},
\]
\[
\log(a^{x})=x\log a~,\quad(a^{x})^{y}=a^{xy}~.
\]

\proofhead
例如。
\[
a^{x}a^{y}=e^{x\log a}e^{y\log a}=e^{(x+y)\log a}=a^{x+y}
\]

以及

\[
(a^{x})^{y}=(e^{x\log a})^{y}=e^{xy\log a}=a^{xy}.
\]

其余断言可类似证明。■

\end{proposition}

\begin{proposition}

对所有 $ \alpha > 0 $，

\[
\lim_{x\to\infty}\frac{\log x}{x^{\alpha}}=0\quad\text{ 且 }\quad\lim_{x\to0+}x^{\alpha}\log x=0.
\]

特别地，对数函数的增长比任何（任意小的）正幂函数都慢。

\proofhead
由于对数递增，由命题 6.5(iii) 可得
\[
\lim_{x\to\infty}\frac{\log x}{x^{\alpha}}=\lim_{x\to\infty}\frac{\log x}{e^{\alpha\log x}}=\lim_{y\to\infty}\frac{y}{e^{\alpha y}}=\frac{1}{\alpha}\lim_{t\to\infty}\frac{t}{e^{t}}=0.
\]

对第二个极限，我们有

\[
\lim_{x\to0+}x^{\alpha}\log x=\lim_{y\to\infty}\left(\frac{1}{y}\right)^{\alpha}\log\frac{1}{y}=-\lim_{y\to\infty}\frac{\log y}{y^{\alpha}}=0,
\]

这证明了断言。■

\end{proposition}

注意到命题 6.5(iii) 对 $ \alpha \in \mathbb{R} $ 也成立。

\subsection*{$ i\mathbb{R} $ 上的指数函数}

函数 $ \exp_{i\mathbb{R}} $ 与实指数函数 $ \exp_{\mathbb{R}} $ 的性质完全不同。例如，$ \exp_{\mathbb{R}} $ 严格递增，而我们将证明 $ \exp_{i\mathbb{R}} $ 是周期函数。在确定其周期的过程中，我们将定义常数 $ \pi $。为证明这些断言，我们首先需要几个引理。

\begin{lemma}

对所有 $ t \in \mathbb{R} $，$ |e^{it}| = 1 $。

\proofhead
由于对所有 $ z \in \mathbb{C} $，$ \overline{e^z} = e^{\overline{z}} $，我们有
\[
|e^{i t}|^{2}=e^{i t}\overline{{(e^{i t})}}=e^{i t}e^{-i t}=e^{0}=1\;,\qquad t\in\mathbb{R}\;,
\]

断言由此得证。■

\end{lemma}

有时考虑如下函数比考虑 $ \exp_{i\mathbb{R}} $ 更为有用：

\[
\mathrm{c i s}\colon\mathbb{R}\to\mathbb{C}~,\quad t\mapsto e^{i t}~.
\]

引理 6.10 表明 cis 的像包含在 $ S^1 := \{ z \in \mathbb{C} ; \, |z| = 1 \} $ 中。在下一个引理中，我们加强这一结果，并证明 cis 的像是整个 $ S^1 $。

\begin{lemma}

cis( $  \mathbb{R}  $) =  $  S^1  $。

证明

\begin{subitems}
\item

第一步我们证明余弦函数的像是

\begin{equation}
\cos(\mathbb{R})=\mathrm{p r}_{1}\left[\mathrm{c i s}(\mathbb{R})\right]=\left[-1,1\right].
\end{equation}

(6.5) 中的第一个等式是欧拉公式的显然推论。为证明第二个等式，令 $ I := \cos(\mathbb{R}) $。于是由波尔查诺介值定理（定理 5.1）可知 $ I $ 是一个区间。此外，由引理 6.10 我们知道

\[
I=\mathrm{p r}_{1}\left(\mathrm{c i s}(\mathbb{R})\right)\subseteq[-1,1].
\]

当然，$1 = \cos(0)$ 属于 $I$，但 $I = \{1\}$ 不可能，因为若对所有 $z \in \mathbb{R}$ 有 $\cos(z) = 1$，则由推论 II.9.9，余弦级数将为 $1 + 0z + 0z^2 + \cdots$。于是 $I$ 具有形式

\[
I=\left[a,1\right]\quad\text{或}\quad I=\left(a,1\right]
\]

对某个适当的 $ a \in [-1,1) $。
假设 $a$ 不等于 $-1$。由于 $a_0 := (a + 1)/2$ 属于 $I$，存在某个 $t_0 \in \mathbb{R}$ 使得 $a_0 = \cos t_0$。令

\[
z_{0}:=\mathrm{c i s}(t_{0})=\cos t_{0}+i\sin t_{0}.
\]

则，由推论 6.4，

\begin{equation*}
\begin{aligned}
\mathrm{pr}_{1}(z_{0}^{2})&=\mathrm{Re}\big((\cos t_{0}+i\sin t_{0})^{2}\big)=\cos^{2}t_{0}-\sin^{2}t_{0}\\&=2\cos^{2}t_{0}-1=2a_{0}^{2}-1=a-\frac{1-a^{2}}{2}<a,
\end{aligned}
\end{equation*}

因为由假设 $ a^{2} < 1 $。不等式 $ \mathrm{pr}_{1}(z_{0}^{2}) < a $ 与 $ \mathrm{pr}_{1}(z_{0}^{2}) = \mathrm{pr}_{1}(e^{2it_{0}}) $ 属于 I 这一事实矛盾。因此我们断定 $ a = -1 $。
为完成 (6.5) 的证明，还需证明 -1 属于 $I$。我们知道存在某个 $t_0 \in \mathbb{R}$ 使得 $\cos t_0 = 0$。由于 $\sin^2 t_0 = 1 - \cos^2 t_0 = 1$，这蕴含 $z_0 = e^{it_0} = i \sin t_0 = \pm i$。因此

\[
-1=\mathrm{pr}_{1}(-1)=\mathrm{pr}_{1}(z_{0}^{2})=\mathrm{pr}_{1}(e^{2i t_{0}})\in I,
\]

如所断言。

\item

接下来我们证明 $ S^1 \subseteq \mathrm{cis}(\mathbb{R}) $。若 $ z \in S^1 $，则

\[
\mathrm{R e}z\in[-1,1]=\mathrm{p r}_{1}\left(\mathrm{c i s}(\mathbb{R})\right),
\]

于是由 (i)，存在某个 $t \in \mathbb{R}$ 使得 $\mathrm{Re} z = \mathrm{Re} e^{it}$。此外，由 $|z| = 1 = |e^{it}|$ 可知，要么 $z = e^{it}$，要么 $\overline{z} = e^{it}$。在第一种情形，$z \in \mathrm{cis}(\mathbb{R})$ 显然。否则 $\overline{z} = e^{it}$，而由 $z = \overline{\overline{z}} = \overline{e^{it}} = e^{-it}$ 可知，此时也有 $z \in \mathrm{cis}(\mathbb{R})$。于是我们证明了 $S^{1} \subseteq \mathrm{cis}(\mathbb{R})$。这连同引理 6.10 蕴含 $\mathrm{cis}(\mathbb{R}) = S^{1}$。■

\end{subitems}

\end{lemma}

\begin{lemma}

集合 $ M := \{t > 0\}; e^{it} = 1 $ 有最小元。

证明

\begin{subitems}
\item

首先我们证明 $ M $ 非空。由引理 6.11，存在某个 $ t \in \mathbb{R}^\times $ 使得 $ e^{it} = -1 $。因为

\[
e^{-it}=\frac{1}{e^{it}}=\frac{1}{-1}=-1,
\]

我们可以假设 $t>0$。则 $e^{2it}=(e^{it})^2=(-1)^2=1$，且 $M$ 非空。

\item

接下来我们证明 $ M $ 在 $ \mathbb{R} $ 中是闭的。为证明这一点，取 $ M $ 中收敛到 $ t^* \in \mathbb{R} $ 的序列 $ (t_n) $。由于 $ t_n $ 对所有 $ n $ 都为正，我们有 $ t^* \geq 0 $。此外，cis 的连续性蕴含

\[
e^{i t^{*}}=\mathrm{c i s}(t^{*})=\mathrm{c i s}(\lim t_{n})=\lim\mathrm{c i s}(t_{n})=1.
\]

为证明 $M$ 是闭的，还需证明 $t^*$ 为正。反设 $t^* = 0$。则存在某个 $m \in \mathbb{N}$ 使得 $t_m \in (0,1)$。由欧拉公式我们有 $1 = e^{it_m} = \cos t_m + i \sin t_m$，从而 $\sin t_m = 0$。
把推论 II.7.9 应用于正弦级数

\[
\sin t=t-\frac{t^{3}}{6}+\frac{t^{5}}{5!}-+\cdots
\]

我们得到

\begin{equation}
\sin t\geq t(1-t^{2}/6)~,\qquad0<t<1~.
\end{equation}

对 $t_m$，这给出 $0 = \sin t_m \geq t_m (1 - t_m^2/6) > 5t_m/6$，矛盾。因此 $M$ 是闭的。

\item

由于 M 是非空且下有界的闭集，它有最小元。

\end{subitems}

\end{lemma}

\subsection*{$ \pi $ 的定义及其推论}

前面的引理使得可以如下定义一个数 $ \pi $：

\[
\pi:=\frac{1}{2}\min\{t>0;e^{it}=1\}.
\]

我们将在第 VI.5 节看到，这样定义的数 $\pi$ 具有通常的几何意义，例如作为单位圆的面积。然而目前，$\pi$ 只是满足 $e^{2\pi i}=1$ 的最小正实数。

考虑数 $ e^{i\pi} $。我们有 $ (e^{i\pi})^2 = e^{2\pi i} = 1 $，于是 $ e^{i\pi} = \pm1 $。由 $ \pi $ 的定义，情形 $ e^{i\pi} = 1 $ 不可能，于是 $ e^{i\pi} = -1 $。这还蕴含 $ e^{-i\pi} = 1/e^{i\pi} = -1 $。这些特殊情形可用来确定所有满足 $ e^z = 1 $ 或 $ e^z = -1 $ 的其他 $ z \in \mathbb{C} $。

\begin{proposition}

\begin{subitems}
\item

$ e^z = 1 \iff z \in 2\pi i \mathbb{Z} $。

\item

$ e^z = -1 \Leftrightarrow z \in \pi i + 2\pi i\mathbb{Z} $。

\proofhead
(i) “⇐” 对所有 $ k \in \mathbb{Z} $，我们有 $ e^{2\pi ik} = (e^{2\pi i})^k = 1 $。
“⇒” 设 $ z = x + iy $（$ x, y \in \mathbb{R} $）满足 $ e^z = 1 $。则
\[
1=\left|e^{z}\right|=\left|e^{x}\right|\left|e^{iy}\right|=e^{x},
\]

从而 $ x = 0 $。若 $ k \in \mathbb{Z} $ 与 $ r \in [0, 2\pi) $ 满足 $ y = 2\pi k + r $，则

\[
1=e^{i y}=e^{2\pi k i}e^{i r}=e^{i r}.
\]

$ \pi $ 的定义蕴含 $ r = 0 $，于是 $ z = 2\pi i k $，其中 $ k \in \mathbb{Z} $。

\end{subitems}

(ii) 由于 $ e^{-i\pi} = -1 $，我们有 $ e^z = -1 $ 当且仅当 $ e^{z-i\pi} = e^z e^{-i\pi} = 1 $。由 (i)，$ e^{z-i\pi} = 1 $ 当且仅当 $ z - i\pi = 2\pi ik $，即对某个 $ k \in \mathbb{Z} $，$ z = i\pi + 2\pi ik $。■

\end{proposition}

由命题 6.13(i) 我们有，对所有 $ k \in \mathbb{Z} $，$ e^{z+2\pi ik} = e^z e^{2\pi ik} = e^z $，从而得到下面的推论。

\qyxffntextdef{1}{若 $ E $ 是向量空间而 $ M $ 是集合，则称 $ f: E \to M $ 以 $ p \in E \setminus \{0\} $ 为\textbf{周期}，若对所有 $ x \in E $ 有 $ f(x+p) = f(x) $。}
\begin{corollary}

指数函数是以 $ 2\pi i $ 为周期的周期函数\qyxffnmark{1}，即

\[
e^{z}=e^{z+2\pi i k},\qquad z\in\mathbb{C},\quad k\in\mathbb{Z}.
\]

利用命题 6.13 我们还可证明，函数 cis 在长度为 $ 2\pi $ 的半开区间上是双射。

\end{corollary}

\begin{proposition}

对每个 $ a \in \mathbb{R} $，函数

\begin{equation*}
\begin{aligned}
&\mathrm{cis}\mid[a,a+2\pi):[a,a+2\pi)\to S^{1},\\&\mathrm{cis}\mid(a,a+2\pi):(a,a+2\pi]\to S^{1}
\end{aligned}
\end{equation*}

都是双射。

证明

\begin{subitems}
\item

设对某个 $ s, t \in \mathbb{R} $ 有 cis $ t = \text{cis } s $。由于 $ e^{i(t-s)} = 1 $，由命题 6.13，存在某个 $ k \in \mathbb{Z} $ 使得 $ t = s + 2\pi k $。这意味着上述每个函数都是单射。

\item

设 $ z \in S^1 $。由引理 6.11，存在某个 $ t \in \mathbb{R} $ 使得 $ \mathrm{cis} t = z $。此外存在 $ k_1, k_2 \in \mathbb{Z} $，$ r_1 \in [0, 2\pi) $ 与 $ r_2 \in (0, 2\pi] $ 使得

\[
t=a+2\pi k_{1}+r_{1}=a+2\pi k_{2}+r_{2}.
\]

由推论 6.14，$ \mathrm{cis}(a+r_1)=\mathrm{cis}(a+r_2)=\mathrm{cis}t=z $，于是这些函数也是满射。■

\end{subitems}

\end{proposition}

\begin{theorem}

\begin{subitems}
\item

$ \cos z = \cos(z + 2k\pi) $，$ \sin z = \sin(z + 2k\pi) $，$ z \in \mathbb{C} $，$ k \in \mathbb{Z} $，
即 cos 与 sin 是以 $ 2\pi $ 为周期的周期函数。

\item

对所有 $ z \in \mathbb{C} $，

\[
\cos z=0\Longleftrightarrow z\in\pi/2+\pi\mathbb{Z},
\]
\[
\sin z=0\Longleftrightarrow z\in\pi\mathbb{Z}.
\]
\item

函数 $ \sin: \mathbb{R} \to \mathbb{R} $ 在 $ (0, \pi) $ 上为正，且在闭区间 $ [0, \pi/2] $ 上严格递增。

\item

$ \cos(z + \pi) = -\cos z $，$ \sin(z + \pi) = -\sin z $，$ z \in \mathbb{C} $。

\item

$ \cos z = \sin(\pi/2 - z) $，$ \sin z = \cos(\pi/2 - z) $，$ z \in \mathbb{C} $。

\item

$ \cos(\mathbb{R}) = \sin(\mathbb{R}) = [-1, 1] $。

\proofhead
断言 (i) 由 (6.3) 与推论 6.14 可得。
\end{subitems}

\begin{subitems}
\item[（ii）]

由 (6.3) 与命题 6.13 我们有

\[
\cos z=0\Longleftrightarrow e^{i z}+e^{-i z}=0\Longleftrightarrow e^{2i z}=-1\Longleftrightarrow z\in\pi/2+\pi\mathbb{Z}.
\]

类似地

\[
\sin z=0\Longleftrightarrow e^{i z}-e^{-i z}=0\Longleftrightarrow e^{2i z}=1\Longleftrightarrow z\in\pi\mathbb{Z}.
\]
\item[（iii）]

由我们刚刚证明的结果，对 $ x \in (0, \pi) $ 有 $ \sin x \neq 0 $。不等式 (6.6) 表明，对所有 $ x \in (0, \sqrt{6}) $，$ \sin x $ 为正。由于介值定理（定理 5.1），实际上我们必须有

\begin{equation}
\sin x>0,\qquad x\in(0,\pi).
\end{equation}

类似地，由于 $ \cos(0)=1 $ 且对所有 $ t\in(-\pi/2,\pi/2) $ 有 $ \cos t\neq0 $，介值定理蕴含对所有在 $ (-\pi/2,\pi/2) $ 中的 $ t $，$ \cos t>0 $。
对第二个断言，设 $ 0 \leq x < y \leq \pi/2 $。由命题 6.3(ii) 我们有

\begin{equation}
\sin y-\sin x=2\cos\frac{y+x}{2}\sin\frac{y-x}{2}.
\end{equation}

由于 $(y+x)/2$ 与 $(y-x)/2$ 都属于 $(0,\pi/2)$，(6.8) 的右边为正，因此 $\sin y > \sin x$。

\item[（iv）]

由 (ii) 我们有 $ \sin\pi=0 $，于是由命题 6.13(ii)，

\[
\cos\pi=\cos\pi+i\sin\pi=e^{i\pi}=-1.
\]

现在设 $ z \in \mathbb{C} $。由命题 6.3(i) 我们得到

\[
\cos(z+\pi)=\cos z\cos\pi-\sin z\sin\pi=-\cos z
\]

以及

\[
\sin(z+\pi)=\sin z\cos\pi+\cos z\sin\pi=-\sin z.
\]
\item[（v）]

由 (ii) 我们有 $ \cos(\pi/2)=0 $，于是利用 (iii) 与推论 6.4，我们得到

\[
0<\sin(\pi/2)=\left|\sin(\pi/2)\right|=\sqrt{1-\cos^{2}(\pi/2)}=1.
\]

由命题 6.3(i) 我们现在有

\[
\cos(\pi/2-z)=\cos(\pi/2)\cos z+\sin(\pi/2)\sin z=\sin z
\]

以及

\[
\sin(\pi/2-z)=\sin(\pi/2)\cos z-\cos(\pi/2)\sin z=\cos z.
\]
\item[（vi）]

我们已在 (6.5) 中证明 $ \cos(\mathbb{R}) = [-1, 1] $。由 (v) 我们还得到 $ \sin(\mathbb{R}) = \cos(\mathbb{R}) $。■

\end{subitems}

\end{theorem}

\begin{remark}

\begin{items}
\item

由于等式

\[
\sin(x+\pi)=-\sin x\;,\quad\cos x=\sin(\pi/2-x)\;,\qquad x\in\mathbb{R}\;,
\]

以及正弦为奇函数这一事实，实正弦函数与实余弦函数完全由 $ \sin x $ 在 $ [0,\pi/2] $ 上的值所确定。

\end{items}

\bookfigure[0.46]{figures/322ab9ea6d5fb2ac.jpg}

\begin{items}
\item[（b）]

$ \pi/2 $ 是余弦函数的最小正零点。
原则上，这一观察连同余弦级数可用来以任意精度逼近数 $ \pi $。例如，由推论 II.7.9，我们有不等式

\[
1-\frac{t^{2}}{2}<\cos t<1-\frac{t^{2}}{2}+\frac{t^{4}}{24},\qquad t\in\mathbb{R}^{\times},
\]

于是 $ \cos 2 < -1/3 $，且对所有 $ 0 \leq t < \sqrt{2} $ 有 $ \cos t > 0 $。由介值定理我们知道，余弦函数在区间 $ [\sqrt{2}, 2) $ 中有一个零点。事实上，由于对所有 $ 0 \leq t < \sqrt{2} $ 有 $ \cos t > 0 $，最小正零点（即 $ \pi/2 $）必在此区间内。于是 $ 2\sqrt{2} \leq \pi < 4 $。由于两个不同的零点相隔距离至少为 $ \pi $，故 $ \pi/2 $ 是区间 $ (0, 2) $ 中唯一的零点。
为得到更好的逼近，取区间 $ [\sqrt{2}, 2) $ 中点处的某个 $ t $，然后利用余弦级数与推论 II.7.9 来计算 $ \cos t $ 的符号。
这将确定 $ \pi/2 $ 是在 $ [\sqrt{2}, t) $ 中还是在 $ (t, 2) $ 中。重复这一过程，$ \pi $ 可被以任意精度确定。经过相当大的努力，人们得到\footnote{记住 $\pi$ 的数字的一个常用口诀是：“How I like a drink, alcoholic of course, after the heavy lectures involving quantum mechanics.” 每个单词的字母个数给出 $\pi$ 的相应数字。}

\[
\pi=3.14159~26535~89793~23846~26433~83279\quad\therefore
\]

我们稍后将发展一种计算 $ \pi $ 的高效得多的方法。

\item[（c）]

若一个复数是某个系数为整数的非常数多项式的零点，则称其为代数数。不是代数数的复数称为超越数。特别地，实超越数是无理数。
1882 年，F. 林德曼证明了 $ \pi $ 是超越数。这连同代数的经典结果，给出了“化圆为方”不可能性的数学证明。也就是说，仅用直尺和圆规不可能作一个面积等于给定圆面积的正方形。

\end{items}

\end{remark}

\subsection*{正切函数与余切函数}

正切函数与余切函数定义为

\[
\tan z:=\frac{\sin z}{\cos z}\;,\quad z\in\mathbb{C}\backslash\left(\frac{\pi}{2}+\pi\mathbb{Z}\right)\;,\qquad\cot z:=\frac{\cos z}{\sin z}\;,\quad z\in\mathbb{C}\backslash\pi\mathbb{Z}\;.
\]

限制到实数上，正切函数与余切函数有如下图像：

\bookfigure[0.24]{figures/25d7278adcbe4c4a.jpg}

正切

\bookfigure[0.24]{figures/1b334d39d1b5e1cb.jpg}

余切

\begin{remark}

\begin{items}
\item

正切函数与余切函数是连续的、以 $ \pi $ 为周期的周期函数，且为奇函数。

\item

正切函数的加法定理成立：

\[
\tan(z\pm w)=\frac{\tan z\pm\tan w}{1\mp\tan z\tan w}
\]

对所有满足 $ z \pm w \in \mathrm{dom}(\tan) $ 的 $ w, z \in \mathrm{dom}(\tan) $。

\proofhead
这由命题 6.3(i) 易得。
\item

对所有 $ z \in \mathbb{C} \setminus \pi \mathbb{Z} $，$ \cot z = -\tan(z - \pi/2) $。

\proofhead
这直接由定理 6.16(iv) 可得。
\end{items}

\end{remark}

\subsection*{复指数函数}

在命题 6.1(v) 与 6.15 中我们已经看到，函数

\[
\mathrm{c i s}\colon[0,2\pi)\to S^{1}~,\quad t\mapsto e^{i t}
\]

是连续的双射，于是对每个 $ z \in S^1 $，存在唯一的 $ \alpha \in [0, 2\pi) $ 使得

\[
z=e^{i\alpha}=\mathrm{c i s}(\alpha)=\cos\alpha+i\sin\alpha.
\]

数 $ \alpha \in [0, 2\pi) $ 可解释为从 1 到 $ z = e^{i\alpha} $ 的圆弧长度（见习题 12），同样也可解释为一个角。此外，由定理 6.16 我们知道 cis: $ \mathbb{R} \to S^1 $ 有周期 $ 2\pi $。因此这个函数把实轴绕 $ S^1 $ 缠绕无穷多次。

\begin{proposition}

对 $ a \in \mathbb{R} $，设 $ I_a $ 是形如 $ [a, a + 2\pi) $ 或 $ (a, a + 2\pi] $ 的区间。则函数

\bookfigure[0.27]{figures/850ff9b9997bab31.jpg}
\begin{equation}
\exp(\mathbb{R}+i I_{a}):\mathbb{R}+i I_{a}\to\mathbb{C}^{\times}~,\quad z\mapsto e^{z}
\end{equation}

是连续的双射。

\proofhead
连续性由定理 6.1(v) 显然。为验证单射性，我们假设存在 $w, z \in \mathbb{R} + iI_a$ 使得 $e^z = e^w$。把 $z = x + iy$ 与 $w = \xi + i\eta$ 写出来，其中 $x, y, \xi$ 与 $\eta$ 为实数。则引理 6.10 蕴含
\begin{equation}
e^{x}=|e^{x}e^{iy}|=|e^{\xi}e^{i\eta}|=e^{\xi}~,
\end{equation}

于是由命题 6.5，$ x=\xi $。我们现在有

\[
e^{i(y-\eta)}=e^{x+i y-(\xi+i y)}=e^{z}/e^{w}=1~,
\]

由命题 6.13，这蕴含 $ y - \eta \in 2\pi\mathbb{Z} $。由假设，$ |y - \eta| < 2\pi $，于是 $ y = \eta $。这样我们就证明了 (6.9) 中的函数是单射。

设 $ w \in \mathbb{C}^\times $。对 $ x := \log|w| \in \mathbb{R} $ 我们有 $ e^x = |w| $。由命题 6.15，存在唯一的 $ y \in I_a $ 使得 $ e^{iy} = w/|w| \in S^1 $。令 $ z = x + iy \in \mathbb{R} + iI_a $，我们有 $ e^z = e^xe^{iy} = |w|(w/|w|) = w $。■

\end{proposition}

上一命题中的函数可如下用图形表示。

\bookfigure[0.53]{figures/05774ddf5192baea.jpg}

最后我们注意到

\[
\mathbb{C}=\bigcup_{k\in\mathbb{Z}}\left\{\mathbb{R}+i\left[a+2k\pi,a+2(k+1)\pi\right)\right\}
\]

是复平面的一个划分，于是由命题 6.19 可知，指数函数 $ \exp: \mathbb{C} \to \mathbb{C}^\times $ 把复平面绕原点缠绕无穷多次，无穷多次地覆盖去心复平面 $ \mathbb{C}^\times $。

\subsection*{极坐标}

利用指数函数，我们可以用极坐标表示复数。在这种表示下，两个复数的乘法有简单的几何解释。

\begin{theorem}

（复数的极坐标表示）每个 $ z \in \mathbb{C}^\times $ 都有如下形式的唯一表示

\[
z=|z|e^{i\alpha}
\]

其中 $ \alpha \in [0, 2\pi) $。

\proofhead
这直接由命题 6.19 可得。■
\end{theorem}

该定理中的实数 $ \alpha \in [0, 2\pi) $ 称为 $ z \in \mathbb{C}^\times $ 的规范化辐角，记作 $ \arg_N(z) $。

\begin{remark}

\begin{items}
\item

（复数的乘积）设 $ w, z \in \mathbb{C}^\times $，$ \alpha := \arg_N(z) $ 与 $ \beta := \arg_N(w) $。把 $ z $ 与 $ w $ 相乘，我们得到 $ zw = |z||w| e^{i(\alpha + \beta)} $，于是由引理 6.10 与推论 6.14，

\end{items}

\bookfigure[0.24]{figures/ffe10fdf4a0dcfb3.jpg}
\begin{equation*}
\begin{aligned}
\left|zw\right|&=\left|z\right|\left|w\right|,\\\arg_{N}(zw)&\equiv\arg_{N}(z)+\arg_{N}(w)
\end{aligned}
\end{equation*}
\bookfigure[0.26]{figures/40fbe67dd0a39219.jpg}

模 2π。

\begin{items}
\item[（b）]

对 $ n \in \mathbb{N}^\times $，方程 $ z^n = 1 $ 恰有 $ n $ 个复数解，即 $ n $ 次单位根，

\[
z_{k}:=e^{2\pi i k/n}~,\qquad k=0,\ldots,n-1~.
\]

点 $ z_k $ 在单位圆上，且是以 1 为一个顶点的正 $ n $ 边形的顶点。

\item[（c）]

对所有 $ a \in \mathbb{C} $ 与 $ k \in \mathbb{N}^\times $，方程 $ z^k = a $ 在 $ \mathbb{C} $ 中可解。\footnote{这填补了例 3.9(b) 中代数基本定理证明中的空缺。它还表明习题 I.11.15 中关于三次方程的解的假设成立。}

\end{items}

\bookfigure[0.19]{figures/d86c104097b930f0.jpg}

\proofhead
若 $a=0$，断言显然。否则，$a=|a|e^{i\alpha}$，其中 $\alpha:=\arg_{N}(a)\in[0,2\pi)$。令 $z:=\sqrt[k]{|a|}e^{i\alpha/k}$。则
\[
z^{k}=(\sqrt[k]{|a|}e^{i\alpha/k})^{k}=(\sqrt[k]{|a|})^{k}(e^{i\alpha/k})^{k}=|a|e^{i\alpha}=a,
\]

于是我们找到了所需的解。■

\begin{items}
\item[（d）]

（平面的极坐标表示）对每个 $ (x, y) \in \mathbb{R}^2 \setminus \{(0,0)\} $，存在唯一的实数 $ r > 0 $ 与 $ \alpha \in [0, 2\pi) $ 使得

\[
x=r\cos\alpha\quad\text{ 且 }\quad y=r\sin\alpha.
\]

\proofhead
设 $ x, y \in \mathbb{R} $，且 $ z := x + iy \in \mathbb{C}^\times $。令
\[
r:=|z|=\sqrt{x^{2}+y^{2}}>0\quad\text{且}\quad\alpha:=\arg_{N}(z)\in[0,2\pi)~.
\]

则，由定理 6.20 与欧拉公式，我们有

\[
x+i y=z=r e^{i\alpha}=r\cos\alpha+i r\sin\alpha~,
\]

断言由此得证。■

\item[（e）]

对所有 $ z \in \mathbb{C} $，$ |e^z| = e^{\mathrm{Re}z} $。

\proofhead
这由 $ |e^z| = |e^{\mathrm{Re}z} e^{i \mathrm{Im} z}| = e^{\mathrm{Re}z} |e^{i \mathrm{Im} z}| = e^{\mathrm{Re}z} $ 可得。■
\end{items}

\end{remark}

\subsection*{复对数}

对给定的 $ w \in \mathbb{C}^\times $，我们要确定方程 $ e^z = w $ 的所有解。由定理 6.20 我们知道这个方程可解，因为

\[
w=e^{\log|w|+i\arg_{N}(w)}.
\]

现在设 $ z \in \mathbb{C} $ 是 $ e^z = w $ 的任意一个解。由推论 6.14 与命题 6.15，存在唯一的 $ k \in \mathbb{Z} $ 使得 $ z = \log |w| + i \arg_N(w) + 2\pi ki $。因此

\[
\left\{\log|w|+i\left(\arg_{N}(w)+2\pi k\right)\in\mathbb{C}\;;\;k\in\mathbb{Z}\right\}
\]

是方程 $ e^{z} = w $ 的所有解的集合。集合

\[
\mathrm{A r g}(w):=\mathrm{a r g}_{N}(w)+2\pi\mathbb{Z}
\]

称为 w 的辐角，而集合

\[
\mathrm{Log}(w):=\log|w|+i\operatorname{Arg}(w)
\]

称为 w 的（复）对数。

这些等式定义了两个集值函数

\[
\operatorname{A r g}\colon\mathbb{C}^{\times}\to\mathcal{P}(\mathbb{C})\;,\quad w\mapsto\operatorname{A r g}(w)\;,
\]
\[
\operatorname{L o g}\colon\mathbb{C}^{\times}\to\mathcal{P}(\mathbb{C})\;,\quad w\mapsto\operatorname{L o g}(w)\;,
\]

分别称为辐角函数与对数函数。

由于集值函数一般相当笨重，我们利用如下事实：对每个 $w \in \mathbb{C}^\times$，存在唯一的 $\varphi =: \arg(w) \in (-\pi, \pi]$ 使得 $w = |w| e^{i\varphi}$。这定义了一个实值函数

\[
\arg:\mathbb{C}^{\times}\to(-\pi,\pi],\quad w\mapsto\arg(w)
\]

称为辐角的主值。对数的主值\footnote{对 $w \in (0, \infty)$，这个定义与实对数 $(\exp|_{\mathbb{R}})^{-1}$ 一致。}定义为

\[
\operatorname{l o g}\colon\mathbb{C}^{\times}\to\mathbb{R}+i(-\pi,\pi]~,\quad w\mapsto\operatorname{l o g}|w|+i\operatorname{a r g}(w)~.
\]

命题 6.5 与 6.15 蕴含 $ \log $ 是双射，且

\begin{equation}
\begin{aligned}
&e^{\log w}=w\ ,\qquad w\in\mathbb{C}^{\times}\ ,\\&\log e^{z}=z\ ,\qquad z\in\mathbb{R}+i\left(-\pi,\pi\right].\\
\end{aligned}
\end{equation}

特别地，$ \log w $ 对 $ w < 0 $ 有定义，且此时 $ \log w = \log |w| + i\pi $。

对集值复对数，我们有

\[
e^{\operatorname{L o g}w}=w~,\quad w\in\mathbb{C}^{\times}~,\qquad\operatorname{L o g}e^{z}=z+2\pi i\mathbb{Z}~,\quad z\in\mathbb{C}~.
\]

最后，\footnote{关于集合上 $+$ 与 $-$ 的含义，见 (I.4.1)。}

\begin{equation}
\operatorname{Log}(zw)=\operatorname{Log}z+\operatorname{Log}w~,\quad\operatorname{Log}(z/w)=\operatorname{Log}z-\operatorname{Log}w
\end{equation}

对所有 $ w, z \in \mathbb{C}^x $。这可以类似于自然对数的加法定理（定理 6.6）来证明。

\subsection*{复幂}

对 $ z \in \mathbb{C}^\times $ 与 $ w \in \mathbb{C} $，

\[
z^{w}:=e^{w\operatorname{L o g}z}
\]

称为 $ z $ 的（复）幂。由于 Log 是集值函数，$ z^{w} $ 是一个集合。具体地，

\[
{z}^{w}=\left\{e^{w(\log|z|+i(\arg_{N}(z)+2\pi k))}\;;\;k\in\mathbb{Z}\right\}.
\]

$ z^{w} $ 的主值当然利用对数的主值来定义：

\[
\mathbb{C}^{\times}\to\mathbb{C}~,\qquad z\mapsto z^{w}:=e^{w\log z}~.
\]

命题 6.8 中的规则容易推广到幂函数的主值：

\begin{equation}
z^{a}z^{b}=z^{a+b}\quad\text{且}\quad z^{a}\cdot w^{a}=(z w)^{a}
\end{equation}

对所有 $ w, z \in \mathbb{C}^\times $ 与 $ a, b \in \mathbb{C} $。

\begin{remark}

\begin{items}
\item

定理 6.1(iii) 表明

\begin{equation}
exp:({\mathbb{C}},+)\to({\mathbb{C}}^{\times},\cdot)
\end{equation}

是阿贝尔群 $(\mathbb{C},+)$ 与 $(\mathbb{C}^{\times},\cdot)$ 之间的群同态。此外，命题 6.13 与 6.15 蕴含 (6.14) 是满射，且其核为 $2\pi i\mathbb{Z}$。由例 I.7.8(c)，商群 $(\mathbb{C},+)/(2\pi i\mathbb{Z})$ 同构于 $(\mathbb{C}^{\times},\cdot)$。

\item

单位圆 $ S^{1} $ 构成阿贝尔群 $ (S^{1}, \cdot) $，即乘法之下的圆群（见习题 I.11.9）。由定理 6.1(iii) 与命题 6.13、6.15 可知

\[
\mathrm{c i s}:(\mathbb{R},+)\to(S^{1},\cdot)
\]

是核为 $ 2\pi\mathbb{Z} $ 的满群同态。因此群 $ (S^1, \cdot) $ 与 $ (\mathbb{R}, +)/(2\pi\mathbb{Z}) $ 同构。

\item

函数

\[
\exp_{\mathbb{R}}:(\mathbb{R},+)\to\big((0,\infty),\cdot\big)
\]

是群同构，其逆为 log: $ (0, \infty) \to \mathbb{R} $。■

\end{items}

\end{remark}

\subsection*{指数函数的进一步表示}

在习题 II.4.3 中我们看到，对有理变元，指数函数由下式给出

\[
e^{r}=\lim_{n\to\infty}\left(1+\frac{r}{n}\right)^{n}.
\]

这一结果可推广到任意复数。

\begin{theorem}

对所有 $ z \in \mathbb{C} $，

\[
e^{z}=\lim_{n\to\infty}\left(1+\frac{z}{n}\right)^{n}.
\]

\proofhead
设 $ z \in \mathbb{C} $。由习题 I.8.1 我们有
\[
a^{n}-b^{n}=(a-b)\sum_{k=0}^{n-1}a^{k}b^{n-k-1},\qquad a,b\in\mathbb{C},
\]

于是

\begin{equation}
\begin{aligned}
e^{z}-(1+z/n)^{n}&=(e^{z/n})^{n}-(1+z/n)^{n}\\&=\left[e^{z/n}-(1+z/n)\right]\sum_{k=0}^{n-1}(e^{z/n})^{k}(1+z/n)^{n-1-k}.
\end{aligned}
\end{equation}

由例 2.25(b) 我们知道

\begin{equation}
r_{n}:=\left[\frac{e^{z/n}-1}{z/n}-1\right]\to0\;,\qquad n\to\infty\;.
\end{equation}

为估计

\begin{equation}
L_{n}:=\sum_{k=0}^{n-1}(e^{z/n})^{k}(1+z/n)^{n-1-k},\qquad n\in\mathbb{N}^{\times},
\end{equation}

我们利用不等式

\[
|e^{w}|=e^{\mathrm{R e}w}\leq e^{|w|},\quad|1+w|\leq1+|w|\leq e^{|w|}
\]

得到

\begin{equation}
|L_{n}|\leq\sum_{k=0}^{n-1}(e^{|z|/n})^{k}(e^{|z|/n})^{n-1-k}=n(e^{|z|/n})^{n-1}\leq n e^{|z|},\qquad n\in\mathbb{N}^{\times}.
\end{equation}

结合 (6.15)、(6.17) 与 (6.18) 我们得到

\[
\left|e^{z}-\left(1+\frac{z}{n}\right)^{n}\right|=\left|\frac{z}{n}r_{n}L_{n}\right|\leq\left|\frac{z}{n}\right|\left|r_{n}\right|n e^{\left|z\right|}=\left|z\right|e^{\left|z\right|}\left|r_{n}\right|,
\]

这连同 (6.16) 证明了断言。

\end{theorem}

\exerciseshead

\begin{exercises}
\item

证明函数 cis: $ \mathbb{R} \to \mathbb{C} $ 与 $ \cos, \sin : \mathbb{R} \to \mathbb{R} $ 是利普希茨连续的，且利普希茨常数为 1。（提示：见例 2.25(b)。）

\item

对 $ z \in \mathbb{C} $ 与 $ m \in \mathbb{N} $，证明棣莫弗公式。

\[
\left(\cos z+i\sin z\right)^{m}=\cos(m z)+i\sin(m z).
\]
\item

证明下列三角恒等式：

\[
(a)\cos^{2}(z/2)=(1+\cos z)/2,\ \sin^{2}(z/2)=(1-\cos z)/2,\ z\in\mathbb{C}.
\]
\[
\mathrm{(b)}\operatorname{t a n}(z/2)=(1-\operatorname{c o s}z)/\operatorname{s i n}z=\operatorname{s i n}z/(1+\operatorname{c o s}z),z\in\mathbb{C}\backslash(\pi\mathbb{Z}).
\]
\item

双曲余弦函数与双曲正弦函数定义为

\[
\operatorname{c o s h}(z):=\frac{e^{z}+e^{-z}}{2}\quad\text{且}\quad\operatorname{s i n h}(z):=\frac{e^{z}-e^{-z}}{2}~,\qquad z\in\mathbb{C}~.
\]

对 $ w, z \in \mathbb{C} $，证明下列各式：

\begin{items}
\item

$ \cosh^{2}z - \sinh^{2}z = 1 $。

\item

$ \cosh(z+w) = \cosh z \cosh w + \sinh z \sinh w $。

\item

$ \sinh(z+w) = \sinh z \cosh w + \cosh z \sinh w $。

\item

$ \cosh z = \cos iz $，$ \sinh z = -i \sin iz $。

\item

$ \cosh z = \sum_{k=0}^{\infty} \frac{z^{2k}}{(2k)!} $，$ \sinh z = \sum_{k=0}^{\infty} \frac{z^{2k+1}}{(2k+1)!} $。

\end{items}

\item

双曲正切函数与双曲余切函数由下式定义：

\[
\operatorname{t a n h}z:=\frac{\sinh(z)}{\cosh(z)}\:,\quad z\in\mathbb{C}\backslash\pi i\left(\mathbb{Z}+1/2\right)\:,\qquad\text{ 且 }\qquad\coth z:=\frac{\cosh(z)}{\sinh(z)}\:,\quad z\in\mathbb{C}\backslash\pi i\mathbb{Z}\:.
\]

函数 $ \cosh $、$ \sinh $、$ \tanh $ 与 $ \coth $ 对实变元取实值。画出这些实值函数的图像。

证明下列各式：

\begin{items}
\item

函数

\[
\mathrm{c o s h}~,\quad\mathrm{s i n h}~,\quad\mathrm{t a n h}:\mathbb{C}\backslash\pi i(\mathbb{Z}+1/2)\to\mathbb{C}~,\quad\mathrm{c o t h}:\mathbb{C}^{\times}\to\mathbb{C}\backslash\pi i\mathbb{Z}
\]

都是连续的。

\item

$ \lim_{x \to \pm \infty} \tanh(x) = \pm 1 $，$ \lim_{x \to \pm 0} \coth(x) = \pm \infty $。

\item

$ \cosh: [0,\infty) \to \mathbb{R} $ 严格递增，且 $ \cosh([0,\infty)) = [1,\infty) $。

\item

$ \sinh: \mathbb{R} \to \mathbb{R} $ 严格递增且是双射。

\item

tanh:  $ \mathbb{R} \to (-1,1) $ 严格递增且是双射。

\item

coth:  $ (0, \infty) \to \mathbb{R} $ 严格递减，且 $ \coth((0, \infty)) = (1, \infty) $。

\item

tanh:  $ \mathbb{R} \to (-1,1) $ 是利普希茨连续的，且利普希茨常数为 1。

\end{items}

\item

确定下列极限：

\[
\mathrm{(a)}\operatorname*{l i m}_{x\to{0+}}x^{x}~,\quad\mathrm{(b)}\operatorname*{l i m}_{x\to{0+}}x^{1/x}~,\quad\mathrm{(c)}\operatorname*{l i m}_{z\to{0}}\frac{\operatorname{l o g}(1+z)}{z}~.
\]

（提示：(c) 见例 2.25(b)。）

\item

对 x, y > 0，证明不等式

\[
\frac{\log x+\log y}{2}\leq\log\Bigl(\frac{x+y}{2}\Bigr).
\]
\item

确定下列极限：

\[
\mathrm{(a)}\operatorname*{l i m}_{z\to0}\frac{\operatorname{s i n}z}{z}~,\qquad\mathrm{(b)}\operatorname*{l i m}_{z\to0}\frac{a^{z}-1}{z}~,\quad a\in\mathbb{C}^{\times}~.
\]
\item

证明函数

\[
\arg:\mathbb{C}\backslash(-\infty,0]\to(-\pi,\pi),\quad\log:\mathbb{C}\backslash(-\infty,0]\to\mathbb{R}+i(-\pi,\pi),
\]

都是连续的。（提示：(i)  $ \arg = \arg \circ \nu $，其中对所有 $ z \in \mathbb{C}^{\times} $ 有 $ \nu(z) := z/|z| $）(ii)  $ \arg\big|_{(S^{1}\setminus\{-1\})}=\big[\mathrm{cis}\big|_{(-\pi,\pi)}\big]^{-1} $ (iii) 对形如 $ [-a,a] $（$ a\in(0,\pi) $）的区间利用习题 3.3(b)。）

\item

证明幂函数主值的下列法则：

\[
z^{a}z^{b}=z^{a+b}~,\quad z^{a}w^{a}=(z w)^{a}~,\qquad z,w\in\mathbb{C}^{\times}~,\quad a,b\in\mathbb{C}~.
\]
\item

计算 $ i^{i} $ 及其主值。

\item

设 $ x \in \mathbb{R} $，$ n \in \mathbb{N}^\times $，且对所有 $ k = 0,1,\ldots,n $ 有 $ z_{n,k} := e^{ixk/n} \in S^1 $。令

\[
L_{n}:=\sum_{k=1}^{n}\left|z_{n,k}-z_{n,k-1}\right|,
\]

即以 $ z_{n,0}, z_{n,1}, \ldots, z_{n,n} $ 为顶点的折线路径的长度。证明

\[
L_{n}=2n\left|\sin\bigl(x/(2n)\bigr)\right|\quad\text{且}\quad\operatorname*{l i m}_{n\to\infty}L_{n}=\left|x\right|.
\]

注 对大的 $n \in \mathbb{N}$ 与 $x \in [0, 2\pi]$，$[0, x]$ 在函数 cis 之下的像由以 $z_{n,0}, z_{n,1}, \ldots, z_{n,n}$ 为顶点的折线路径逼近。于是 $L_n$ 是圆上 1 与 cis($x) = e^{ix}$ 之间弧长的近似。本习题表明，函数 cis: $\mathbb{R} \to S^1$ 把直线 $\mathbb{R}$ 缠绕到 $S^1$ 上的方式使长度得以保持。

\item

研究函数 $ \mathbb{C} \to \mathbb{C} $，$ z \mapsto z^2 $ 的性质。特别地，对 $ z = x + iy $，计算双曲线 $ x^2 - y^2 = \text{const} $、$ xy = \text{const} $ 以及直线 $ x = \text{const} $、$ y = \text{const} $ 的像。

\item

确定下列方程在 C 中的全部解：

\begin{items}
\item

$ z^{4} = (\sqrt{2}/2)(1 + i) $。

\item

$ z^{5}=i $。

\item

$ z^{3} + 6z + 2 = 0 $

\item

$ z^{3}+(1-2i)z^{2}-(1+2i)z-1=0 $
（提示：对 (c) 与 (d) 中的三次方程，利用习题 I.11.15。）

\end{items}

\item

对 $ x \in \mathbb{R} $ 与 $ n \in \mathbb{N} $，令

\[
f_{n}(x):=\lim_{k\to\infty}\left(\cos(n!\pi x)\right)^{2k}.
\]

确定 $ \lim_{n \to \infty} f_n(x) $。（提示：分别考虑 $ x \in \mathbb{R} \setminus \mathbb{Q} $ 的情形，并利用 $ |\cos(m\pi)| = 1 $ 当且仅当 $ m \in \mathbb{Z} $ 这一事实。）

\item

证明 $ \cosh 1 $ 是无理数。（提示：习题 II.7.10。）

\end{exercises}
